\documentclass[12pt,a4paper]{article}
\usepackage[utf8]{inputenc}
\usepackage{amsmath}
\usepackage{amssymb}
\usepackage{enumerate}
\usepackage{xcolor}
\usepackage{geometry}
\geometry{a4paper, margin=2cm}
\usepackage{hyperref}

\hypersetup{
    colorlinks=true,
    linkcolor=blue,
    filecolor=magenta,      
    urlcolor=cyan,
    pdftitle={CEA201 - Ultra Edition Final v5}
}

\begin{document}

\title{\textbf{CEA201 - Ultra Edition Final v5: Reclassified Question Bank}\\\large (Complete 531 Questions Categorized Across 21 Chapters - COA 11th Edition)}
\author{\textbf{Gemini Notebook Study Workbench}}
\date{\today}
\maketitle

\begin{abstract}
\noindent This document aggregates all 531 multiple-choice questions from the CEA201 (Computer Organization and Architecture 11th Edition) question bank. Every single question has been audited, cross-verified, and reclassified into its corresponding chapter (Chapters 1 to 21) based on COA 11th Edition syllabus standards.
\end{abstract}

\tableofcontents
\newpage

\section{Chapter 01: Basic Concepts and Computer Evolution}
\textit{Total questions in chapter: 55}

\textbf{Question 40:} Second generation computers used \underline{\hspace{1.5cm}}.

\begin{itemize}
    \item[\textbf{A.}] integrated circuits
    \item[\textbf{B.}] large-scale integration
    \item[\textbf{C.}] \textcolor{red}{\textbf{transistors}}
    \item[\textbf{D.}] vacuum tubes
\end{itemize}

\vspace{0.5cm}

\vspace{0.5cm}

\textbf{Question 159:} The \underline{\hspace{1.5cm}} defines the third generation of computers.

\begin{itemize}
    \item[\textbf{A.}] \textcolor{red}{\textbf{integrated circuit}}
    \item[\textbf{B.}] vacuum tube
    \item[\textbf{C.}] transistor
    \item[\textbf{D.}] VLSI
\end{itemize}

\vspace{0.5cm}

\vspace{0.5cm}

\textbf{Question 162:} The end user is concerned mainly with the computer's architecture.

\begin{itemize}
    \item[\textbf{A.}] True
    \item[\textbf{B.}] \textcolor{red}{\textbf{False}}
\end{itemize}

\vspace{0.5cm}

\vspace{0.5cm}

\textbf{Question 208:} The ENIAC used \underline{\hspace{1.5cm}}.

\begin{itemize}
    \item[\textbf{A.}] \textcolor{red}{\textbf{vacuum tubes}}
    \item[\textbf{B.}] transistors
    \item[\textbf{C.}] integrated circuits
    \item[\textbf{D.}] IAS
\end{itemize}

\vspace{0.5cm}

\vspace{0.5cm}

\textbf{Question 233:} Computer organization refers to attributes of a system visible to the programmer.

\begin{itemize}
    \item[\textbf{A.}] True
    \item[\textbf{B.}] \textcolor{red}{\textbf{False}}
\end{itemize}

\vspace{0.5cm}

\vspace{0.5cm}

\textbf{Question 236:} A vacuum tube is a solid-state device made from silicon.

\begin{itemize}
    \item[\textbf{A.}] True
    \item[\textbf{B.}] \textcolor{red}{\textbf{False}}
\end{itemize}

\vspace{0.5cm}

\vspace{0.5cm}

\textbf{Question 245:} The ENIAC is an example of a \underline{\hspace{1.5cm}} generation computer.

\begin{itemize}
    \item[\textbf{A.}] \textcolor{red}{\textbf{first}}
    \item[\textbf{B.}] second
    \item[\textbf{C.}] third
    \item[\textbf{D.}] fourth
\end{itemize}

\vspace{0.5cm}

\vspace{0.5cm}

\textbf{Question 256:} Which of the following components was used in the first ENIAC computer?

\begin{itemize}
    \item[\textbf{A.}] Bipolar transistors
    \item[\textbf{B.}] Field transistors
    \item[\textbf{C.}] \textcolor{red}{\textbf{Vacuum tubes}}
    \item[\textbf{D.}] Semiconductor Ics
\end{itemize}

\vspace{0.5cm}

\vspace{0.5cm}

\textbf{Question 257:} Which of the following statements is true for Von Neumann architecture?

\begin{itemize}
    \item[\textbf{A.}] \textcolor{red}{\textbf{Shared bus between the program memory and data memory}}
    \item[\textbf{B.}] Separate bus between the program memory and data memory
    \item[\textbf{C.}] External bus for program memory and data memory
    \item[\textbf{D.}] External bus for data memory only
\end{itemize}

\vspace{0.5cm}

\vspace{0.5cm}

\textbf{Question 260:} The first generation of computers used \underline{\hspace{1.5cm}} for digital logic elements and memory?

\begin{itemize}
    \item[\textbf{A.}] Transistor
    \item[\textbf{B.}] Integrated Circuits
    \item[\textbf{C.}] Large-scale integration
    \item[\textbf{D.}] \textcolor{red}{\textbf{Vacuum Tubes}}
\end{itemize}

\vspace{0.5cm}

\vspace{0.5cm}

\textbf{Question 305:} What is the distinction between Computer Architecture and Computer Organization?

\begin{itemize}
    \item[\textbf{A.}] Computer Architecture is the way the system is structured, while Computer Organization is those attributes of a system that are visible to the user
    \item[\textbf{B.}] \textcolor{red}{\textbf{Computer Architecture is those attributes of a system that are visible to the user, while Computer Organization is the way the system is structured}}
    \item[\textbf{C.}] Computer Architecture and Computer Organization are the same
    \item[\textbf{D.}] Computer Architecture is slower than Computer Organization
\end{itemize}

\vspace{0.5cm}

\vspace{0.5cm}

\textbf{Question 357:} What is false about the von Neumann architecture?

\begin{itemize}
    \item[\textbf{A.}] Data and instructions are stored in a single read-write memory.
    \item[\textbf{B.}] The contents of this memory are addressable by location, without regard to the type of data contained there.
    \item[\textbf{C.}] Execution occurs in a sequential fashion (unless explicitly modified) from one instruction to the next.
    \item[\textbf{D.}] \textcolor{red}{\textbf{Data is stored in main memory and instructions are stored in cache memory}}
\end{itemize}

\vspace{0.5cm}

\vspace{0.5cm}

\textbf{Question 380:} Consider some computer generations 1: the first computer generation 2: the second computer generation 3: the third computer generation And some technologies are used in computer generations A: Vacuum Tubes B: Transistors C: Integrated Circuits (IC) D: Microprocessors Select the main technology applied in each computer generation

\begin{itemize}
    \item[\textbf{A.}] 1-D;2-C;3-B
    \item[\textbf{B.}] 1-A;2-C;3-D
    \item[\textbf{C.}] \textcolor{red}{\textbf{1-A;2-B;3-C}}
    \item[\textbf{D.}] 1-D;2-B;3-C
\end{itemize}

\vspace{0.5cm}

\vspace{0.5cm}

\textbf{Question 457:} Which of the following attributes do not belong to computer architecture?

\begin{itemize}
    \item[\textbf{A.}] Number of bits used to represent data types
    \item[\textbf{B.}] The instruction set
    \item[\textbf{C.}] \textcolor{red}{\textbf{Interface between the computer and peripherals}}
    \item[\textbf{D.}] Techniques for addressing memory
    \item[\textbf{E.}] \textcolor{red}{\textbf{The memory technology used
F: I/O mechanisms}}
\end{itemize}

\vspace{0.5cm}

\vspace{0.5cm}

\textbf{Question 488:} What is the first evolutionary step in the evolution of the I/O function?

\begin{itemize}
    \item[\textbf{A.}] The I/O module is enhanced to become a processor in its own right.
    \item[\textbf{B.}] A controller or I/O module is added, and programmed I/O without interrupts is used.
    \item[\textbf{C.}] \textcolor{red}{\textbf{The CPU directly controls a peripheral device.}}
    \item[\textbf{D.}] The I/O module is given direct access to memory via DMA.
    \item[\textbf{E.}] The I/O module has a local memory of its own.
\end{itemize}

\vspace{0.5cm}

\vspace{0.5cm}

\textbf{Question 165:} A computer must be able to process, store, move, and control data.

\begin{itemize}
    \item[\textbf{A.}] \textcolor{red}{\textbf{True}}
    \item[\textbf{B.}] False
\end{itemize}

\vspace{0.5cm}

\vspace{0.5cm}

\textbf{Question 308:} During the development process of the computer, which of the following statements is false?

\begin{itemize}
    \item[\textbf{A.}] The second generation uses transistors
    \item[\textbf{B.}] The first generation uses vacuum tubes
    \item[\textbf{C.}] The fourth generation uses integrated circuit
    \item[\textbf{D.}] \textcolor{red}{\textbf{The third generation uses transistor}}
\end{itemize}

\vspace{0.5cm}

\vspace{0.5cm}

\textbf{Question 26:} Interfaces between the computer and peripherals is an example of an organizational attribute.

\begin{itemize}
    \item[\textbf{A.}] \textcolor{red}{\textbf{True}}
    \item[\textbf{B.}] False
\end{itemize}

\vspace{0.5cm}

\vspace{0.5cm}

\textbf{Question 48:} The \underline{\hspace{1.5cm}} controls the movement of data and instructions into and out of the processor.

\begin{itemize}
    \item[\textbf{A.}] \textcolor{red}{\textbf{control unit}}
    \item[\textbf{B.}] ALU
    \item[\textbf{C.}] shifter
    \item[\textbf{D.}] branch
\end{itemize}

\vspace{0.5cm}

\vspace{0.5cm}

\textbf{Question 52:} The \underline{\hspace{1.5cm}} interprets the instructions in memory and causes them to be executed.

\begin{itemize}
    \item[\textbf{A.}] I/O
    \item[\textbf{B.}] \textcolor{red}{\textbf{control unit}}
    \item[\textbf{C.}] main memory
    \item[\textbf{D.}] arithmetic and logic unit
\end{itemize}

\vspace{0.5cm}

\vspace{0.5cm}

\textbf{Question 147:} Changes in technology not only influence organization but also result in the introduction of more powerful and more complex architectures.

\begin{itemize}
    \item[\textbf{A.}] \textcolor{red}{\textbf{True}}
    \item[\textbf{B.}] False
\end{itemize}

\vspace{0.5cm}

\vspace{0.5cm}

\textbf{Question 205:} A wafer is made of silicon and is broken up into chips which consists of many gates and/or memory cells plus a number of input and output attachment points.

\begin{itemize}
    \item[\textbf{A.}] \textcolor{red}{\textbf{True}}
    \item[\textbf{B.}] False
\end{itemize}

\vspace{0.5cm}

\vspace{0.5cm}

\textbf{Question 354:} \underline{\hspace{1.5cm}} refers to the operational units and their interconnections that realize the architectural specifications.

\begin{itemize}
    \item[\textbf{A.}] Computer architecture
    \item[\textbf{B.}] Computer function
    \item[\textbf{C.}] \textcolor{red}{\textbf{Computer organization}}
    \item[\textbf{D.}] Instruction set architecture
\end{itemize}

\vspace{0.5cm}

\vspace{0.5cm}

\textbf{Question 378:} Which of the following statements is NOT part of the Von Newmann principle?

\begin{itemize}
    \item[\textbf{A.}] Computers can operate according to a stored program
    \item[\textbf{B.}] The computer uses a program counter to indicate the location of the next statement
    \item[\textbf{C.}] A computer's memory is addressable
    \item[\textbf{D.}] \textcolor{red}{\textbf{Each statement must have a memory area containing the address of the next instruction}}
\end{itemize}

\vspace{0.5cm}

\vspace{0.5cm}

\textbf{Question 356:} What is Memory Address Register (MAR)?

\begin{itemize}
    \item[\textbf{A.}] Contains a word to be stored in memory or sent to the I/O unit, or is used to receive a word from memory or from the I/O unit.
    \item[\textbf{B.}] Employed to hold temporarily the righthand instruction from a word in memory.
    \item[\textbf{C.}] \textcolor{red}{\textbf{Contains the address in memory of the word to be written from or read into the MBR.}}
    \item[\textbf{D.}] Contains the address of the next instruction pair to be fetched from memory.
\end{itemize}

\vspace{0.5cm}

\vspace{0.5cm}

\textbf{Question 264:} The basic components of a computer are:

\begin{itemize}
    \item[\textbf{A.}] \textcolor{red}{\textbf{Main memory, CPU, I/O modules and system interconnection}}
    \item[\textbf{B.}] Main memory, CPU, I/O modules and Storage device
    \item[\textbf{C.}] Main Memory, CPU, Peripherals and Storage device
    \item[\textbf{D.}] Main memory, CPU, I/O modules and Storage device
\end{itemize}

\vspace{0.5cm}

\vspace{0.5cm}

\textbf{Question 519:} Choose an operation in which no operating system's utility is called.

\begin{itemize}
    \item[\textbf{A.}] \textcolor{red}{\textbf{$Y = 10 + 7;$}}
    \item[\textbf{B.}] Read an integer from the keyboard.
    \item[\textbf{C.}] Print data to monitor.
    \item[\textbf{D.}] Read data from a file.
    \item[\textbf{E.}] None of the others
\end{itemize}

\vspace{0.5cm}

\vspace{0.5cm}

\textbf{Question 235:} The control unit (CU) does the actual computation or processing of data.

\begin{itemize}
    \item[\textbf{A.}] True
    \item[\textbf{B.}] \textcolor{red}{\textbf{False}}
\end{itemize}

\vspace{0.5cm}

\vspace{0.5cm}

\textbf{Question 307:} Which one of four basic functions of computer describes the following statement? "The paths among components are used to move data from memory to memory and from memory through gates to memory".

\begin{itemize}
    \item[\textbf{A.}] Data storage
    \item[\textbf{B.}] Data processing
    \item[\textbf{C.}] \textcolor{red}{\textbf{Data movement}}
    \item[\textbf{D.}] Control
\end{itemize}

\vspace{0.5cm}

\vspace{0.5cm}

\textbf{Question 341:} Which of the following PDP series computers is known for its use of 12-bit instructions and a single general-purpose register, the accumulator?

\begin{itemize}
    \item[\textbf{A.}] \textcolor{red}{\textbf{PDP-8}}
    \item[\textbf{B.}] PDP-10
    \item[\textbf{C.}] PDP-11
    \item[\textbf{D.}] PDP-6
\end{itemize}

\vspace{0.5cm}

\vspace{0.5cm}

\textbf{Question 41:} Architectural attributes include \underline{\hspace{1.5cm}}.

\begin{itemize}
    \item[\textbf{A.}] memory technology used
    \item[\textbf{B.}] interfaces
    \item[\textbf{C.}] control signals
    \item[\textbf{D.}] \textcolor{red}{\textbf{I/O mechanisms}}
\end{itemize}

\vspace{0.5cm}

\vspace{0.5cm}

\textbf{Question 84:} The number of bits used to represent various data types is an example of an architectural attribute.

\begin{itemize}
    \item[\textbf{A.}] \textcolor{red}{\textbf{True}}
    \item[\textbf{B.}] False
\end{itemize}

\vspace{0.5cm}

\vspace{0.5cm}

\textbf{Question 89:} The most common means of computer/user interaction is a \underline{\hspace{1.5cm}}.

\begin{itemize}
    \item[\textbf{A.}] \textcolor{red}{\textbf{keyboard/monitor}}
    \item[\textbf{B.}] mouse/printer
    \item[\textbf{C.}] modem/printer
    \item[\textbf{D.}] monitor/printer
\end{itemize}

\vspace{0.5cm}

\vspace{0.5cm}

\textbf{Question 97:} What kind of Information is stored inside the computer?

\begin{itemize}
    \item[\textbf{A.}] \textcolor{red}{\textbf{Binary form}}
    \item[\textbf{B.}] ASCII code form
    \item[\textbf{C.}] Decimal form
    \item[\textbf{D.}] Alpha numeric
    \item[\textbf{E.}] Numeric form
\end{itemize}

\vspace{0.5cm}

\vspace{0.5cm}

\textbf{Question 101:} \underline{\hspace{1.5cm}} attributes include hardware details transparent to the programmer.

\begin{itemize}
    \item[\textbf{A.}] Interface
    \item[\textbf{B.}] \textcolor{red}{\textbf{Organizational}}
    \item[\textbf{C.}] Memory
    \item[\textbf{D.}] Architectural
\end{itemize}

\vspace{0.5cm}

\vspace{0.5cm}

\textbf{Question 114:} It is a(n) \underline{\hspace{1.5cm}} design issue whether a computer will have a multiply instruction.

\begin{itemize}
    \item[\textbf{A.}] \textcolor{red}{\textbf{architectural}}
    \item[\textbf{B.}] memory
    \item[\textbf{C.}] elementary
    \item[\textbf{D.}] organizational
\end{itemize}

\vspace{0.5cm}

\vspace{0.5cm}

\textbf{Question 120:} The \underline{\hspace{1.5cm}} performs the computer's data processing functions.

\begin{itemize}
    \item[\textbf{A.}] Register
    \item[\textbf{B.}] CPU interconnection
    \item[\textbf{C.}] \textcolor{red}{\textbf{ALU}}
    \item[\textbf{D.}] system bus
\end{itemize}

\vspace{0.5cm}

\vspace{0.5cm}

\textbf{Question 143:} Computer \underline{\hspace{1.5cm}} refers to those attributes that have a direct impact on the logical execution of a program.

\begin{itemize}
    \item[\textbf{A.}] organization
    \item[\textbf{B.}] specifics
    \item[\textbf{C.}] design
    \item[\textbf{D.}] \textcolor{red}{\textbf{architecture}}
\end{itemize}

\vspace{0.5cm}

\vspace{0.5cm}

\textbf{Question 144:} A \underline{\hspace{1.5cm}} system is a set of interrelated subsystems.

\begin{itemize}
    \item[\textbf{A.}] secondary
    \item[\textbf{B.}] \textcolor{red}{\textbf{hierarchical}}
    \item[\textbf{C.}] complex
    \item[\textbf{D.}] functional
\end{itemize}

\vspace{0.5cm}

\vspace{0.5cm}

\textbf{Question 153:} It is a(n) \underline{\hspace{1.5cm}} issue whether the multiply instruction will be implemented by a special r mechanism that makes repeated use of the add unit of the system.

\begin{itemize}
    \item[\textbf{A.}] architectural
    \item[\textbf{B.}] memory
    \item[\textbf{C.}] mechanical
    \item[\textbf{D.}] \textcolor{red}{\textbf{organizational}}
\end{itemize}

\vspace{0.5cm}

\vspace{0.5cm}

\textbf{Question 169:} Architectural attributes include \underline{\hspace{1.5cm}}.

\begin{itemize}
    \item[\textbf{A.}] \textcolor{red}{\textbf{I/O mechanisms}}
    \item[\textbf{B.}] control signals
    \item[\textbf{C.}] interfaces
    \item[\textbf{D.}] memory technology used
\end{itemize}

\vspace{0.5cm}

\vspace{0.5cm}

\textbf{Question 179:} The \underline{\hspace{1.5cm}} stores data.

\begin{itemize}
    \item[\textbf{A.}] \textcolor{red}{\textbf{main memory}}
    \item[\textbf{B.}] I/O
    \item[\textbf{C.}] system bus
    \item[\textbf{D.}] control unit
\end{itemize}

\vspace{0.5cm}

\vspace{0.5cm}

\textbf{Question 215:} \underline{\hspace{1.5cm}} attributes include hardware details transparent to the programmer.

\begin{itemize}
    \item[\textbf{A.}] \textcolor{red}{\textbf{Organizational}}
    \item[\textbf{B.}] Memory
    \item[\textbf{C.}] Architectural
    \item[\textbf{D.}] Interface
\end{itemize}

\vspace{0.5cm}

\vspace{0.5cm}

\textbf{Question 222:} Architectural attributes include \underline{\hspace{1.5cm}}.

\begin{itemize}
    \item[\textbf{A.}] interfaces
    \item[\textbf{B.}] \textcolor{red}{\textbf{I/O mechanisms}}
    \item[\textbf{C.}] control signals
    \item[\textbf{D.}] memory technology used
\end{itemize}

\vspace{0.5cm}

\vspace{0.5cm}

\textbf{Question 248:} The \underline{\hspace{1.5cm}} stores data.

\begin{itemize}
    \item[\textbf{A.}] system bus
    \item[\textbf{B.}] I/O
    \item[\textbf{C.}] \textcolor{red}{\textbf{main memory}}
    \item[\textbf{D.}] control unit
\end{itemize}

\vspace{0.5cm}

\vspace{0.5cm}

\textbf{Question 249:} Computers are classified into generations based on the fundamental hardware technology employed.

\begin{itemize}
    \item[\textbf{A.}] \textcolor{red}{\textbf{True}}
    \item[\textbf{B.}] False
\end{itemize}

\vspace{0.5cm}

\vspace{0.5cm}

\textbf{Question 259:} Central processing unit (CPU) of IAS computer consists of \underline{\hspace{1.5cm}}.

\begin{itemize}
    \item[\textbf{A.}] Main memory and ALU (arithmetic and logic unit)
    \item[\textbf{B.}] \textcolor{red}{\textbf{ALU (Arithmetic and Logic Unit) and CU (Control Unit)}}
    \item[\textbf{C.}] CU (Control Unit) and IO Module
    \item[\textbf{D.}] ALU (Arithmetic and Logic Unit) and IO Module
\end{itemize}

\vspace{0.5cm}

\vspace{0.5cm}

\textbf{Question 304:} Which of the following statements is part of the Von Newmann principle?

\begin{itemize}
    \item[\textbf{A.}] \textcolor{red}{\textbf{The computer uses a program counter to indicate the location of the next statement}}
    \item[\textbf{B.}] Computer can control all operations with a single program
    \item[\textbf{C.}] Computer memory is not addressable
    \item[\textbf{D.}] Each instruction must have a memory area containing the address of the next instruction
\end{itemize}

\vspace{0.5cm}

\vspace{0.5cm}

\textbf{Question 306:} Which of the following are the four basic functions that a computer performs?

\begin{itemize}
    \item[\textbf{A.}] \textcolor{red}{\textbf{Data processing, Data storage, Data movement, Control}}
    \item[\textbf{B.}] Data processing, Data storage, Data movement, Interrupt
    \item[\textbf{C.}] Data processing, Data storage, Interrupt, Control
    \item[\textbf{D.}] Data processing, Interrupt, Data movement, Control
\end{itemize}

\vspace{0.5cm}

\vspace{0.5cm}

\textbf{Question 313:} What is true about IAS Memory Formats?

\begin{itemize}
    \item[\textbf{A.}] The memory of the IAS consists of 1000 storage locations (called words) of 32 bits each
    \item[\textbf{B.}] Only data is stored in the memory
    \item[\textbf{C.}] \textcolor{red}{\textbf{Both data and instructions are stored in the memory}}
    \item[\textbf{D.}] Only instructions are stored in the memory
\end{itemize}

\vspace{0.5cm}

\vspace{0.5cm}

\textbf{Question 353:} Choose the most four basic functions of a computer.

\begin{itemize}
    \item[\textbf{A.}] \textcolor{red}{\textbf{Moving data, storing data, processing data, controlling.}}
    \item[\textbf{B.}] Supporting operating system, accessing hard disks, supporting network connections, supporting HDMI ports.
    \item[\textbf{C.}] Reading disks, accessing network resources.
    \item[\textbf{D.}] Reading data from keyboard, printing data to monitor, allowing network connections.
\end{itemize}

\vspace{0.5cm}

\vspace{0.5cm}

\textbf{Question 479:} Which of the following are the 4 basic functions of a computer?

\begin{itemize}
    \item[\textbf{A.}] Data processing; Data storage; Data movement; CPU.
    \item[\textbf{B.}] CPU; Main memory; I/O; System interconnection.
    \item[\textbf{C.}] Control unit; ALU; Registers; CPU interconnection.
    \item[\textbf{D.}] \textcolor{red}{\textbf{Data processing; Data storage; Data movement; Control.}}
\end{itemize}

\vspace{0.5cm}

\vspace{0.5cm}

\textbf{Question 527:} Which of the following are examples of architectural attributes in computer systems? (choose two correct answers)

\begin{itemize}
    \item[\textbf{A.}] \textcolor{red}{\textbf{Instruction set}}
    \item[\textbf{B.}] Control signals
    \item[\textbf{C.}] Memory technology
    \item[\textbf{D.}] \textcolor{red}{\textbf{Number of bits used for data types}}
    \item[\textbf{E.}] Interfaces between peripherals and the computer
\end{itemize}

\vspace{0.5cm}

\vspace{0.5cm}

\textbf{Question 232:} Which of the following is/are NOT part(s) of the CPU?

\begin{itemize}
    \item[\textbf{A.}] ALU
    \item[\textbf{B.}] The Control unit
    \item[\textbf{C.}] The Registers
    \item[\textbf{D.}] \textcolor{red}{\textbf{System bus}}
\end{itemize}

\vspace{0.5cm}

\vspace{0.5cm}

\textbf{Question 258:} Which of the following component does not belong to central processing unit?

\begin{itemize}
    \item[\textbf{A.}] \textcolor{red}{\textbf{System interconnection}}
    \item[\textbf{B.}] Arithmetic and logic unit
    \item[\textbf{C.}] Registers
    \item[\textbf{D.}] Control unit
    \item[\textbf{E.}] CPU interconnection
\end{itemize}

\vspace{0.5cm}

\vspace{0.5cm}

\section{Chapter 02: Performance Concepts}
\textit{Total questions in chapter: 12}

\textbf{Question 7:} The use of multiple processors on the same chip is referred to as \underline{\hspace{1.5cm}} and provides the potential to increase performance without increasing the clock rate.

\begin{itemize}
    \item[\textbf{A.}] \textcolor{red}{\textbf{multicore}}
    \item[\textbf{B.}] GPU
    \item[\textbf{C.}] data channels
    \item[\textbf{D.}] MPC
\end{itemize}

\vspace{0.5cm}

\vspace{0.5cm}

\textbf{Question 12:} It is possible to improve pipeline performance by automatically rearranging instructions within a program so that branch instructions occur later than actually desired.

\begin{itemize}
    \item[\textbf{A.}] \textcolor{red}{\textbf{True}}
    \item[\textbf{B.}] False
\end{itemize}

\vspace{0.5cm}

\vspace{0.5cm}

\textbf{Question 88:} Designers wrestle with the challenge of balancing processor performance with that of main memory and other computer components.

\begin{itemize}
    \item[\textbf{A.}] \textcolor{red}{\textbf{True}}
    \item[\textbf{B.}] False
\end{itemize}

\vspace{0.5cm}

\vspace{0.5cm}

\textbf{Question 301:} Follow the Amdahl's law for multiprocessors, if only 10\% of the code is inherently serial (f = 0.9), running the program on a multicore system with 4 processors, a performance gain (speedup factor) would be \underline{\hspace{1.5cm}}.

\begin{itemize}
    \item[\textbf{A.}] \textcolor{red}{\textbf{307\%}}
    \item[\textbf{B.}] 297\%
    \item[\textbf{C.}] 317\%
    \item[\textbf{D.}] 327\%
\end{itemize}

\vspace{0.5cm}

\vspace{0.5cm}

\textbf{Question 319:} Why is cache design used in high-performance computing (HPC)? (Choose three correct answers)

\begin{itemize}
    \item[\textbf{A.}] \textcolor{red}{\textbf{Because there is a significant speed gap between the processor and the internal memory in HPC.}}
    \item[\textbf{B.}] \textcolor{red}{\textbf{Because applications in HPC often require a large bandwidth to support intensive data processing}}
    \item[\textbf{C.}] Because power consumption can be a significant operational cost in HPC environments.
    \item[\textbf{D.}] \textcolor{red}{\textbf{Because multiple processors are often working in parallel, caches provide a way to efficiently manage data required by these processors}}
\end{itemize}

\vspace{0.5cm}

\vspace{0.5cm}

\textbf{Question 391:} A benchmark program is running on a 400 MHz processor. The executed program consists of 500 instruction executions, with the following instruction mix and clock cycle count in the table below:

\begin{center}
\renewcommand{\arraystretch}{1.3}
\begin{tabular}{|l|c|c|}
\hline
\textbf{Instruction type} & \textbf{Instruction count} & \textbf{Cycles per Instruction} \\ \hline
Arithmetic & 300 & 1 \\ \hline
Data transfer & 100 & 2 \\ \hline
Control transfer & 100 & 2 \\ \hline
\end{tabular}
\end{center}

Calculate MIPS rate for this case.

\begin{itemize}
    \item[\textbf{A.}] \textcolor{red}{\textbf{MIPS rate = 285.7}}
    \item[\textbf{B.}] MIPS rate = 275.7
    \item[\textbf{C.}] MIPS rate = 265.7
    \item[\textbf{D.}] MIPS rate = 295.7
\end{itemize}

\vspace{0.5cm}

\textbf{Question 478:} What is one way to control power density in microprocessor chips?

\begin{itemize}
    \item[\textbf{A.}] Increase the clock speed of the chip.
    \item[\textbf{B.}] \textcolor{red}{\textbf{Use more of the chip area for cache memory.}}
    \item[\textbf{C.}] Reduce the size of the chip die.
    \item[\textbf{D.}] Use more transistors for logic operations.
\end{itemize}

\vspace{0.5cm}

\vspace{0.5cm}

\textbf{Question 505:} Which of the following types of processors typically offer better performance for parallel processing?

\begin{itemize}
    \item[\textbf{A.}] Single-core processors
    \item[\textbf{B.}] \textcolor{red}{\textbf{Multi-core processors}}
    \item[\textbf{C.}] Microelectronics
    \item[\textbf{D.}] Digital signal processors (DSPs)
\end{itemize}

\vspace{0.5cm}

\vspace{0.5cm}

\textbf{Question 506:} Follow the Amdahl's law for multiprocessors, if only 20\% of the code is inherently serial (f = 0.8), running the program on a multicore system with 8 processors, a performance gain (speedup factor) would be \underline{\hspace{1.5cm}}.

\begin{itemize}
    \item[\textbf{A.}] \textcolor{red}{\textbf{333\%}}
    \item[\textbf{B.}] 303\%
    \item[\textbf{C.}] 313\%
    \item[\textbf{D.}] 323\%
\end{itemize}

\vspace{0.5cm}

\vspace{0.5cm}

\textbf{Question 138:} The memory transfer rate has not kept up with increases in processor speed.

\begin{itemize}
    \item[\textbf{A.}] \textcolor{red}{\textbf{True}}
    \item[\textbf{B.}] False
\end{itemize}

\vspace{0.5cm}

\vspace{0.5cm}

\textbf{Question 310:} The pulse rate in the ck system is known as the \underline{\hspace{1.5cm}}.

\begin{itemize}
    \item[\textbf{A.}] ck cycle
    \item[\textbf{B.}] \textcolor{red}{\textbf{ck speed}}
    \item[\textbf{C.}] ck time
    \item[\textbf{D.}] ck tick
\end{itemize}

\vspace{0.5cm}

\vspace{0.5cm}

\textbf{Question 507:} Choose an INCORRECT trend in contemporary processor designs.

\begin{itemize}
    \item[\textbf{A.}] \textcolor{red}{\textbf{The number of processes accepted is infinitive.}}
    \item[\textbf{B.}] The number of integrated transistors is getting increasing and more larger.
    \item[\textbf{C.}] Higher clock frequency is applied.
    \item[\textbf{D.}] More cores are integrated in one chip.
\end{itemize}

\vspace{0.5cm}

\vspace{0.5cm}

\section{Chapter 03: A Top-Level View of Computer Function and Interconnection}
\textit{Total questions in chapter: 51}

\textbf{Question 38:} A bus that connects major computer components (processor, memory, I/O) is called a \underline{\hspace{1.5cm}}

\begin{itemize}
    \item[\textbf{A.}] control bus
    \item[\textbf{B.}] \textcolor{red}{\textbf{system bus}}
    \item[\textbf{C.}] data bus
    \item[\textbf{D.}] address bus
\end{itemize}

\vspace{0.5cm}

\vspace{0.5cm}

\textbf{Question 39:} A common example of system interconnection is by means of a \underline{\hspace{1.5cm}}.

\begin{itemize}
    \item[\textbf{A.}] data transport
    \item[\textbf{B.}] control device
    \item[\textbf{C.}] register
    \item[\textbf{D.}] \textcolor{red}{\textbf{system bus}}
\end{itemize}

\vspace{0.5cm}

\vspace{0.5cm}

\textbf{Question 49:} The routine executed in response to an interrupt request is called \underline{\hspace{1.5cm}} routine.

\begin{itemize}
    \item[\textbf{A.}] \textcolor{red}{\textbf{Interrupt Service}}
    \item[\textbf{B.}] Interrupt acknowledge
    \item[\textbf{C.}] Serial interrupt
    \item[\textbf{D.}] Vectored interrupt
    \item[\textbf{E.}] Sub-routine
\end{itemize}

\vspace{0.5cm}

\vspace{0.5cm}

\textbf{Question 57:} The data lines provide a path for moving data among system modules and are collectively called the \underline{\hspace{1.5cm}}.

\begin{itemize}
    \item[\textbf{A.}] control bus
    \item[\textbf{B.}] system bus
    \item[\textbf{C.}] address bus
    \item[\textbf{D.}] \textcolor{red}{\textbf{data bus}}
\end{itemize}

\vspace{0.5cm}

\vspace{0.5cm}

\textbf{Question 76:} The System bus is made up of \underline{\hspace{1.5cm}}

\begin{itemize}
    \item[\textbf{A.}] Control bus
    \item[\textbf{B.}] Address bus
    \item[\textbf{C.}] Both Control bus and Address bus
    \item[\textbf{D.}] \textcolor{red}{\textbf{Control bus, Data bus and Address bus}}
    \item[\textbf{E.}] Data bus
\end{itemize}

\vspace{0.5cm}

\vspace{0.5cm}

\textbf{Question 107:} During the \underline{\hspace{1.5cm}} the opcode of the next instruction is loaded into the IR and the addres into the MAR.

\begin{itemize}
    \item[\textbf{A.}] execute cycle
    \item[\textbf{B.}] \textcolor{red}{\textbf{fetch cycle}}
    \item[\textbf{C.}] instruction cycle
    \item[\textbf{D.}] clock cycle
\end{itemize}

\vspace{0.5cm}

\vspace{0.5cm}

\textbf{Question 110:} Interrupt processing allows an application program to be suspended in order that a variety of conditions can be serviced and later resumed.

\begin{itemize}
    \item[\textbf{A.}] \textcolor{red}{\textbf{True}}
    \item[\textbf{B.}] False
\end{itemize}

\vspace{0.5cm}

\vspace{0.5cm}

\textbf{Question 113:} The \underline{\hspace{1.5cm}} is connected to the data lines of the system bus.

\begin{itemize}
    \item[\textbf{A.}] MAR
    \item[\textbf{B.}] PC
    \item[\textbf{C.}] \textcolor{red}{\textbf{MBR}}
    \item[\textbf{D.}] IR
\end{itemize}

\vspace{0.5cm}

\vspace{0.5cm}

\textbf{Question 116:} The interconnection structure must support which transfer?

\begin{itemize}
    \item[\textbf{A.}] memory to processor
    \item[\textbf{B.}] processor to memory
    \item[\textbf{C.}] I/O to or from memory
    \item[\textbf{D.}] \textcolor{red}{\textbf{all of the above}}
\end{itemize}

\vspace{0.5cm}

\vspace{0.5cm}

\textbf{Question 124:} The \underline{\hspace{1.5cm}} are used to designate the source or destination of the data on the data bus.

\begin{itemize}
    \item[\textbf{A.}] system lines
    \item[\textbf{B.}] data lines
    \item[\textbf{C.}] control lines
    \item[\textbf{D.}] \textcolor{red}{\textbf{address lines}}
\end{itemize}

\vspace{0.5cm}

\vspace{0.5cm}

\textbf{Question 134:} Timing refers to the way in which events are coordinated on the bus.

\begin{itemize}
    \item[\textbf{A.}] \textcolor{red}{\textbf{True}}
    \item[\textbf{B.}] False
\end{itemize}

\vspace{0.5cm}

\vspace{0.5cm}

\textbf{Question 137:} An interrupt is generated from software and it is provoked by the execution of an instruction.

\begin{itemize}
    \item[\textbf{A.}] True
    \item[\textbf{B.}] \textcolor{red}{\textbf{False}}
\end{itemize}

\vspace{0.5cm}

\vspace{0.5cm}

\textbf{Question 172:} An interrupt is a hardware-generated signal to the processor.

\begin{itemize}
    \item[\textbf{A.}] \textcolor{red}{\textbf{True}}
    \item[\textbf{B.}] False
\end{itemize}

\vspace{0.5cm}

\vspace{0.5cm}

\textbf{Question 178:} The \underline{\hspace{1.5cm}} is connected to the address lines of the system bus.

\begin{itemize}
    \item[\textbf{A.}] MBR
    \item[\textbf{B.}] \textcolor{red}{\textbf{MAR}}
    \item[\textbf{C.}] PC
    \item[\textbf{D.}] IR
\end{itemize}

\vspace{0.5cm}

\vspace{0.5cm}

\textbf{Question 188:} With asynchronous timing the occurrence of events on the bus is determined by a clock.

\begin{itemize}
    \item[\textbf{A.}] \textcolor{red}{\textbf{False}}
    \item[\textbf{B.}] True
\end{itemize}

\vspace{0.5cm}

\vspace{0.5cm}

\textbf{Question 198:} An interrupt is generated from software and it is provoked by the execution of an instruction.

\begin{itemize}
    \item[\textbf{A.}] \textcolor{red}{\textbf{False}}
    \item[\textbf{B.}] True
\end{itemize}

\vspace{0.5cm}

\vspace{0.5cm}

\textbf{Question 201:} The System bus is made up of \underline{\hspace{1.5cm}}

\begin{itemize}
    \item[\textbf{A.}] Address bus
    \item[\textbf{B.}] Control bus
    \item[\textbf{C.}] \textcolor{red}{\textbf{Control bus, Data bus and Address bus}}
    \item[\textbf{D.}] Both Control bus
\end{itemize}

\vspace{0.5cm}

\vspace{0.5cm}

\textbf{Question 203:} During the \underline{\hspace{1.5cm}} the opcode of the next instruction is loaded into the IR and the address portion is loaded into the MAR.

\begin{itemize}
    \item[\textbf{A.}] \textcolor{red}{\textbf{fetch cycle}}
    \item[\textbf{B.}] clock cycle
    \item[\textbf{C.}] instruction cycle
    \item[\textbf{D.}] execute cycle
\end{itemize}

\vspace{0.5cm}

\vspace{0.5cm}

\textbf{Question 204:} The \underline{\hspace{1.5cm}} are used to designate the source or destination of the data on the data bus.

\begin{itemize}
    \item[\textbf{A.}] system lines
    \item[\textbf{B.}] data lines
    \item[\textbf{C.}] control lines
    \item[\textbf{D.}] \textcolor{red}{\textbf{address lines}}
\end{itemize}

\vspace{0.5cm}

\vspace{0.5cm}

\textbf{Question 216:} The data lines provide a path for moving data among system modules and are collectively called the \underline{\hspace{1.5cm}}.

\begin{itemize}
    \item[\textbf{A.}] system bus
    \item[\textbf{B.}] address bus
    \item[\textbf{C.}] control bus
    \item[\textbf{D.}] \textcolor{red}{\textbf{data bus}}
\end{itemize}

\vspace{0.5cm}

\vspace{0.5cm}

\textbf{Question 262:} Which of the following determines the Bus Width?

\begin{itemize}
    \item[\textbf{A.}] The clock speed of the CPU
    \item[\textbf{B.}] The number of cores in the processor
    \item[\textbf{C.}] The size of the motherboard
    \item[\textbf{D.}] \textcolor{red}{\textbf{The number of parallel lines in the data bus}}
    \item[\textbf{E.}] Number of components connected to Bus
\end{itemize}

\vspace{0.5cm}

\vspace{0.5cm}

\textbf{Question 303:} What is the significance of the program counter (PC) in the fetch phase of the instruction cycle?

\begin{itemize}
    \item[\textbf{A.}] The program counter (PC) is not used in the fetch phase, and its role is limited to tracking the number of instructions executed by the CPU
    \item[\textbf{B.}] \textcolor{red}{\textbf{The program counter (PC) in the fetch phase holds the memory address of the next instruction to be fetched and executed}}
    \item[\textbf{C.}] The program counter (PC) is responsible for executing instructions and has no specific role during the fetch phase
    \item[\textbf{D.}] The program counter (PC) is only relevant in multi-core processors and does not contribute to the fetch phase of the instruction cycle in single-core systems
\end{itemize}

\vspace{0.5cm}

\vspace{0.5cm}

\textbf{Question 314:} What is the most important characteristic of the Synchronous Bus?

\begin{itemize}
    \item[\textbf{A.}] Data is transmitted at the same time
    \item[\textbf{B.}] The occurrence of one event on a bus follows and depends on the occurrence of a previous event.
    \item[\textbf{C.}] \textcolor{red}{\textbf{The occurrence of events on the bus is determined by a clock}}
    \item[\textbf{D.}] No common clock signal controlling operation
\end{itemize}

\vspace{0.5cm}

\vspace{0.5cm}

\textbf{Question 368:} What is an interrupt vector?

\begin{itemize}
    \item[\textbf{A.}] \textcolor{red}{\textbf{Part of memory which contains the addresses of interrupt handlers}}
    \item[\textbf{B.}] A signal an I/O device sends to CPU
    \item[\textbf{C.}] A signal an I/O software sends to CPU
    \item[\textbf{D.}] None of the mentioned
\end{itemize}

\vspace{0.5cm}

\vspace{0.5cm}

\textbf{Question 382:} What is the function of the bus system in the computer?

\begin{itemize}
    \item[\textbf{A.}] Extend the communication function of the computer
    \item[\textbf{B.}] \textcolor{red}{\textbf{Connect components in the computer}}
    \item[\textbf{C.}] Control peripherals
    \item[\textbf{D.}] Transform signals in the computer
    \item[\textbf{E.}] All of the mentioned
\end{itemize}

\vspace{0.5cm}

\vspace{0.5cm}

\textbf{Question 439:} The routine executed in response to an interrupt request is called \underline{\hspace{1.5cm}} routine.

\begin{itemize}
    \item[\textbf{A.}] Vectored interrupt
    \item[\textbf{B.}] Sub-routine
    \item[\textbf{C.}] Serial interrupt
    \item[\textbf{D.}] Interrupt acknowledge
    \item[\textbf{E.}] \textcolor{red}{\textbf{Interrupt Service}}
\end{itemize}

\vspace{0.5cm}

\vspace{0.5cm}

\textbf{Question 460:} What is the role of the memory address register (MAR) in a computer system?

\begin{itemize}
    \item[\textbf{A.}] It holds data being transferred between memory and the processor.
    \item[\textbf{B.}] \textcolor{red}{\textbf{It specifies the address in memory for the next read or write operation.}}
    \item[\textbf{C.}] It performs arithmetic operations on data stored in memory.
    \item[\textbf{D.}] It stores program instructions after they are executed.
\end{itemize}

\vspace{0.5cm}

\vspace{0.5cm}

\textbf{Question 481:} The instruction, which adds 1 to the value in a memory location, has five stages: fetch opcode (2 bus cycles), fetch operand address (2 bus cycles), fetch operand (2 bus cycles), add 1 to operand (1 bus cycle), and store operand (3 bus cycles). An interrupt sends request at beginning of fetch operand address stage. How many cycles does the processor enter the interrupt processing cycle?

\begin{itemize}
    \item[\textbf{A.}] 6
    \item[\textbf{B.}] 7
    \item[\textbf{C.}] \textcolor{red}{\textbf{8}}
    \item[\textbf{D.}] 9
    \item[\textbf{E.}] 5
\end{itemize}

\vspace{0.5cm}

\vspace{0.5cm}

\textbf{Question 482:} Which type of interrupt is exemplified by "division by zero, attempt to execute an illegal machine instruction"?

\begin{itemize}
    \item[\textbf{A.}] \textcolor{red}{\textbf{Program}}
    \item[\textbf{B.}] Timer
    \item[\textbf{C.}] I/O
    \item[\textbf{D.}] Hardware failure
\end{itemize}

\vspace{0.5cm}

\vspace{0.5cm}

\textbf{Question 483:} Which type is not an Interconnection bus?

\begin{itemize}
    \item[\textbf{A.}] Data lines
    \item[\textbf{B.}] Control lines
    \item[\textbf{C.}] \textcolor{red}{\textbf{Signal lines}}
    \item[\textbf{D.}] Address lines
\end{itemize}

\vspace{0.5cm}

\vspace{0.5cm}

\textbf{Question 529:} When an interrupt is about to execute next, where is the Program Counter (PC) positioned?

\begin{itemize}
    \item[\textbf{A.}] Add 1 to PC
    \item[\textbf{B.}] At the address of the last executed instruction
    \item[\textbf{C.}] Jump to random unknown position
    \item[\textbf{D.}] \textcolor{red}{\textbf{At the started address of an interrupt handler routine}}
    \item[\textbf{E.}] At the beginning of the memory address space
\end{itemize}

\vspace{0.5cm}

\vspace{0.5cm}

\textbf{Question 63:} A(n) \underline{\hspace{1.5cm}} is generated by a failure such as power failure or memory parity error.

\begin{itemize}
    \item[\textbf{A.}] timer interrupt
    \item[\textbf{B.}] \textcolor{red}{\textbf{hardware failure interrupt}}
    \item[\textbf{C.}] I/O interrupt
    \item[\textbf{D.}] program interrupt
\end{itemize}

\vspace{0.5cm}

\vspace{0.5cm}

\textbf{Question 181:} Program execution consists of repeating the process of instruction fetch and instruction execution

\begin{itemize}
    \item[\textbf{A.}] False
    \item[\textbf{B.}] \textcolor{red}{\textbf{True}}
\end{itemize}

\vspace{0.5cm}

\vspace{0.5cm}

\textbf{Question 15:} The \underline{\hspace{1.5cm}} cycle occurs at the beginning of each instruction cycle and causes an instruction to be fetched from memory.

\begin{itemize}
    \item[\textbf{A.}] execute
    \item[\textbf{B.}] indirect
    \item[\textbf{C.}] \textcolor{red}{\textbf{fetch}}
    \item[\textbf{D.}] interrupt
\end{itemize}

\vspace{0.5cm}

\vspace{0.5cm}

\textbf{Question 70:} Interrupts do not improve processing efficiency.

\begin{itemize}
    \item[\textbf{A.}] True
    \item[\textbf{B.}] \textcolor{red}{\textbf{False}}
\end{itemize}

\vspace{0.5cm}

\vspace{0.5cm}

\textbf{Question 225:} The \underline{\hspace{1.5cm}} contains a word of data to be written to memory or the word most recently read.

\begin{itemize}
    \item[\textbf{A.}] \textcolor{red}{\textbf{MBR}}
    \item[\textbf{B.}] IR
    \item[\textbf{C.}] PC
    \item[\textbf{D.}] MAR
\end{itemize}

\vspace{0.5cm}

\vspace{0.5cm}

\textbf{Question 1:} A sequence of codes or instructions is called \underline{\hspace{1.5cm}}.

\begin{itemize}
    \item[\textbf{A.}] \textcolor{red}{\textbf{software}}
    \item[\textbf{B.}] memory
    \item[\textbf{C.}] an interconnect
    \item[\textbf{D.}] a register
\end{itemize}

\vspace{0.5cm}

\vspace{0.5cm}

\textbf{Question 458:} Choose the appropriate words to fill in the corresponding blanks in the following definitions:
\underline{\hspace{1.5cm}} contains a word to be stored in memory or sent to the I/O unit, or is used to receive a word from memory or from the I/O unit. \underline{\hspace{1.5cm}} specifies the address in memory of the word to be written from or read into the MBR. \underline{\hspace{1.5cm}} is employed to hold temporarily the right hand instruction from a word in memory.

\begin{itemize}
    \item[\textbf{A.}] \textcolor{red}{\textbf{MBR; MAR; IBR}}
    \item[\textbf{B.}] AC; MAR; IR
    \item[\textbf{C.}] MAR; MBR; IBR
    \item[\textbf{D.}] AC; MBR; MAR
\end{itemize}

\vspace{0.5cm}

\vspace{0.5cm}

\textbf{Question 459:} The stored-program concept means a computer could get its instructions by reading them from memory, and a program could be set or altered by setting the values of a portion of memory.

\begin{itemize}
    \item[\textbf{A.}] \textcolor{red}{\textbf{True}}
    \item[\textbf{B.}] False
\end{itemize}

\vspace{0.5cm}

\vspace{0.5cm}

\textbf{Question 244:} Because all devices on a synchronous bus are tied to a fixed clock rate, the system cannot take advantage of advances in device performance.

\begin{itemize}
    \item[\textbf{A.}] \textcolor{red}{\textbf{True}}
    \item[\textbf{B.}] False
\end{itemize}

\vspace{0.5cm}

\vspace{0.5cm}

\textbf{Question 46:} The method of using the same lines for multiple purposes is known as time multiplexing.

\begin{itemize}
    \item[\textbf{A.}] \textcolor{red}{\textbf{True}}
    \item[\textbf{B.}] False
\end{itemize}

\vspace{0.5cm}

\vspace{0.5cm}

\textbf{Question 56:} The basic function of a computer is to execute programs.

\begin{itemize}
    \item[\textbf{A.}] False
    \item[\textbf{B.}] \textcolor{red}{\textbf{True}}
\end{itemize}

\vspace{0.5cm}

\vspace{0.5cm}

\textbf{Question 69:} The \underline{\hspace{1.5cm}} performs the computer's data processing functions.

\begin{itemize}
    \item[\textbf{A.}] system bus
    \item[\textbf{B.}] Register
    \item[\textbf{C.}] \textcolor{red}{\textbf{ALU}}
    \item[\textbf{D.}] CPU interconnection
\end{itemize}

\vspace{0.5cm}

\vspace{0.5cm}

\textbf{Question 98:} A(n) \underline{\hspace{1.5cm}} is generated by some condition that occurs as a result of an instruction exec

\begin{itemize}
    \item[\textbf{A.}] timer interrupt
    \item[\textbf{B.}] I/O interrupt
    \item[\textbf{C.}] \textcolor{red}{\textbf{program interrupt}}
    \item[\textbf{D.}] hardware failure interrupt
\end{itemize}

\vspace{0.5cm}

\vspace{0.5cm}

\textbf{Question 106:} A \underline{\hspace{1.5cm}} is a mechanism that provides for communication among CPU, main memory, i

\begin{itemize}
    \item[\textbf{A.}] \textcolor{red}{\textbf{system interconnection}}
    \item[\textbf{B.}] CPU interconnection
    \item[\textbf{C.}] peripheral
    \item[\textbf{D.}] processor
\end{itemize}

\vspace{0.5cm}

\vspace{0.5cm}

\textbf{Question 148:} The \underline{\hspace{1.5cm}} contains the 8-bit opcode instruction being executed.

\begin{itemize}
    \item[\textbf{A.}] memory buffer register
    \item[\textbf{B.}] instruction buffer register
    \item[\textbf{C.}] \textcolor{red}{\textbf{instruction register}}
    \item[\textbf{D.}] memory address register
\end{itemize}

\vspace{0.5cm}

\vspace{0.5cm}

\textbf{Question 224:} The \underline{\hspace{1.5cm}} performs the computer's data processing functions.

\begin{itemize}
    \item[\textbf{A.}] \textcolor{red}{\textbf{ALU}}
    \item[\textbf{B.}] Register
    \item[\textbf{C.}] system bus
    \item[\textbf{D.}] CPU interconnection
\end{itemize}

\vspace{0.5cm}

\vspace{0.5cm}

\textbf{Question 311:} Which register is the memory address register?

\begin{itemize}
    \item[\textbf{A.}] \textcolor{red}{\textbf{MAR}}
    \item[\textbf{B.}] MBR
    \item[\textbf{C.}] IR
    \item[\textbf{D.}] PC
\end{itemize}

\vspace{0.5cm}

\vspace{0.5cm}

\textbf{Question 374:} What is the primary purpose of the "Fetch instruction" phase in the operation of a processor?

\begin{itemize}
    \item[\textbf{A.}] \textcolor{red}{\textbf{To read an instruction from memory}}
    \item[\textbf{B.}] To interpret the instruction
    \item[\textbf{C.}] To perform arithmetic operations on data
    \item[\textbf{D.}] To write data to memory
\end{itemize}

\vspace{0.5cm}

\vspace{0.5cm}

\textbf{Question 394:} Performs the computer's data processing functions is \underline{\hspace{1.5cm}}?

\begin{itemize}
    \item[\textbf{A.}] \textcolor{red}{\textbf{Arithmetic and logic unit (ALU)}}
    \item[\textbf{B.}] Control unit
    \item[\textbf{C.}] CPU interconnection
    \item[\textbf{D.}] Registers
\end{itemize}

\vspace{0.5cm}

\vspace{0.5cm}

\textbf{Question 490:} \underline{\hspace{1.5cm}}: Read instruction from its memory location into the processor.

\begin{itemize}
    \item[\textbf{A.}] \textcolor{red}{\textbf{Instruction fetch}}
    \item[\textbf{B.}] Instruction operation decoding
    \item[\textbf{C.}] Operand address calculation
    \item[\textbf{D.}] Data operation
\end{itemize}

\vspace{0.5cm}

\vspace{0.5cm}

\section{Chapter 04: The Memory Hierarchy Locality and Performance}
\textit{Total questions in chapter: 18}

\textbf{Question 123:} Within the processor there is a set of registers that function as a level of memory above main memory and cache in the hierarchy.

\begin{itemize}
    \item[\textbf{A.}] \textcolor{red}{\textbf{True}}
    \item[\textbf{B.}] False
\end{itemize}

\vspace{0.5cm}

\vspace{0.5cm}

\textbf{Question 80:} Cache memory is a much faster memory than the register file.

\begin{itemize}
    \item[\textbf{A.}] True
    \item[\textbf{B.}] \textcolor{red}{\textbf{False}}
\end{itemize}

\vspace{0.5cm}

\vspace{0.5cm}

\textbf{Question 118:} Cache is not a form of internal memory.

\begin{itemize}
    \item[\textbf{A.}] True
    \item[\textbf{B.}] \textcolor{red}{\textbf{False}}
\end{itemize}

\vspace{0.5cm}

\vspace{0.5cm}

\textbf{Question 397:} Which method of accessing units of data is used to copy a block in main memory into cache memory?

\begin{itemize}
    \item[\textbf{A.}] Direct access
    \item[\textbf{B.}] Random access
    \item[\textbf{C.}] Sequential access
    \item[\textbf{D.}] \textcolor{red}{\textbf{Associative}}
\end{itemize}

\vspace{0.5cm}

\vspace{0.5cm}

\textbf{Question 132:} For internal memory, the \underline{\hspace{1.5cm}} is equal to the number of electrical lines into and out of the memory module.

\begin{itemize}
    \item[\textbf{A.}] access time
    \item[\textbf{B.}] \textcolor{red}{\textbf{unit of transfer}}
    \item[\textbf{C.}] capacity
    \item[\textbf{D.}] memory ratio
\end{itemize}

\vspace{0.5cm}

\vspace{0.5cm}

\textbf{Question 5:} individual blocks or records have a unique address based on physical location with \underline{\hspace{1.5cm}}.

\begin{itemize}
    \item[\textbf{A.}] associative
    \item[\textbf{B.}] physical access
    \item[\textbf{C.}] \textcolor{red}{\textbf{direct access}}
    \item[\textbf{D.}] sequential access
\end{itemize}

\vspace{0.5cm}

\vspace{0.5cm}

\textbf{Question 190:} "Memory is organized into records and access must be made in a specific linear sequence" is \underline{\hspace{1.5cm}}.

\begin{itemize}
    \item[\textbf{A.}] \textcolor{red}{\textbf{sequential access}}
    \item[\textbf{B.}] direct access
    \item[\textbf{C.}] random access
    \item[\textbf{D.}] associative
\end{itemize}

\vspace{0.5cm}

\vspace{0.5cm}

\textbf{Question 238:} \underline{\hspace{1.5cm}} refers to whether memory is internal or external to the computer.

\begin{itemize}
    \item[\textbf{A.}] \textcolor{red}{\textbf{Location}}
    \item[\textbf{B.}] Access
    \item[\textbf{C.}] Hierarchy
    \item[\textbf{D.}] Tag
\end{itemize}

\vspace{0.5cm}

\vspace{0.5cm}

\textbf{Question 272:} Which memory has the fastest speed and smallest capacity?

\begin{itemize}
    \item[\textbf{A.}] \textcolor{red}{\textbf{Cache}}
    \item[\textbf{B.}] Main memory
    \item[\textbf{C.}] HDD
    \item[\textbf{D.}] Magnetic Disk
\end{itemize}

\vspace{0.5cm}

\vspace{0.5cm}

\textbf{Question 59:} The performance of the cache memory is measured in terms of a quantity called \underline{\hspace{1.5cm}}

\begin{itemize}
    \item[\textbf{A.}] \textcolor{red}{\textbf{Hit Ratio}}
    \item[\textbf{B.}] Instruction Ratio
    \item[\textbf{C.}] Miss Ratio
    \item[\textbf{D.}] Initialization Ratio
    \item[\textbf{E.}] Address Ratio
\end{itemize}

\vspace{0.5cm}

\vspace{0.5cm}

\textbf{Question 322:} What is correct about increasing performance and endurance?

\begin{itemize}
    \item[\textbf{A.}] Hard Disk- DRAM - NAND Flash - SRAM
    \item[\textbf{B.}] \textcolor{red}{\textbf{Hard Disk - NAND Flash - DRAM - SRAM}}
    \item[\textbf{C.}] NAND Flash - Hard Disk - SRAM - DRAM
    \item[\textbf{D.}] Hard Disk - DRAM - NAND Flash - SRAM
\end{itemize}

\vspace{0.5cm}

\vspace{0.5cm}

\textbf{Question 58:} The \underline{\hspace{1.5cm}} consists of the access time plus any additional time required before a second commence.

\begin{itemize}
    \item[\textbf{A.}] direct access
    \item[\textbf{B.}] transfer rate
    \item[\textbf{C.}] \textcolor{red}{\textbf{memory cycle time}}
    \item[\textbf{D.}] latency
\end{itemize}

\vspace{0.5cm}

\vspace{0.5cm}

\textbf{Question 195:} The \underline{\hspace{1.5cm}} consists of the access time plus any additional time required before a second access can commence.

\begin{itemize}
    \item[\textbf{A.}] latency
    \item[\textbf{B.}] \textcolor{red}{\textbf{memory cycle time}}
    \item[\textbf{C.}] transfer rate
    \item[\textbf{D.}] direct access
\end{itemize}

\vspace{0.5cm}

\vspace{0.5cm}

\textbf{Question 265:} What is the correct order of memory access speed from fastest to slowest?

\begin{itemize}
    \item[\textbf{A.}] \textcolor{red}{\textbf{Registers > Cache > RAM > SSD}}
    \item[\textbf{B.}] Cache > Registers > RAM > SSD
    \item[\textbf{C.}] Registers > Cache > SSD > RAM
    \item[\textbf{D.}] Cache > Registers > SSD > RAM
    \item[\textbf{E.}] All of the mentioned are wrong
\end{itemize}

\vspace{0.5cm}

\vspace{0.5cm}

\textbf{Question 274:} Sort the following memory types in ascending order by access speed:

\begin{itemize}
    \item[\textbf{A.}] \textcolor{red}{\textbf{HDD -Main Memory - L2 cache - L1 cache}}
    \item[\textbf{B.}] HDD - Main Memory - L1 cache - L2 cache
    \item[\textbf{C.}] HDD - L2 cache - L1 cache - Main Memory
    \item[\textbf{D.}] Main Memory - L2 cache - L1 cache - HDD
\end{itemize}

\vspace{0.5cm}

\vspace{0.5cm}

\textbf{Question 362:} A byte addressable microprocessor has 24 bit address. What is maximum memory capacity?

\begin{itemize}
    \item[\textbf{A.}] 4 MegaByte
    \item[\textbf{B.}] 8 MegaByte
    \item[\textbf{C.}] \textcolor{red}{\textbf{16 MegaByte}}
    \item[\textbf{D.}] 32 MegaByte
\end{itemize}

\vspace{0.5cm}

\vspace{0.5cm}

\textbf{Question 399:} In modern computers, which devices among the memory types usually have the smallest capacity?

\begin{itemize}
    \item[\textbf{A.}] RAM
    \item[\textbf{B.}] Magnetic Tape
    \item[\textbf{C.}] Hard Disk Drive
    \item[\textbf{D.}] \textcolor{red}{\textbf{Cache Memory}}
\end{itemize}

\vspace{0.5cm}

\vspace{0.5cm}

\textbf{Question 411:} What methods can be used to access units of data in memory?

\begin{itemize}
    \item[\textbf{A.}] Indirect access, Direct access, Random access, Associative
    \item[\textbf{B.}] Sequential access, Indirect access, Random access, Associative
    \item[\textbf{C.}] \textcolor{red}{\textbf{Sequential access, Direct access, Random access, Associative}}
    \item[\textbf{D.}] Sequential access, Direct access, Random access, Indirect access
\end{itemize}

\vspace{0.5cm}

\vspace{0.5cm}

\section{Chapter 05: Cache Memory}
\textit{Total questions in chapter: 25}

\textbf{Question 8:} A line includes a \underline{\hspace{1.5cm}} that identifies which particular block is currently being stored.

\begin{itemize}
    \item[\textbf{A.}] cache
    \item[\textbf{B.}] hit
    \item[\textbf{C.}] \textcolor{red}{\textbf{tag}}
    \item[\textbf{D.}] locality
\end{itemize}

\vspace{0.5cm}

\vspace{0.5cm}

\textbf{Question 14:} \underline{\hspace{1.5cm}} is the simplest mapping technique and maps each block of main memory into only one possible cache line.

\begin{itemize}
    \item[\textbf{A.}] \textcolor{red}{\textbf{Direct mapping}}
    \item[\textbf{B.}] Associative mapping
    \item[\textbf{C.}] Set associative mapping
    \item[\textbf{D.}] None of the above
\end{itemize}

\vspace{0.5cm}

\vspace{0.5cm}

\textbf{Question 29:} The cache is capable of handling global as well as local variables.

\begin{itemize}
    \item[\textbf{A.}] False
    \item[\textbf{B.}] \textcolor{red}{\textbf{True}}
\end{itemize}

\vspace{0.5cm}

\vspace{0.5cm}

\textbf{Question 50:} The correspondence between the main memory blocks and those in the cache is specified by

\begin{itemize}
    \item[\textbf{A.}] Segment function
    \item[\textbf{B.}] Hit rate
    \item[\textbf{C.}] Replacement algorithm
    \item[\textbf{D.}] \textcolor{red}{\textbf{Mapping function}}
    \item[\textbf{E.}] Miss penalty
\end{itemize}

\vspace{0.5cm}

\vspace{0.5cm}

\textbf{Question 115:} The L1 cache is slower than the L3 cache.

\begin{itemize}
    \item[\textbf{A.}] True
    \item[\textbf{B.}] \textcolor{red}{\textbf{False}}
\end{itemize}

\vspace{0.5cm}

\vspace{0.5cm}

\textbf{Question 127:} When using \underline{\hspace{1.5cm}} technique, all write operations are made to main memory as well as to the cache, ensuring that main memory is always valid.

\begin{itemize}
    \item[\textbf{A.}] write back
    \item[\textbf{B.}] LRU
    \item[\textbf{C.}] \textcolor{red}{\textbf{write through}}
    \item[\textbf{D.}] unified cache
\end{itemize}

\vspace{0.5cm}

\vspace{0.5cm}

\textbf{Question 177:} Cache memory enhances

\begin{itemize}
    \item[\textbf{A.}] Memory capacity
    \item[\textbf{B.}] \textcolor{red}{\textbf{Memory access time}}
    \item[\textbf{C.}] Secondary storage capacity
    \item[\textbf{D.}] Secondary storage access time
\end{itemize}

\vspace{0.5cm}

\vspace{0.5cm}

\textbf{Question 266:} What is the cache memory level that is integrated into the processor chip and has the lowest latency?

\begin{itemize}
    \item[\textbf{A.}] \textcolor{red}{\textbf{L1 cache}}
    \item[\textbf{B.}] L2 cache
    \item[\textbf{C.}] L3 cache
    \item[\textbf{D.}] L4 cache
\end{itemize}

\vspace{0.5cm}

\vspace{0.5cm}

\textbf{Question 268:} Consider a machine with a byte addressable main memory of $2^{16}$ bytes and block size of 8 bytes. Assume that a direct mapped cache consisting of 32 lines is used with this machine. How many bits are there in the line field of the cache?

\begin{itemize}
    \item[\textbf{A.}] 3
    \item[\textbf{B.}] 4
    \item[\textbf{C.}] \textcolor{red}{\textbf{5}}
    \item[\textbf{D.}] 6
    \item[\textbf{E.}] 7
\end{itemize}

\vspace{0.5cm}

\vspace{0.5cm}

\textbf{Question 317:} What is the primary purpose of cache memory in a computer system?

\begin{itemize}
    \item[\textbf{A.}] \textcolor{red}{\textbf{To store frequently used data for quick access}}
    \item[\textbf{B.}] To store the operating system
    \item[\textbf{C.}] To store user files and documents
    \item[\textbf{D.}] To store the CPU registers
\end{itemize}

\vspace{0.5cm}

\vspace{0.5cm}

\textbf{Question 318:} A set-associative cache consists of 64 lines, divided into four-line sets. $2^{19}$-words main memory contains 4K blocks of 128 words each. How many bits are there in the tag field of the cache?

\begin{itemize}
    \item[\textbf{A.}] 5
    \item[\textbf{B.}] 6
    \item[\textbf{C.}] 7
    \item[\textbf{D.}] \textcolor{red}{\textbf{8}}
    \item[\textbf{E.}] 9
\end{itemize}

\vspace{0.5cm}

\vspace{0.5cm}

\textbf{Question 352:} What is the most common mapping technique used in cache memory in modern computers?

\begin{itemize}
    \item[\textbf{A.}] Direct Mapping
    \item[\textbf{B.}] Fully Associative
    \item[\textbf{C.}] \textcolor{red}{\textbf{Set Associative}}
    \item[\textbf{D.}] None of the mentioned
\end{itemize}

\vspace{0.5cm}

\vspace{0.5cm}

\textbf{Question 360:} What is the special feature of Memory Cache?

\begin{itemize}
    \item[\textbf{A.}] \textcolor{red}{\textbf{Allows faster access than DRAM memory}}
    \item[\textbf{B.}] Memory cache is outboard storage memory
    \item[\textbf{C.}] Allows faster access than CPU registers
    \item[\textbf{D.}] Fixed memory - Read Only Memory
    \item[\textbf{E.}] Has a larger capacity than HDD
\end{itemize}

\vspace{0.5cm}

\vspace{0.5cm}

\textbf{Question 361:} In the direct mapping method from 256MB main memory with 512KB cache, what is the number of bits for the TAG element in the address?

\begin{itemize}
    \item[\textbf{A.}] 6
    \item[\textbf{B.}] 7
    \item[\textbf{C.}] 8
    \item[\textbf{D.}] \textcolor{red}{\textbf{9}}
\end{itemize}

\vspace{0.5cm}

\vspace{0.5cm}

\textbf{Question 398:} Which of the following algorithms is not typically used in cache memory replacement?

\begin{itemize}
    \item[\textbf{A.}] Least Recently Used (LRU)
    \item[\textbf{B.}] First-In-First-Out (FIFO)
    \item[\textbf{C.}] Least Frequently Used (LFU)
    \item[\textbf{D.}] \textcolor{red}{\textbf{Round Robin (RR)}}
\end{itemize}

\vspace{0.5cm}

\vspace{0.5cm}

\textbf{Question 412:} In one-cache system using direct mapping technique, with the cache containing 10 lines, which memory block's data can be transfered to the cache line 7?

\begin{itemize}
    \item[\textbf{A.}] \textcolor{red}{\textbf{7, 17, 27, 37. .....}}
    \item[\textbf{B.}] 0, 1, 2, 3, 4, 5, 6, 7, ...
    \item[\textbf{C.}] 0, 7, 17, 27, 37, ...
    \item[\textbf{D.}] 0, 7, 77, 777, ...
\end{itemize}

\vspace{0.5cm}

\vspace{0.5cm}

\textbf{Question 461:} A two-way set-associative cache has lines of 16 bytes and a total size of 8 Kbytes. The 64-Mbyte main memory is byte addressable. What are the values of Tag, Set and Byte offset fields?

\begin{itemize}
    \item[\textbf{A.}] \textcolor{red}{\textbf{14/8/4}}
    \item[\textbf{B.}] 16/6/4
    \item[\textbf{C.}] 12/10/4
    \item[\textbf{D.}] 15/7/4
    \item[\textbf{E.}] 13/9/4
\end{itemize}

\vspace{0.5cm}

\vspace{0.5cm}

\textbf{Question 484:} Which kind of cache mapping maps each block of main memory into one unique line of the cache?

\begin{itemize}
    \item[\textbf{A.}] \textcolor{red}{\textbf{Direct mapping}}
    \item[\textbf{B.}] Set associative mapping
    \item[\textbf{C.}] Fully associative mapping
    \item[\textbf{D.}] Direct access
\end{itemize}

\vspace{0.5cm}

\vspace{0.5cm}

\textbf{Question 485:} If the address received when accessing the cache memory includes Tag, Line, Word. At that time, what mapping technique is the system using?

\begin{itemize}
    \item[\textbf{A.}] \textcolor{red}{\textbf{Direct mapping}}
    \item[\textbf{B.}] Set-direct mapping
    \item[\textbf{C.}] Associative mapping
    \item[\textbf{D.}] Set-associative mapping
\end{itemize}

\vspace{0.5cm}

\vspace{0.5cm}

\textbf{Question 508:} In modern processors, why is cache memory typically organized into multiple levels, such as L1, L2, and L3 caches?

\begin{itemize}
    \item[\textbf{A.}] \textcolor{red}{\textbf{To balance the trade-off between speed and size by having different cache levels}}
    \item[\textbf{B.}] To conserve chip area by reducing the cache size and increasing the memory size
    \item[\textbf{C.}] To provide redundancy in case of cache failures
    \item[\textbf{D.}] To improve cache coherence in multi-processor systems
\end{itemize}

\vspace{0.5cm}

\vspace{0.5cm}

\textbf{Question 517:} How many techniques are there for dealing with an update to a cache line? What are their names?

\begin{itemize}
    \item[\textbf{A.}] \textcolor{red}{\textbf{There are 2 techniques and their names are written through, write back}}
    \item[\textbf{B.}] There are 3 techniques and their names are written through, write back, and write allocate
    \item[\textbf{C.}] There is only one technique and its name is "write not allocate"
    \item[\textbf{D.}] There are 4 techniques and their names are written through, write back, write allocate, write not allocate
\end{itemize}

\vspace{0.5cm}

\vspace{0.5cm}

\textbf{Question 530:} A two-way set-associative cache has lines of 16 bytes and a total size of 8 Kbytes. The 64-Mbyte main memory is byte addressable. How many bits are there in the set field of the cache?

\begin{itemize}
    \item[\textbf{A.}] 5
    \item[\textbf{B.}] 6
    \item[\textbf{C.}] 7
    \item[\textbf{D.}] \textcolor{red}{\textbf{8}}
    \item[\textbf{E.}] 9
\end{itemize}

\vspace{0.5cm}

\vspace{0.5cm}

\textbf{Question 396:} What is the benefit of cache memory in terms of computer performance?

\begin{itemize}
    \item[\textbf{A.}] It reduces the computer's power consumption.
    \item[\textbf{B.}] \textcolor{red}{\textbf{It speeds up data access and improves system performance.}}
    \item[\textbf{C.}] It increases the storage capacity of the computer.
    \item[\textbf{D.}] It serves as a backup storage for critical files.
\end{itemize}

\vspace{0.5cm}

\vspace{0.5cm}

\textbf{Question 518:} A EPROM has 20 address lines and 8 bit data lines, what is maximum memory capacity of this chip?

\begin{itemize}
    \item[\textbf{A.}] 1 Mbits
    \item[\textbf{B.}] 2 Mbits
    \item[\textbf{C.}] 4 Mbits
    \item[\textbf{D.}] \textcolor{red}{\textbf{8 Mbits}}
    \item[\textbf{E.}] 16 Mbits
\end{itemize}

\vspace{0.5cm}

\vspace{0.5cm}

\textbf{Question 79:} Internal memory capacity is typically expressed in terms of \underline{\hspace{1.5cm}}

\begin{itemize}
    \item[\textbf{A.}] hertz
    \item[\textbf{B.}] nanos
    \item[\textbf{C.}] \textcolor{red}{\textbf{bytes}}
    \item[\textbf{D.}] LOR
\end{itemize}

\vspace{0.5cm}

\vspace{0.5cm}

\section{Chapter 06: Internal Memory}
\textit{Total questions in chapter: 37}

\textbf{Question 91:} RAM must be provided with a constant power supply.

\begin{itemize}
    \item[\textbf{A.}] \textcolor{red}{\textbf{True}}
    \item[\textbf{B.}] False
\end{itemize}

\vspace{0.5cm}

\vspace{0.5cm}

\textbf{Question 92:} A characteristic of ROM is that it is volatile.

\begin{itemize}
    \item[\textbf{A.}] True
    \item[\textbf{B.}] \textcolor{red}{\textbf{False}}
\end{itemize}

\vspace{0.5cm}

\vspace{0.5cm}

\textbf{Question 164:} A static RAM will hold its data as long as power is supplied to it.

\begin{itemize}
    \item[\textbf{A.}] \textcolor{red}{\textbf{True}}
    \item[\textbf{B.}] False
\end{itemize}

\vspace{0.5cm}

\vspace{0.5cm}

\textbf{Question 192:} The two traditional forms of RAM used in computers are DRAM and SRAM.

\begin{itemize}
    \item[\textbf{A.}] \textcolor{red}{\textbf{True}}
    \item[\textbf{B.}] False
\end{itemize}

\vspace{0.5cm}

\vspace{0.5cm}

\textbf{Question 270:} What is the main idea of using Hamming code for error correction?

\begin{itemize}
    \item[\textbf{A.}] Adding extra parity bits to the data bits such that the number of 1s in each subset of bits is even
    \item[\textbf{B.}] Adding extra parity bits to the data bits such that the number of 1s in each subset of bits is odd
    \item[\textbf{C.}] \textcolor{red}{\textbf{Adding extra parity bits to the data bits such that the parity bits form a binary number indicating the position of the error bit}}
    \item[\textbf{D.}] Adding extra parity bits to the data bits such that the parity bits form a binary number indicating the number of error bits
\end{itemize}

\vspace{0.5cm}

\vspace{0.5cm}

\textbf{Question 321:} Using Hamming Code with one error corection to store an 12-bit word in memory, the stored word 111001001101 consists of 8 bits data and 4 bit bits parity check. What are the parity bits?

\begin{itemize}
    \item[\textbf{A.}] \textcolor{red}{\textbf{0110}}
    \item[\textbf{B.}] 0111
    \item[\textbf{C.}] 1110
    \item[\textbf{D.}] 0101
    \item[\textbf{E.}] None of the mentioned
\end{itemize}

\vspace{0.5cm}

\vspace{0.5cm}

\textbf{Question 383:} Which one of the following is an invalid about RAM?

\begin{itemize}
    \item[\textbf{A.}] Both static and dynamic RAMs are volatile; that is, power must be continuously supplied to the memory to preserve the bit values.
    \item[\textbf{B.}] A dynamic memory cell is simpler and smaller than a static memory cell.
    \item[\textbf{C.}] \textcolor{red}{\textbf{Both static and dynamic RAMs requires the supporting refresh circuitry}}
    \item[\textbf{D.}] SRAMs are somewhat faster than DRAMs
\end{itemize}

\vspace{0.5cm}

\vspace{0.5cm}

\textbf{Question 432:} How many check bits are needed if the Hamming error correction code is used to detect single bit errors in a 2030-bit data?

\begin{itemize}
    \item[\textbf{A.}] \textcolor{red}{\textbf{11}}
    \item[\textbf{B.}] 10
    \item[\textbf{C.}] 12
    \item[\textbf{D.}] 9
\end{itemize}

\vspace{0.5cm}

\vspace{0.5cm}

\textbf{Question 434:} Given the 11-bit data word 001 0010 1010, generate the corresponding 15-bit Hamming code word.

\begin{itemize}
    \item[\textbf{A.}] 001 0010 01011011
    \item[\textbf{B.}] 001 0010 01011001
    \item[\textbf{C.}] 001 0010 01010101
    \item[\textbf{D.}] \textcolor{red}{\textbf{001 0010 01101001}}
\end{itemize}

\vspace{0.5cm}

\vspace{0.5cm}

\textbf{Question 444:} How many parity check bits must be included with the data word to achieve single error correction when the data word contains 1023 bits?

\begin{itemize}
    \item[\textbf{A.}] 10
    \item[\textbf{B.}] \textcolor{red}{\textbf{11}}
    \item[\textbf{C.}] 12
    \item[\textbf{D.}] 9
\end{itemize}

\vspace{0.5cm}

\vspace{0.5cm}

\textbf{Question 462:} How many check bits are needed for a Hamming error correction code to detect single-error correction in a 1 Kbyte data word?

\begin{itemize}
    \item[\textbf{A.}] 11 bits
    \item[\textbf{B.}] 12 bits
    \item[\textbf{C.}] 13 bits
    \item[\textbf{D.}] \textcolor{red}{\textbf{14 bits}}
    \item[\textbf{E.}] 15 bits
\end{itemize}

\vspace{0.5cm}

\vspace{0.5cm}

\textbf{Question 463:} The use of Hamming code on 4-bit words (M = 4) is illustrated in the below Venn diagrams. With three intersecting circles, there are seven compartments. We assign the 4 data bits to the inner compartments. The remaining compartments are filled with what are called parity bits. Each parity bit is chosen so that the total number of 1s in its circle is even. Choose the correct value of the corresponding parity bit in circles A, B, C.

\begin{itemize}
    \item[\textbf{A.}] \textcolor{red}{\textbf{1; 0; 0}}
    \item[\textbf{B.}] 0; 1; 0
    \item[\textbf{C.}] 0; 0; 1
    \item[\textbf{D.}] 1; 0; 1
\end{itemize}

\vspace{0.5cm}

\vspace{0.5cm}

\textbf{Question 464:} Using Hamming Code with one error correction to store a 12-bit word in memory, the stored word 101001011101 consists of 8 bits data and 4 bits parity check. What are the data bits?

\begin{itemize}
    \item[\textbf{A.}] 01001101
    \item[\textbf{B.}] 10100100
    \item[\textbf{C.}] 10101011
    \item[\textbf{D.}] 01011101
    \item[\textbf{E.}] \textcolor{red}{\textbf{None of the mentioned}}
\end{itemize}

\vspace{0.5cm}

\vspace{0.5cm}

\textbf{Question 512:} EEPROM is read-mostly-memory and volatile

\begin{itemize}
    \item[\textbf{A.}] True
    \item[\textbf{B.}] \textcolor{red}{\textbf{False}}
\end{itemize}

\vspace{0.5cm}

\vspace{0.5cm}

\textbf{Question 122:} One distinguishing characteristic of memory that is designated as \underline{\hspace{1.5cm}} is that it is possible to read data from the memory and to write new data into the memory easily and rapidly.

\begin{itemize}
    \item[\textbf{A.}] \textcolor{red}{\textbf{RAM}}
    \item[\textbf{B.}] ROM
    \item[\textbf{C.}] EPROM
    \item[\textbf{D.}] EEPROM
\end{itemize}

\vspace{0.5cm}

\vspace{0.5cm}

\textbf{Question 175:} In a volatile memory, information decays naturally or is lost when electrical power is switched

\begin{itemize}
    \item[\textbf{A.}] \textcolor{red}{\textbf{True}}
    \item[\textbf{B.}] False
\end{itemize}

\vspace{0.5cm}

\vspace{0.5cm}

\textbf{Question 206:} A \underline{\hspace{1.5cm}} contains a permanent pattern of data that cannot be changed, is nonvolatile, and cannot have new data written into it.

\begin{itemize}
    \item[\textbf{A.}] flash memory
    \item[\textbf{B.}] \textcolor{red}{\textbf{ROM}}
    \item[\textbf{C.}] RAM
    \item[\textbf{D.}] SRAM
\end{itemize}

\vspace{0.5cm}

\vspace{0.5cm}

\textbf{Question 213:} One distinguishing characteristic of memory that is designated as \underline{\hspace{1.5cm}} is that it is possible both to read data from the memory and to write new data into the memory easily and rapidly.

\begin{itemize}
    \item[\textbf{A.}] \textcolor{red}{\textbf{RAM}}
    \item[\textbf{B.}] EPROM
    \item[\textbf{C.}] ROM
    \item[\textbf{D.}] EEPROM
\end{itemize}

\vspace{0.5cm}

\vspace{0.5cm}

\textbf{Question 242:} A \underline{\hspace{1.5cm}} contains a permanent pattern of data that cannot be changed, is nonvolatile, and cannot have new data written into it.

\begin{itemize}
    \item[\textbf{A.}] RAM
    \item[\textbf{B.}] SRAM
    \item[\textbf{C.}] \textcolor{red}{\textbf{ROM}}
    \item[\textbf{D.}] flash memory
\end{itemize}

\vspace{0.5cm}

\vspace{0.5cm}

\textbf{Question 271:} What are the key differences in the architecture of NOR and NAND flash memory?

\begin{itemize}
    \item[\textbf{A.}] NOR flash memory cells are connected in series, while NAND flash memory cells are connected in parallel
    \item[\textbf{B.}] \textcolor{red}{\textbf{NOR flash memory cells are connected in parallel, while NAND flash memory cells are connected in series}}
    \item[\textbf{C.}] Both NOR and NAND flash memory cells are connected in series
    \item[\textbf{D.}] Both NOR and NAND flash memory cells are connected in parallel
    \item[\textbf{E.}] All of the mentioned are wrong
\end{itemize}

\vspace{0.5cm}

\vspace{0.5cm}

\textbf{Question 11:} Which properties do all semiconductor memory cells share?

\begin{itemize}
    \item[\textbf{A.}] they exhibit two stable states which can be used to represent binary 1 and 0
    \item[\textbf{B.}] they are capable of being written into to set the state
    \item[\textbf{C.}] they are capable of being read to sense the state
    \item[\textbf{D.}] \textcolor{red}{\textbf{all of the above}}
\end{itemize}

\vspace{0.5cm}

\vspace{0.5cm}

\textbf{Question 16:} A number of chips can be grouped together to form a memory bank.

\begin{itemize}
    \item[\textbf{A.}] \textcolor{red}{\textbf{True}}
    \item[\textbf{B.}] False
\end{itemize}

\vspace{0.5cm}

\vspace{0.5cm}

\textbf{Question 60:} All DRAMs require a refresh operation.

\begin{itemize}
    \item[\textbf{A.}] \textcolor{red}{\textbf{True}}
    \item[\textbf{B.}] False
\end{itemize}

\vspace{0.5cm}

\vspace{0.5cm}

\textbf{Question 65:} Which of the following memory types are nonvolatile?

\begin{itemize}
    \item[\textbf{A.}] erasable PROM
    \item[\textbf{B.}] programmable ROM
    \item[\textbf{C.}] flash memory
    \item[\textbf{D.}] \textcolor{red}{\textbf{all of the above}}
\end{itemize}

\vspace{0.5cm}

\vspace{0.5cm}

\textbf{Question 81:} The basic element of a semiconductor memory is the memory cell.

\begin{itemize}
    \item[\textbf{A.}] \textcolor{red}{\textbf{True}}
    \item[\textbf{B.}] False
\end{itemize}

\vspace{0.5cm}

\vspace{0.5cm}

\textbf{Question 158:} Semiconductor memory comes in packaged chips.

\begin{itemize}
    \item[\textbf{A.}] \textcolor{red}{\textbf{True}}
    \item[\textbf{B.}] False
\end{itemize}

\vspace{0.5cm}

\vspace{0.5cm}

\textbf{Question 255:} With \underline{\hspace{1.5cm}} the microchip is organized so that a section of memory cells are erased in a single action.

\begin{itemize}
    \item[\textbf{A.}] \textcolor{red}{\textbf{flash memory}}
    \item[\textbf{B.}] SDRAM
    \item[\textbf{C.}] DRAM
    \item[\textbf{D.}] EEPROM
\end{itemize}

\vspace{0.5cm}

\vspace{0.5cm}

\textbf{Question 436:} Semiconductor memory comes in packaged chips.

\begin{itemize}
    \item[\textbf{A.}] False
    \item[\textbf{B.}] \textcolor{red}{\textbf{True}}
\end{itemize}

\vspace{0.5cm}

\vspace{0.5cm}

\textbf{Question 486:} What is the key advantage of double-data-rate SDRAM compared to standard SDRAM?

\begin{itemize}
    \item[\textbf{A.}] It has a larger memory capacity.
    \item[\textbf{B.}] \textcolor{red}{\textbf{It can send data twice per clock cycle.}}
    \item[\textbf{C.}] It uses less power than SDRAM.
    \item[\textbf{D.}] It is only compatible with desktop computers.
\end{itemize}

\vspace{0.5cm}

\vspace{0.5cm}

\textbf{Question 516:} \underline{\hspace{1.5cm}} is made with cells that store data as charge on capacitors.

\begin{itemize}
    \item[\textbf{A.}] SRAM
    \item[\textbf{B.}] \textcolor{red}{\textbf{DRAM}}
    \item[\textbf{C.}] RDRAM
    \item[\textbf{D.}] CDRAM
\end{itemize}

\vspace{0.5cm}

\vspace{0.5cm}

\textbf{Question 509:} In terms of organization, what is the basic element of semiconductor memory?

\begin{itemize}
    \item[\textbf{A.}] Memory capacity
    \item[\textbf{B.}] Memory block
    \item[\textbf{C.}] \textcolor{red}{\textbf{Memory cell}}
    \item[\textbf{D.}] Memory location
\end{itemize}

\vspace{0.5cm}

\vspace{0.5cm}

\textbf{Question 365:} In terms of performance, what is the main advantage of a solid state drive over a magnetic disk?

\begin{itemize}
    \item[\textbf{A.}] \textcolor{red}{\textbf{A solid state drive has faster access time, lower latency, and higher reliability}}
    \item[\textbf{B.}] A solid state drive has larger capacity, lower power consumption, and lower cost
    \item[\textbf{C.}] A solid state drive has better compatibility, longer lifespan, and higher security
    \item[\textbf{D.}] A solid state drive has none of the mentioned advantages over a magnetic disk
\end{itemize}

\vspace{0.5cm}

\vspace{0.5cm}

\textbf{Question 253:} In a \underline{\hspace{1.5cm}}, binary values are stored using traditional flip-flop logic-gate configurations.

\begin{itemize}
    \item[\textbf{A.}] ROM
    \item[\textbf{B.}] \textcolor{red}{\textbf{SRAM}}
    \item[\textbf{C.}] DRAM
    \item[\textbf{D.}] RAM
\end{itemize}

\vspace{0.5cm}

\vspace{0.5cm}

\textbf{Question 130:} An error-correcting code enhances the reliability of the memory at the cost of added complexity.

\begin{itemize}
    \item[\textbf{A.}] \textcolor{red}{\textbf{True}}
    \item[\textbf{B.}] False
\end{itemize}

\vspace{0.5cm}

\vspace{0.5cm}

\textbf{Question 363:} In error correcting code (single ECC), how many bits are used to correct one bit in 8-bit data?

\begin{itemize}
    \item[\textbf{A.}] \textcolor{red}{\textbf{4.}}
    \item[\textbf{B.}] 5
    \item[\textbf{C.}] 6
    \item[\textbf{D.}] 7
\end{itemize}

\vspace{0.5cm}

\vspace{0.5cm}

\textbf{Question 422:} The data word that was read from memory 001001001011110.Show data stored in memory

\begin{itemize}
    \item[\textbf{A.}] 101 0010 1001
    \item[\textbf{B.}] 001 0010 1110
    \item[\textbf{C.}] 101 0010 1011
    \item[\textbf{D.}] \textcolor{red}{\textbf{001 0010 1011}}
\end{itemize}

\vspace{0.5cm}

\vspace{0.5cm}

\textbf{Question 453:} For the data 10110101100, calculated check bits?

\begin{itemize}
    \item[\textbf{A.}] 0101
    \item[\textbf{B.}] 0011
    \item[\textbf{C.}] \textcolor{red}{\textbf{1100}}
    \item[\textbf{D.}] 1010
\end{itemize}

\vspace{0.5cm}

\vspace{0.5cm}

\section{Chapter 07: External Memory}
\textit{Total questions in chapter: 40}

\textbf{Question 103:} The SSDs now on the market use a type of semiconductor memory referred to as flash memory

\begin{itemize}
    \item[\textbf{A.}] \textcolor{red}{\textbf{True}}
    \item[\textbf{B.}] False
\end{itemize}

\vspace{0.5cm}

\vspace{0.5cm}

\textbf{Question 364:} Which RAID level uses striping technique with a minimum of 3 disks and provides fault tolerance through the use of parity bit?

\begin{itemize}
    \item[\textbf{A.}] RAID 0
    \item[\textbf{B.}] RAID 1
    \item[\textbf{C.}] RAID 2
    \item[\textbf{D.}] \textcolor{red}{\textbf{RAID 3}}
\end{itemize}

\vspace{0.5cm}

\vspace{0.5cm}

\textbf{Question 21:} RAID level 0 is not a true member of the RAID family because it does not include redundancy to improve performance.

\begin{itemize}
    \item[\textbf{A.}] \textcolor{red}{\textbf{True}}
    \item[\textbf{B.}] False
\end{itemize}

\vspace{0.5cm}

\vspace{0.5cm}

\textbf{Question 30:} Consider a magnetic disk drive with 4 surfaces, 512 tracks per surface, and 32 sectors per track. If the sector size is 512 byte. What is the disk capacity?

\begin{itemize}
    \item[\textbf{A.}] \textcolor{red}{\textbf{32 MB}}
    \item[\textbf{B.}] 512 KB
    \item[\textbf{C.}] 64 KB
    \item[\textbf{D.}] 16 GB
    \item[\textbf{E.}] 64 MB
\end{itemize}

\vspace{0.5cm}

\vspace{0.5cm}

\textbf{Question 36:} Because data are striped in very small strips, RAID 3 cannot achieve very high data transfer

\begin{itemize}
    \item[\textbf{A.}] \textcolor{red}{\textbf{False}}
    \item[\textbf{B.}] True
\end{itemize}

\vspace{0.5cm}

\vspace{0.5cm}

\textbf{Question 99:} What is the purpose of RAID system?

\begin{itemize}
    \item[\textbf{A.}] It increases the processor speed
    \item[\textbf{B.}] It increases the disk storage capacity
    \item[\textbf{C.}] \textcolor{red}{\textbf{It increases the disk storage capacity and availability}}
    \item[\textbf{D.}] It increases operating system efficiency
    \item[\textbf{E.}] It decreases operating system efficiency
\end{itemize}

\vspace{0.5cm}

\vspace{0.5cm}

\textbf{Question 136:} In most contemporary systems fixed-length sectors are used, with \underline{\hspace{1.5cm}} bytes being the nearly universal sector size.

\begin{itemize}
    \item[\textbf{A.}] 64
    \item[\textbf{B.}] 128
    \item[\textbf{C.}] 256
    \item[\textbf{D.}] \textcolor{red}{\textbf{512}}
\end{itemize}

\vspace{0.5cm}

\vspace{0.5cm}

\textbf{Question 145:} The sum of the seek time and the rotational delay equals the \underline{\hspace{1.5cm}}, which is the time it takes to get into position to read or write.

\begin{itemize}
    \item[\textbf{A.}] \textcolor{red}{\textbf{access time}}
    \item[\textbf{B.}] gap time
    \item[\textbf{C.}] transfer time
    \item[\textbf{D.}] constant angular velocity
\end{itemize}

\vspace{0.5cm}

\vspace{0.5cm}

\textbf{Question 167:} When the magnetizable coating is applied to both sides of the platter the disk is then referred

\begin{itemize}
    \item[\textbf{A.}] multiple sided
    \item[\textbf{B.}] substrate
    \item[\textbf{C.}] \textcolor{red}{\textbf{double sided}}
    \item[\textbf{D.}] all of the above
\end{itemize}

\vspace{0.5cm}

\vspace{0.5cm}

\textbf{Question 187:} The set of all the tracks in the same relative position on the platter is referred to as a \underline{\hspace{1.5cm}}

\begin{itemize}
    \item[\textbf{A.}] \textcolor{red}{\textbf{cylinder}}
    \item[\textbf{B.}] floppy disk
    \item[\textbf{C.}] single-sided disk
    \item[\textbf{D.}] sector
\end{itemize}

\vspace{0.5cm}

\vspace{0.5cm}

\textbf{Question 193:} The sum of the seek time and the rotational delay equals the \underline{\hspace{1.5cm}}, which is the time it takes to get into position to read or write.

\begin{itemize}
    \item[\textbf{A.}] constant angular velocity
    \item[\textbf{B.}] transfer time
    \item[\textbf{C.}] \textcolor{red}{\textbf{access time}}
    \item[\textbf{D.}] gap time
\end{itemize}

\vspace{0.5cm}

\vspace{0.5cm}

\textbf{Question 210:} What is the purpose of RAID system?

\begin{itemize}
    \item[\textbf{A.}] \textcolor{red}{\textbf{It increases the disk storage capacity and availability}}
    \item[\textbf{B.}] It increases the disk storage capacity
    \item[\textbf{C.}] It increases operating system efficiency
    \item[\textbf{D.}] It decreases operating system efficiency
    \item[\textbf{E.}] It increases the processor speed
\end{itemize}

\vspace{0.5cm}

\vspace{0.5cm}

\textbf{Question 320:} How many bytes of data does each sector in the Winchester hard drive disk have?

\begin{itemize}
    \item[\textbf{A.}] 128 bytes
    \item[\textbf{B.}] 256 bytes
    \item[\textbf{C.}] \textcolor{red}{\textbf{512 bytes}}
    \item[\textbf{D.}] 1024 bytes
    \item[\textbf{E.}] 4096 bytes
\end{itemize}

\vspace{0.5cm}

\vspace{0.5cm}

\textbf{Question 323:} With the hard disk data layout, the set of all the tracks in the same relative position on the platter, is called \underline{\hspace{1.5cm}}.

\begin{itemize}
    \item[\textbf{A.}] \textcolor{red}{\textbf{Cylinder}}
    \item[\textbf{B.}] Tracks
    \item[\textbf{C.}] Inter-track gap
    \item[\textbf{D.}] Sector
\end{itemize}

\vspace{0.5cm}

\vspace{0.5cm}

\textbf{Question 324:} Consider a 5-drive, 200 GBytes-per-drive RAID array. What is the available data storage capacity for each of the RAID levels 5?

\begin{itemize}
    \item[\textbf{A.}] 200 GBytes
    \item[\textbf{B.}] 400 GBytes
    \item[\textbf{C.}] 600 GBytes
    \item[\textbf{D.}] \textcolor{red}{\textbf{800 GBytes}}
    \item[\textbf{E.}] None of the mentioned
\end{itemize}

\vspace{0.5cm}

\vspace{0.5cm}

\textbf{Question 384:} The speed of data delivery is your main concern when configuring a RAID drive for a Media Streaming Server. This server has two hard disks installed. What type of RAID should you install, and what type of data will be stored on Disk 1 and Disk 2?

\begin{itemize}
    \item[\textbf{A.}] \textcolor{red}{\textbf{RAID 0 - Disk 1 (Stripe) and Disk 2 (Stripe)}}
    \item[\textbf{B.}] RAID 0 - Disk 1 (Mirror) and Disk 2 (Mirror)
    \item[\textbf{C.}] RAID 1 - Disk 1 (Stripe) and Disk 2 (Stripe)
    \item[\textbf{D.}] RAID 1 - Disk 1 (Mirror) and Disk 2 (Mirror)
\end{itemize}

\vspace{0.5cm}

\vspace{0.5cm}

\textbf{Question 421:} A 6 drive, 320GB-per-drive RAID array. What is the available data storage capacity for each of the RAID levels 5?

\begin{itemize}
    \item[\textbf{A.}] 960 GB
    \item[\textbf{B.}] 1920 GB
    \item[\textbf{C.}] \textcolor{red}{\textbf{1600 GB}}
    \item[\textbf{D.}] 1280 GB
\end{itemize}

\vspace{0.5cm}

\vspace{0.5cm}

\textbf{Question 423:} What is the purpose of RAID system?

\begin{itemize}
    \item[\textbf{A.}] It increases the processor speed
    \item[\textbf{B.}] It increases operating system efficiency
    \item[\textbf{C.}] It decreases operating system efficiency
    \item[\textbf{D.}] \textcolor{red}{\textbf{It increases the disk storage capacity and availability}}
    \item[\textbf{E.}] It increases the disk storage capacity
\end{itemize}

\vspace{0.5cm}

\vspace{0.5cm}

\textbf{Question 451:} Consider the RAID array, 250GB each. Need 1000 GB storage data memory. How many drives for each RAID level 3?

\begin{itemize}
    \item[\textbf{A.}] \textcolor{red}{\textbf{5}}
    \item[\textbf{B.}] 3
    \item[\textbf{C.}] 6
    \item[\textbf{D.}] 4
\end{itemize}

\vspace{0.5cm}

\vspace{0.5cm}

\textbf{Question 465:} In terms of data storage, what is the main distinction between a solid state drive and a magnetic disk?

\begin{itemize}
    \item[\textbf{A.}] \textcolor{red}{\textbf{A magnetic disk uses a spinning platter and a read/write head, while a solid state drive uses flash memory chips}}
    \item[\textbf{B.}] A magnetic disk uses flash memory chips, while a solid state drive uses a spinning platter and a read/write head
    \item[\textbf{C.}] A magnetic disk and a solid state drive both use flash memory chips, but with different interfaces
    \item[\textbf{D.}] A magnetic disk and a solid state drive both use a spinning platter and a read/write head, but with different speeds
\end{itemize}

\vspace{0.5cm}

\vspace{0.5cm}

\textbf{Question 510:} Which is incorrect about RAID?

\begin{itemize}
    \item[\textbf{A.}] RAID is a set of physical disk drives viewed by the operating system as a single logical drive.
    \item[\textbf{B.}] Data are distributed across the physical drives of an array in a scheme known as striping, described subsequently.
    \item[\textbf{C.}] \textcolor{red}{\textbf{RAID is only used for Windows OS}}
    \item[\textbf{D.}] Redundant disk capacity is used to store parity information, which guarantees data recoverability in case of a disk failure.
\end{itemize}

\vspace{0.5cm}

\vspace{0.5cm}

\textbf{Question 273:} \underline{\hspace{1.5cm}} is a set of physical disk drives viewed by the operating system as a single logical drive.

\begin{itemize}
    \item[\textbf{A.}] CLV
    \item[\textbf{B.}] SSD
    \item[\textbf{C.}] \textcolor{red}{\textbf{RAID}}
    \item[\textbf{D.}] CAV
\end{itemize}

\vspace{0.5cm}

\vspace{0.5cm}

\textbf{Question 17:} Data are transferred to and from the disk in \underline{\hspace{1.5cm}}.

\begin{itemize}
    \item[\textbf{A.}] tracks
    \item[\textbf{B.}] gaps
    \item[\textbf{C.}] \textcolor{red}{\textbf{sectors}}
    \item[\textbf{D.}] pits
\end{itemize}

\vspace{0.5cm}

\vspace{0.5cm}

\textbf{Question 173:} There are typically hundreds of sectors per track and they may be either fixed or variable length

\begin{itemize}
    \item[\textbf{A.}] \textcolor{red}{\textbf{True}}
    \item[\textbf{B.}] False
\end{itemize}

\vspace{0.5cm}

\vspace{0.5cm}

\textbf{Question 160:} Magnetic disks are the foundation of external memory on virtually all computer systems.

\begin{itemize}
    \item[\textbf{A.}] \textcolor{red}{\textbf{True}}
    \item[\textbf{B.}] False
\end{itemize}

\vspace{0.5cm}

\vspace{0.5cm}

\textbf{Question 184:} Which of the following is an example of sequential access media?

\begin{itemize}
    \item[\textbf{A.}] Main Memory
    \item[\textbf{B.}] CD
    \item[\textbf{C.}] \textcolor{red}{\textbf{Magnetic tape}}
    \item[\textbf{D.}] Cache memory
    \item[\textbf{E.}] Magnetic disk
\end{itemize}

\vspace{0.5cm}

\vspace{0.5cm}

\textbf{Question 275:} Which is the correct choice for sorting in increasing speed average of memory?

\begin{itemize}
    \item[\textbf{A.}] SSD $\square$ Main Memory $\square$ Cache Memory $\square$ Magnetic Tape
    \item[\textbf{B.}] Magnetic Tape $\square$ SSD $\square$ Cache Memory $\square$ Main Memory
    \item[\textbf{C.}] Magnetic Disk $\square$ SSD $\square$ Cache Memory $\square$ Main Memory
    \item[\textbf{D.}] Magnetic Disk $\square$ Magnetic Tape $\square$ Main Memory $\square$ Cache Memory
    \item[\textbf{E.}] \textcolor{red}{\textbf{Magnetic Disk $\square$ SSD $\square$ Main Memory $\square$ Cache Memory}}
\end{itemize}

\vspace{0.5cm}

\vspace{0.5cm}

\textbf{Question 447:} How many parity check bits must be included with the data word to achieve single error correction when the data word contains 2048 bits?

\begin{itemize}
    \item[\textbf{A.}] 9
    \item[\textbf{B.}] 10
    \item[\textbf{C.}] \textcolor{red}{\textbf{12}}
    \item[\textbf{D.}] 11
\end{itemize}

\vspace{0.5cm}

\vspace{0.5cm}

\textbf{Question 315:} Which of the following memory devices has the lowest access speed?

\begin{itemize}
    \item[\textbf{A.}] ROM
    \item[\textbf{B.}] Flash memory
    \item[\textbf{C.}] \textcolor{red}{\textbf{Magnetic tape}}
    \item[\textbf{D.}] HDD
    \item[\textbf{E.}] Cache
\end{itemize}

\vspace{0.5cm}

\vspace{0.5cm}

\textbf{Question 487:} Which of the following issues are peculiar to SSDs compared to HDDs? (choose two correct answers)

\begin{itemize}
    \item[\textbf{A.}] \textcolor{red}{\textbf{Performance slows down as the device is used.}}
    \item[\textbf{B.}] \textcolor{red}{\textbf{A certain number of writes make the flash memory unusable.}}
    \item[\textbf{C.}] They are more vulnerable to physical damage.
    \item[\textbf{D.}] They require defragmentation to maintain performance.
    \item[\textbf{E.}] They are slower at data retrieval compared to HDDs.
\end{itemize}

\vspace{0.5cm}

\vspace{0.5cm}

\textbf{Question 51:} \underline{\hspace{1.5cm}} is when the disk rotates more slowly for accesses near the outer edge than for those near the center.

\begin{itemize}
    \item[\textbf{A.}] Seek time
    \item[\textbf{B.}] Magnetoresistive
    \item[\textbf{C.}] Constant angular velocity (CAV)
    \item[\textbf{D.}] \textcolor{red}{\textbf{Constant linear velocity (CLV)}}
\end{itemize}

\vspace{0.5cm}

\vspace{0.5cm}

\textbf{Question 64:} A \underline{\hspace{1.5cm}} disk is permanently mounted in the disk drive, such as the hard disk in a pers

\begin{itemize}
    \item[\textbf{A.}] removable
    \item[\textbf{B.}] double sided
    \item[\textbf{C.}] \textcolor{red}{\textbf{nonremovable}}
    \item[\textbf{D.}] movable-head
\end{itemize}

\vspace{0.5cm}

\vspace{0.5cm}

\textbf{Question 87:} \underline{\hspace{1.5cm}} is the standardized scheme for multiple-disk database design.

\begin{itemize}
    \item[\textbf{A.}] \textcolor{red}{\textbf{RAID}}
    \item[\textbf{B.}] CAV
    \item[\textbf{C.}] CLV
    \item[\textbf{D.}] SSD
\end{itemize}

\vspace{0.5cm}

\vspace{0.5cm}

\textbf{Question 117:} Scanning information at the same rate by rotating the disk at a fixed speed is known as the \_

\begin{itemize}
    \item[\textbf{A.}] \textcolor{red}{\textbf{constant angular velocity}}
    \item[\textbf{B.}] magnetoresistive
    \item[\textbf{C.}] rotational delay
    \item[\textbf{D.}] constant linear velocity
\end{itemize}

\vspace{0.5cm}

\vspace{0.5cm}

\textbf{Question 140:} A \underline{\hspace{1.5cm}} is a high-definition video disk that can store 25 Gbytes on a single layer on a single side.

\begin{itemize}
    \item[\textbf{A.}] DVD
    \item[\textbf{B.}] DVD-R
    \item[\textbf{C.}] DVD-RW
    \item[\textbf{D.}] \textcolor{red}{\textbf{Blu-ray DVD}}
\end{itemize}

\vspace{0.5cm}

\vspace{0.5cm}

\textbf{Question 176:} External memory is often equated with main memory.

\begin{itemize}
    \item[\textbf{A.}] True
    \item[\textbf{B.}] \textcolor{red}{\textbf{False}}
\end{itemize}

\vspace{0.5cm}

\vspace{0.5cm}

\textbf{Question 191:} The disadvantage of \underline{\hspace{1.5cm}} is that the amount of data that can be stored on the long outer tracks is the same as what can be stored on the short inner tracks.

\begin{itemize}
    \item[\textbf{A.}] \textcolor{red}{\textbf{CAV}}
    \item[\textbf{B.}] SSD
    \item[\textbf{C.}] ROM
    \item[\textbf{D.}] CLV
\end{itemize}

\vspace{0.5cm}

\vspace{0.5cm}

\textbf{Question 207:} \underline{\hspace{1.5cm}} is when the disk rotates more slowly for accesses near the outer edge than for those near the center.

\begin{itemize}
    \item[\textbf{A.}] \textcolor{red}{\textbf{Constant linear velocity (CLV)}}
    \item[\textbf{B.}] Constant angular velocity (CAV)
    \item[\textbf{C.}] Seek time
    \item[\textbf{D.}] Magnetoresistive
\end{itemize}

\vspace{0.5cm}

\vspace{0.5cm}

\textbf{Question 251:} Adjacent tracks are separated by \underline{\hspace{1.5cm}}.

\begin{itemize}
    \item[\textbf{A.}] sectors
    \item[\textbf{B.}] \textcolor{red}{\textbf{gaps}}
    \item[\textbf{C.}] pits
    \item[\textbf{D.}] heads
\end{itemize}

\vspace{0.5cm}

\vspace{0.5cm}

\textbf{Question 325:} What is incorrect about SSDs have the following advantages over HDDs?

\begin{itemize}
    \item[\textbf{A.}] \textcolor{red}{\textbf{Higher access times and latency rates: Over 10 times slower than the HDD.}}
    \item[\textbf{B.}] Durability: Less susceptible to physical shock and vibration.
    \item[\textbf{C.}] Longer lifespan: SSDs are not susceptible to mechanical wear.
    \item[\textbf{D.}] Lower power consumption: SSDs use considerably less power than comparable-size HDDs.
    \item[\textbf{E.}] Quieter and cooler running capabilities: Less space required, lower energy costs, and a greener enterprise.
\end{itemize}

\vspace{0.5cm}

\vspace{0.5cm}

\section{Chapter 08: Input \& Output}
\textit{Total questions in chapter: 42}

\textbf{Question 3:} An I/O device is referred to as a \underline{\hspace{1.5cm}}.

\begin{itemize}
    \item[\textbf{A.}] CPU
    \item[\textbf{B.}] control device
    \item[\textbf{C.}] \textcolor{red}{\textbf{peripheral}}
    \item[\textbf{D.}] register
\end{itemize}

\vspace{0.5cm}

\vspace{0.5cm}

\textbf{Question 9:} The I/O function includes a \underline{\hspace{1.5cm}} requirement to coordinate the flow of traffic between internal resources and external devices.

\begin{itemize}
    \item[\textbf{A.}] cycle
    \item[\textbf{B.}] status reporting
    \item[\textbf{C.}] \textcolor{red}{\textbf{control and timing}}
    \item[\textbf{D.}] data
\end{itemize}

\vspace{0.5cm}

\vspace{0.5cm}

\textbf{Question 18:} With isolated I/O there is a single address space for memory locations and I/O devices.

\begin{itemize}
    \item[\textbf{A.}] True
    \item[\textbf{B.}] \textcolor{red}{\textbf{False}}
\end{itemize}

\vspace{0.5cm}

\vspace{0.5cm}

\textbf{Question 55:} An I/O module must recognize one unique address for each peripheral it controls.

\begin{itemize}
    \item[\textbf{A.}] \textcolor{red}{\textbf{True}}
    \item[\textbf{B.}] False
\end{itemize}

\vspace{0.5cm}

\vspace{0.5cm}

\textbf{Question 61:} I/O channels are commonly seen on microcomputers, whereas I/O controllers are used on m

\begin{itemize}
    \item[\textbf{A.}] \textcolor{red}{\textbf{False}}
    \item[\textbf{B.}] True
\end{itemize}

\vspace{0.5cm}

\vspace{0.5cm}

\textbf{Question 72:} Because the 82C55A is programmable via the control register, it can be used to control variety of simple peripheral devices.

\begin{itemize}
    \item[\textbf{A.}] \textcolor{red}{\textbf{True}}
    \item[\textbf{B.}] False
\end{itemize}

\vspace{0.5cm}

\vspace{0.5cm}

\textbf{Question 102:} Which one of the following can be called as a peripheral?

\begin{itemize}
    \item[\textbf{A.}] Control Unit
    \item[\textbf{B.}] Arithmetic Unit
    \item[\textbf{C.}] \textcolor{red}{\textbf{Speakers}}
    \item[\textbf{D.}] Logic Unit
    \item[\textbf{E.}] Main Memory
\end{itemize}

\vspace{0.5cm}

\vspace{0.5cm}

\textbf{Question 125:} The \underline{\hspace{1.5cm}} command is used to activate a peripheral and tell it what to do.

\begin{itemize}
    \item[\textbf{A.}] \textcolor{red}{\textbf{control}}
    \item[\textbf{B.}] test
    \item[\textbf{C.}] read
    \item[\textbf{D.}] write
\end{itemize}

\vspace{0.5cm}

\vspace{0.5cm}

\textbf{Question 133:} What interface is used to connect the processor to I/O devices that require transmission of data one bit at a time?

\begin{itemize}
    \item[\textbf{A.}] Parallel
    \item[\textbf{B.}] \textcolor{red}{\textbf{Serial}}
    \item[\textbf{C.}] Output
    \item[\textbf{D.}] Input
    \item[\textbf{E.}] Bus
\end{itemize}

\vspace{0.5cm}

\vspace{0.5cm}

\textbf{Question 146:} An I/O module cannot exchange data directly with the processor.

\begin{itemize}
    \item[\textbf{A.}] True
    \item[\textbf{B.}] \textcolor{red}{\textbf{False}}
\end{itemize}

\vspace{0.5cm}

\vspace{0.5cm}

\textbf{Question 170:} A set of I/O modules is a key element of a computer system.

\begin{itemize}
    \item[\textbf{A.}] \textcolor{red}{\textbf{True}}
    \item[\textbf{B.}] False
\end{itemize}

\vspace{0.5cm}

\vspace{0.5cm}

\textbf{Question 174:} Swapping is an I/O operation.

\begin{itemize}
    \item[\textbf{A.}] \textcolor{red}{\textbf{True}}
    \item[\textbf{B.}] False
\end{itemize}

\vspace{0.5cm}

\vspace{0.5cm}

\textbf{Question 182:} The \underline{\hspace{1.5cm}} command causes the I/O module to take an item of data from the data bus and transmit that data item to the peripheral.

\begin{itemize}
    \item[\textbf{A.}] \textcolor{red}{\textbf{write}}
    \item[\textbf{B.}] test
    \item[\textbf{C.}] read
    \item[\textbf{D.}] control
\end{itemize}

\vspace{0.5cm}

\vspace{0.5cm}

\textbf{Question 211:} The \underline{\hspace{1.5cm}} command is used to activate a peripheral and tell it what to do.

\begin{itemize}
    \item[\textbf{A.}] read
    \item[\textbf{B.}] test
    \item[\textbf{C.}] write
    \item[\textbf{D.}] \textcolor{red}{\textbf{control}}
\end{itemize}

\vspace{0.5cm}

\vspace{0.5cm}

\textbf{Question 228:} An I/O module that takes on most of the detailed processing burden, presenting a high-level interface to the processor, is usually referred to as an \underline{\hspace{1.5cm}}.

\begin{itemize}
    \item[\textbf{A.}] I/O command
    \item[\textbf{B.}] \textcolor{red}{\textbf{I/O channel}}
    \item[\textbf{C.}] device controller
    \item[\textbf{D.}] I/O controller
\end{itemize}

\vspace{0.5cm}

\vspace{0.5cm}

\textbf{Question 237:} An I/O module that takes on most of the detailed processing burden, presenting a high-level interface to the processor, is usually referred to as an \underline{\hspace{1.5cm}}.

\begin{itemize}
    \item[\textbf{A.}] \textcolor{red}{\textbf{I/O channel}}
    \item[\textbf{B.}] I/O command
    \item[\textbf{C.}] I/O controller
    \item[\textbf{D.}] device controller
\end{itemize}

\vspace{0.5cm}

\vspace{0.5cm}

\textbf{Question 278:} In isolated I/O,\underline{\hspace{1.5cm}}

\begin{itemize}
    \item[\textbf{A.}] The I/O devices and the memory share the same address space
    \item[\textbf{B.}] \textcolor{red}{\textbf{The I/O devices have a separate address space from memory}}
    \item[\textbf{C.}] The memory and I/O devices have an associated address space
    \item[\textbf{D.}] A part of the memory is specifically set aside for the I/O operation
    \item[\textbf{E.}] None of the mentioned
\end{itemize}

\vspace{0.5cm}

\vspace{0.5cm}

\textbf{Question 367:} (1) An I/O module must recognize one unique address for each peripheral it controls. (2) I/O channels are commonly seen on microcomputers, whereas I/O controllers are used on mainframes. The statement (1) is \underline{\hspace{1.5cm}} and (2) is \underline{\hspace{1.5cm}}

\begin{itemize}
    \item[\textbf{A.}] \textcolor{red}{\textbf{true, false}}
    \item[\textbf{B.}] true, true
    \item[\textbf{C.}] false, true
    \item[\textbf{D.}] false, false
\end{itemize}

\vspace{0.5cm}

\vspace{0.5cm}

\textbf{Question 386:} What is the main distinction between Interrupt-Driven I/O and Direct Memory Access (DMA)?

\begin{itemize}
    \item[\textbf{A.}] \textcolor{red}{\textbf{Interrupt-Driven I/O involves the CPU in every data transfer, while DMA bypasses the CPU and transfers data directly between I/O device and memory}}
    \item[\textbf{B.}] Interrupt-Driven I/O requires special hardware and software support, while DMA does not need any additional components
    \item[\textbf{C.}] Interrupt-Driven I/O is suitable for small and frequent data transfers, while DMA is suitable for large and infrequent data transfers
    \item[\textbf{D.}] All of the mentioned
\end{itemize}

\vspace{0.5cm}

\vspace{0.5cm}

\textbf{Question 400:} What are the most basic functions of an I/O module?

\begin{itemize}
    \item[\textbf{A.}] \textcolor{red}{\textbf{Control and timing, processor communication, device communication, data buffering, error detection.}}
    \item[\textbf{B.}] Control and timing, processor communication, data buffering, error detection, writing data to memory.
    \item[\textbf{C.}] Control and timing, processor communication, device communication, data buffering, blocking device.
    \item[\textbf{D.}] Control and timing, processor communication, device communication, error detection, data processing.
\end{itemize}

\vspace{0.5cm}

\vspace{0.5cm}

\textbf{Question 401:} The I/O technique where the processor busy waits for an I/O operation to complete is called \underline{\hspace{1.5cm}}

\begin{itemize}
    \item[\textbf{A.}] Programmed I/O or DMA
    \item[\textbf{B.}] Interrupt-driven I/O
    \item[\textbf{C.}] Direct Memory Access (DMA)
    \item[\textbf{D.}] \textcolor{red}{\textbf{Programmed I/O}}
\end{itemize}

\vspace{0.5cm}

\vspace{0.5cm}

\textbf{Question 409:} What are basic registers that help CPU establish a connection with an I/O device?

\begin{itemize}
    \item[\textbf{A.}] \textcolor{red}{\textbf{I/OAR, I/OBR}}
    \item[\textbf{B.}] I/OAR, MAR
    \item[\textbf{C.}] I/OBR, MBR
    \item[\textbf{D.}] MAR, MBR
\end{itemize}

\vspace{0.5cm}

\vspace{0.5cm}

\textbf{Question 415:} The decision as to which process's pending I/O request shall be handled by an available I/O device

\begin{itemize}
    \item[\textbf{A.}] Medium-term scheduling
    \item[\textbf{B.}] Short-term scheduling
    \item[\textbf{C.}] \textcolor{red}{\textbf{I/O scheduling}}
    \item[\textbf{D.}] Long-term scheduling
\end{itemize}

\vspace{0.5cm}

\vspace{0.5cm}

\textbf{Question 467:} In which I/O technique does an I/O channel or processor manage transfer of data, thus freeing the CPU?

\begin{itemize}
    \item[\textbf{A.}] Programmed I/O
    \item[\textbf{B.}] Interrupt-driven I/O
    \item[\textbf{C.}] \textcolor{red}{\textbf{Channel I/O}}
    \item[\textbf{D.}] Memory-mapped I/O
\end{itemize}

\vspace{0.5cm}

\vspace{0.5cm}

\textbf{Question 468:} Which of the following is a DISADVANTAGE of using programmed I/O?

\begin{itemize}
    \item[\textbf{A.}] Increases system efficiency
    \item[\textbf{B.}] Reduces CPU involvement in I/O operations
    \item[\textbf{C.}] \textcolor{red}{\textbf{Requires the CPU to be constantly busy with I/O tasks}}
    \item[\textbf{D.}] Increases system complexity
\end{itemize}

\vspace{0.5cm}

\vspace{0.5cm}

\textbf{Question 489:} In programmed I/O, who does initialize the transfer of data?

\begin{itemize}
    \item[\textbf{A.}] I/O controller
    \item[\textbf{B.}] \textcolor{red}{\textbf{CPU}}
    \item[\textbf{C.}] DMA controller
    \item[\textbf{D.}] Operating system
\end{itemize}

\vspace{0.5cm}

\vspace{0.5cm}

\textbf{Question 528:} Which statement is correct in memory-mapped I/O?

\begin{itemize}
    \item[\textbf{A.}] \textcolor{red}{\textbf{The I/O devices and the memory share the same address space}}
    \item[\textbf{B.}] The I/O devices have a separate address space
    \item[\textbf{C.}] The memory and I/O devices have an associated address space
    \item[\textbf{D.}] A part of the memory is specifically set aside for the I/O operation
\end{itemize}

\vspace{0.5cm}

\vspace{0.5cm}

\textbf{Question 532:} The advantage of interrupt driven I/O over programmed I/O is \underline{\hspace{1.5cm}}.

\begin{itemize}
    \item[\textbf{A.}] the latter offers faster transfer of data
    \item[\textbf{B.}] \textcolor{red}{\textbf{the former does not have to wait until I/O operation is complete}}
    \item[\textbf{C.}] the devices have to deal with fewer address lines
    \item[\textbf{D.}] no advantage as such
\end{itemize}

\vspace{0.5cm}

\vspace{0.5cm}

\textbf{Question 533:} What is the main advantage of using DMA over Programmed I/O or Interrupt-Driven I/O?

\begin{itemize}
    \item[\textbf{A.}] DMA reduces the CPU involvement and overhead in data transfer
    \item[\textbf{B.}] DMA increases the data transfer rate and efficiency
    \item[\textbf{C.}] DMA allows concurrent execution of CPU and I/O operations
    \item[\textbf{D.}] \textcolor{red}{\textbf{All of the mentioned}}
\end{itemize}

\vspace{0.5cm}

\vspace{0.5cm}

\textbf{Question 6:} When data are moved over longer distances, to or from a remote device, the process is known as \underline{\hspace{1.5cm}}.

\begin{itemize}
    \item[\textbf{A.}] \textcolor{red}{\textbf{data communications}}
    \item[\textbf{B.}] registering
    \item[\textbf{C.}] structuring
    \item[\textbf{D.}] data transport
\end{itemize}

\vspace{0.5cm}

\vspace{0.5cm}

\textbf{Question 413:} \underline{\hspace{1.5cm}} can exchange data directly with the processor. Just as the processor can initiate a read or write with memory, designating the address of a specific location, the processor can also read data from or write data to \underline{\hspace{1.5cm}}.

\begin{itemize}
    \item[\textbf{A.}] \textcolor{red}{\textbf{I/O module / Memory}}
    \item[\textbf{B.}] Interrupts / I/O module
    \item[\textbf{C.}] I/O module / Interrupts
    \item[\textbf{D.}] Interrupts / Interrupts
    \item[\textbf{E.}] Memory / I/O module
\end{itemize}

\vspace{0.5cm}

\vspace{0.5cm}

\textbf{Question 13:} Cycle stealing is/are used in which concept?

\begin{itemize}
    \item[\textbf{A.}] Programmed I/O
    \item[\textbf{B.}] \textcolor{red}{\textbf{DMA}}
    \item[\textbf{C.}] Interrupts
    \item[\textbf{D.}] Memory mapped I/O
    \item[\textbf{E.}] All of the above
\end{itemize}

\vspace{0.5cm}

\vspace{0.5cm}

\textbf{Question 221:} Cycle stealing is/are used in which concept?

\begin{itemize}
    \item[\textbf{A.}] Programmed I/O
    \item[\textbf{B.}] \textcolor{red}{\textbf{DMA}}
    \item[\textbf{C.}] Interrupts
    \item[\textbf{D.}] Memory mapped I/O
    \item[\textbf{E.}] All of the above
\end{itemize}

\vspace{0.5cm}

\vspace{0.5cm}

\textbf{Question 466:} From a structural point of view, the external devices are controlled by \underline{\hspace{1.5cm}}.

\begin{itemize}
    \item[\textbf{A.}] Control signals
    \item[\textbf{B.}] Peripheral device
    \item[\textbf{C.}] \textcolor{red}{\textbf{I/O modules}}
    \item[\textbf{D.}] Remote devices
\end{itemize}

\vspace{0.5cm}

\vspace{0.5cm}

\textbf{Question 214:} A \underline{\hspace{1.5cm}} is a small, very-high-speed memory maintained by the instruction fetch stage of the pipeline and containing the n most recently fetched instructions in sequence.

\begin{itemize}
    \item[\textbf{A.}] multiple stream
    \item[\textbf{B.}] delayed branch
    \item[\textbf{C.}] branch prediction
    \item[\textbf{D.}] \textcolor{red}{\textbf{loop buffer}}
\end{itemize}

\vspace{0.5cm}

\vspace{0.5cm}

\textbf{Question 326:} Which registers can interact with the secondary storage?

\begin{itemize}
    \item[\textbf{A.}] \textcolor{red}{\textbf{MAR}}
    \item[\textbf{B.}] PC
    \item[\textbf{C.}] IR
    \item[\textbf{D.}] R0
    \item[\textbf{E.}] All of the mentioned
\end{itemize}

\vspace{0.5cm}

\vspace{0.5cm}

\textbf{Question 329:} For interrupts, all I/O modules share a common interrupt request line. When the processor senses an interrupt, it sends out an interrupt acknowledge. This signal propagates through a series of I/O modules until it gets to a requesting module. Which kind of interrupt technique is it?

\begin{itemize}
    \item[\textbf{A.}] Multiple interrupt lines
    \item[\textbf{B.}] Software poll
    \item[\textbf{C.}] \textcolor{red}{\textbf{Daisy chain}}
    \item[\textbf{D.}] Bus arbitration
\end{itemize}

\vspace{0.5cm}

\vspace{0.5cm}

\textbf{Question 157:} When large volumes of data are to be moved, a more efficient technique is direct memory ac

\begin{itemize}
    \item[\textbf{A.}] \textcolor{red}{\textbf{True}}
    \item[\textbf{B.}] False
\end{itemize}

\vspace{0.5cm}

\vspace{0.5cm}

\textbf{Question 261:} In the computer, what categories do external devices include? (choose 3 correct answers)

\begin{itemize}
    \item[\textbf{A.}] \textcolor{red}{\textbf{Human readable}}
    \item[\textbf{B.}] \textcolor{red}{\textbf{Communication}}
    \item[\textbf{C.}] Data Conversion
    \item[\textbf{D.}] \textcolor{red}{\textbf{Machine readable}}
\end{itemize}

\vspace{0.5cm}

\vspace{0.5cm}

\textbf{Question 286:} The hardware mechanism that allows a device to notify the CPU is called \underline{\hspace{1.5cm}} .

\begin{itemize}
    \item[\textbf{A.}] polling
    \item[\textbf{B.}] \textcolor{red}{\textbf{interrupt}}
    \item[\textbf{C.}] driver
    \item[\textbf{D.}] controlling
\end{itemize}

\vspace{0.5cm}

\vspace{0.5cm}

\textbf{Question 381:} The instruction, which adds 1 to thevalue in a memory location, has five stages: fetch opcode (four cycles), fetch operand address (three cycles), fetch operand (three cycles), add 1 to operand (three cycles), and store operand (three cycles). An interrupt sends request at beginning of fetch operand stage. How many cycles does the processor enter the interrupt processing cycle?

\begin{itemize}
    \item[\textbf{A.}] 6
    \item[\textbf{B.}] 7
    \item[\textbf{C.}] 8
    \item[\textbf{D.}] 9
    \item[\textbf{E.}] \textcolor{red}{\textbf{10}}
\end{itemize}

\vspace{0.5cm}

\vspace{0.5cm}

\textbf{Question 456:} What is the primary role of the IO module? (choose two correct answers)

\begin{itemize}
    \item[\textbf{A.}] \textcolor{red}{\textbf{Its role is data transfer, it will transfer data to and from peripheral devices.}}
    \item[\textbf{B.}] \textcolor{red}{\textbf{Its role is device communication control, it manages and controls the flow of data between the CPU and peripherals}}
    \item[\textbf{C.}] Its role is data formatting and converting analog signals into digital audio data for CPU
    \item[\textbf{D.}] Its role is protecting data from user, ensuring that sensitive information is not lost or intercepted.
\end{itemize}

\vspace{0.5cm}

\vspace{0.5cm}

\section{Chapter 09: Operating System Support}
\textit{Total questions in chapter: 43}

\textbf{Question 23:} The OS must determine how much processor time is to be devoted to the execution of a particular user program.

\begin{itemize}
    \item[\textbf{A.}] \textcolor{red}{\textbf{True}}
    \item[\textbf{B.}] False
\end{itemize}

\vspace{0.5cm}

\vspace{0.5cm}

\textbf{Question 24:} The most important system program is the OS.

\begin{itemize}
    \item[\textbf{A.}] \textcolor{red}{\textbf{True}}
    \item[\textbf{B.}] False
\end{itemize}

\vspace{0.5cm}

\vspace{0.5cm}

\textbf{Question 34:} With a batch operating system the user does not have direct access to the processor.

\begin{itemize}
    \item[\textbf{A.}] \textcolor{red}{\textbf{True}}
    \item[\textbf{B.}] False
\end{itemize}

\vspace{0.5cm}

\vspace{0.5cm}

\textbf{Question 93:} \underline{\hspace{1.5cm}} is when the processor spends most of its time swapping pages rather than execut

\begin{itemize}
    \item[\textbf{A.}] Swapping
    \item[\textbf{B.}] \textcolor{red}{\textbf{Thrashing}}
    \item[\textbf{C.}] Paging
    \item[\textbf{D.}] Multitasking
\end{itemize}

\vspace{0.5cm}

\vspace{0.5cm}

\textbf{Question 111:} In a system without virtual memory, the effective address is a virtual address or a register.

\begin{itemize}
    \item[\textbf{A.}] True
    \item[\textbf{B.}] \textcolor{red}{\textbf{False}}
\end{itemize}

\vspace{0.5cm}

\vspace{0.5cm}

\textbf{Question 151:} With demand paging it is necessary to load an entire process into main memory.

\begin{itemize}
    \item[\textbf{A.}] True
    \item[\textbf{B.}] \textcolor{red}{\textbf{False}}
\end{itemize}

\vspace{0.5cm}

\vspace{0.5cm}

\textbf{Question 200:} The \underline{\hspace{1.5cm}} scheduler is also known as the dispatcher.

\begin{itemize}
    \item[\textbf{A.}] medium-term
    \item[\textbf{B.}] I/O
    \item[\textbf{C.}] \textcolor{red}{\textbf{short-term}}
    \item[\textbf{D.}] long-term
\end{itemize}

\vspace{0.5cm}

\vspace{0.5cm}

\textbf{Question 220:} \underline{\hspace{1.5cm}} is when the processor spends most of its time swapping pages rather than executing instructions.

\begin{itemize}
    \item[\textbf{A.}] Swapping
    \item[\textbf{B.}] Paging
    \item[\textbf{C.}] \textcolor{red}{\textbf{Thrashing}}
    \item[\textbf{D.}] Multitasking
\end{itemize}

\vspace{0.5cm}

\vspace{0.5cm}

\textbf{Question 254:} The \underline{\hspace{1.5cm}} scheduler is also known as the dispatcher.

\begin{itemize}
    \item[\textbf{A.}] long-term
    \item[\textbf{B.}] medium-term
    \item[\textbf{C.}] \textcolor{red}{\textbf{short-term}}
    \item[\textbf{D.}] I/O
\end{itemize}

\vspace{0.5cm}

\vspace{0.5cm}

\textbf{Question 279:} Which component defines the system call interface to the operating system and facilitates binary portability?

\begin{itemize}
    \item[\textbf{A.}] \textcolor{red}{\textbf{Application Binary Interface}}
    \item[\textbf{B.}] Application Programming Interface
    \item[\textbf{C.}] Instruction Set Architecture
    \item[\textbf{D.}] Central Processing Unit
\end{itemize}

\vspace{0.5cm}

\vspace{0.5cm}

\textbf{Question 280:} What is the initial state of a process when it is admitted by the high-level scheduler, but not yet ready to execute?

\begin{itemize}
    \item[\textbf{A.}] \textcolor{red}{\textbf{New}}
    \item[\textbf{B.}] Ready
    \item[\textbf{C.}] Running
    \item[\textbf{D.}] Halted
\end{itemize}

\vspace{0.5cm}

\vspace{0.5cm}

\textbf{Question 331:} Which state indicates that a process is currently being executed by the processor?

\begin{itemize}
    \item[\textbf{A.}] \textcolor{red}{\textbf{Running}}
    \item[\textbf{B.}] Ready
    \item[\textbf{C.}] NewBorn
    \item[\textbf{D.}] Halted
\end{itemize}

\vspace{0.5cm}

\vspace{0.5cm}

\textbf{Question 332:} The task of subdivision is carried out dynamically by the OS and is known as\underline{\hspace{1.5cm}} ?

\begin{itemize}
    \item[\textbf{A.}] Scheduling
    \item[\textbf{B.}] \textcolor{red}{\textbf{Memory management}}
    \item[\textbf{C.}] Virtual Memory
    \item[\textbf{D.}] All of the mentioned
\end{itemize}

\vspace{0.5cm}

\vspace{0.5cm}

\textbf{Question 389:} We have a long-term queue of process requests, typically stored on \underline{\hspace{1.5cm}}.

\begin{itemize}
    \item[\textbf{A.}] main memory
    \item[\textbf{B.}] \textcolor{red}{\textbf{disk}}
    \item[\textbf{C.}] cache memory
    \item[\textbf{D.}] registers
\end{itemize}

\vspace{0.5cm}

\vspace{0.5cm}

\textbf{Question 425:} Facilities and services provided by the OS that assist the programmer in creating programs are in the form of \_ programs that are not actually part of the OS but are accessible through the OS.

\begin{itemize}
    \item[\textbf{A.}] \textcolor{red}{\textbf{utility}}
    \item[\textbf{B.}] logical address
    \item[\textbf{C.}] JCL
    \item[\textbf{D.}] multitasking
\end{itemize}

\vspace{0.5cm}

\vspace{0.5cm}

\textbf{Question 450:} \underline{\hspace{1.5cm}} is when the processor spends most of its time swapping pages rather than executing instructions.

\begin{itemize}
    \item[\textbf{A.}] Multitasking
    \item[\textbf{B.}] Swapping
    \item[\textbf{C.}] \textcolor{red}{\textbf{Thrashing}}
    \item[\textbf{D.}] Paging
\end{itemize}

\vspace{0.5cm}

\vspace{0.5cm}

\textbf{Question 19:} The \underline{\hspace{1.5cm}} defines the system call interface to the operating system and the hardware resources and services available in a system through the user instruction set architecture.

\begin{itemize}
    \item[\textbf{A.}] HLL
    \item[\textbf{B.}] API
    \item[\textbf{C.}] \textcolor{red}{\textbf{ABI}}
    \item[\textbf{D.}] ISA
\end{itemize}

\vspace{0.5cm}

\vspace{0.5cm}

\textbf{Question 27:} The ABI is the boundary between hardware and software.

\begin{itemize}
    \item[\textbf{A.}] True
    \item[\textbf{B.}] \textcolor{red}{\textbf{False}}
\end{itemize}

\vspace{0.5cm}

\vspace{0.5cm}

\textbf{Question 43:} Both batch multiprogramming and time sharing use multiprogramming.

\begin{itemize}
    \item[\textbf{A.}] False
    \item[\textbf{B.}] \textcolor{red}{\textbf{True}}
\end{itemize}

\vspace{0.5cm}

\vspace{0.5cm}

\textbf{Question 86:} Privileged instructions are certain instructions that are designated special and can be executed only by the monitor.

\begin{itemize}
    \item[\textbf{A.}] \textcolor{red}{\textbf{True}}
    \item[\textbf{B.}] False
\end{itemize}

\vspace{0.5cm}

\vspace{0.5cm}

\textbf{Question 166:} Techniques that automatically move program and data blocks into the physical main memory when they are required for execution are called \underline{\hspace{1.5cm}}.

\begin{itemize}
    \item[\textbf{A.}] Associative-Mapping techniques
    \item[\textbf{B.}] Main Memory techniques
    \item[\textbf{C.}] \textcolor{red}{\textbf{Virtual Memory techniques}}
    \item[\textbf{D.}] Cache Memory techniques
    \item[\textbf{E.}] Paging techniques
\end{itemize}

\vspace{0.5cm}

\vspace{0.5cm}

\textbf{Question 189:} The \underline{\hspace{1.5cm}} is a program that controls the execution of application programs and acts as an interface between applications and the computer hardware.

\begin{itemize}
    \item[\textbf{A.}] nucleus
    \item[\textbf{B.}] batch system
    \item[\textbf{C.}] \textcolor{red}{\textbf{operating system}}
    \item[\textbf{D.}] job control language
\end{itemize}

\vspace{0.5cm}

\vspace{0.5cm}

\textbf{Question 218:} Techniques that automatically move program and data blocks into the physical main memory when they are required for execution are called \underline{\hspace{1.5cm}}.

\begin{itemize}
    \item[\textbf{A.}] Associative-Mapping techniques
    \item[\textbf{B.}] Main Memory techniques
    \item[\textbf{C.}] Cache Memory techniques
    \item[\textbf{D.}] Paging techniques
    \item[\textbf{E.}] \textcolor{red}{\textbf{Virtual-memory techniques}}
\end{itemize}

\vspace{0.5cm}

\vspace{0.5cm}

\textbf{Question 269:} For reads to and writes from main memory, \underline{\hspace{1.5cm}} translates each virtual address into a physical address in main memory.

\begin{itemize}
    \item[\textbf{A.}] MAR
    \item[\textbf{B.}] \textcolor{red}{\textbf{MMU}}
    \item[\textbf{C.}] Overlays
    \item[\textbf{D.}] TLB
    \item[\textbf{E.}] Accumulator
\end{itemize}

\vspace{0.5cm}

\vspace{0.5cm}

\textbf{Question 402:} Which method allows the programmer to view memory as consisting of multiple address spaces and is used to map logical addresses of variable length onto physical memory?

\begin{itemize}
    \item[\textbf{A.}] Paging
    \item[\textbf{B.}] Overlays
    \item[\textbf{C.}] \textcolor{red}{\textbf{Segmentation}}
    \item[\textbf{D.}] Paging with segmentation
\end{itemize}

\vspace{0.5cm}

\vspace{0.5cm}

\textbf{Question 491:} The \underline{\hspace{1.5cm}} scheduler executes relatively infrequently and makes the coarse-grained decision of whether or not to take on a new process, and which one to take. The \underline{\hspace{1.5cm}} scheduler, also known as the dispatcher, executes frequently and makes the fine-grained decision of which job to execute next.

\begin{itemize}
    \item[\textbf{A.}] short-term; medium-term
    \item[\textbf{B.}] medium-term; short-term
    \item[\textbf{C.}] \textcolor{red}{\textbf{long-term; short-term}}
    \item[\textbf{D.}] long-term; medium-term
\end{itemize}

\vspace{0.5cm}

\vspace{0.5cm}

\textbf{Question 520:} Which technique involves dividing physical memory into fixed-size or variable-size blocks to allocate memory to processes?

\begin{itemize}
    \item[\textbf{A.}] \textcolor{red}{\textbf{Partitioning}}
    \item[\textbf{B.}] Swapping
    \item[\textbf{C.}] Paging
    \item[\textbf{D.}] Translation Lookaside Buffer
\end{itemize}

\vspace{0.5cm}

\vspace{0.5cm}

\textbf{Question 281:} The chunks of a program, known as pages, could be assigned to available chunks of memory, known as frames, is called \underline{\hspace{1.5cm}}.

\begin{itemize}
    \item[\textbf{A.}] Swapping
    \item[\textbf{B.}] Partitioning
    \item[\textbf{C.}] \textcolor{red}{\textbf{Paging}}
    \item[\textbf{D.}] Virtual Memory
    \item[\textbf{E.}] Segmentation
\end{itemize}

\vspace{0.5cm}

\vspace{0.5cm}

\textbf{Question 333:} Which of the following statements is incorrect about Translation Look-aside Buffer (TLB)?

\begin{itemize}
    \item[\textbf{A.}] \textcolor{red}{\textbf{The use of TLB eliminates the need for keeping a page table in memory}}
    \item[\textbf{B.}] TLB only maintains a subset of the entries stored in the full memory-based page table
    \item[\textbf{C.}] When there is a TLB miss the system needs to access the page table
    \item[\textbf{D.}] A translation lookaside buffer (TLB) is a memory cache that stores the recent translations of virtual memory to physical memory
\end{itemize}

\vspace{0.5cm}

\vspace{0.5cm}

\textbf{Question 109:} Scheduling and memory management are the two OS functions that are most relevant to the computer organization and architecture.

\begin{itemize}
    \item[\textbf{A.}] \textcolor{red}{\textbf{True}}
    \item[\textbf{B.}] False
\end{itemize}

\vspace{0.5cm}

\vspace{0.5cm}

\textbf{Question 446:} The \underline{\hspace{1.5cm}} gives a program access to the hardware resources and services available in a system through the user instruction set architecture supplemented with high-level language library calls.

\begin{itemize}
    \item[\textbf{A.}] \textcolor{red}{\textbf{API}}
    \item[\textbf{B.}] ISA
    \item[\textbf{C.}] JCL
    \item[\textbf{D.}] ABI
\end{itemize}

\vspace{0.5cm}

\vspace{0.5cm}

\textbf{Question 20:} It is common for programs, both system and application, to continue to exhibit new bugs after years of operation.

\begin{itemize}
    \item[\textbf{A.}] \textcolor{red}{\textbf{True}}
    \item[\textbf{B.}] False
\end{itemize}

\vspace{0.5cm}

\vspace{0.5cm}

\textbf{Question 4:} A \underline{\hspace{1.5cm}} is an actual location in main memory.

\begin{itemize}
    \item[\textbf{A.}] logical address
    \item[\textbf{B.}] partition address
    \item[\textbf{C.}] base address
    \item[\textbf{D.}] \textcolor{red}{\textbf{physical address}}
\end{itemize}

\vspace{0.5cm}

\vspace{0.5cm}

\textbf{Question 104:} INT 21h / AH=1 is \underline{\hspace{1.5cm}}

\begin{itemize}
    \item[\textbf{A.}] write character from standard input, entry is stored in AL.
    \item[\textbf{B.}] \textcolor{red}{\textbf{read character from standard input, with echo, result is stored in AL.}}
    \item[\textbf{C.}] write character from standard input, entry is stored in AH.
    \item[\textbf{D.}] read character from standard input, with echo, result is stored in AH.
\end{itemize}

\vspace{0.5cm}

\vspace{0.5cm}

\textbf{Question 129:} A \underline{\hspace{1.5cm}} system works only one program at a time.

\begin{itemize}
    \item[\textbf{A.}] batch
    \item[\textbf{B.}] \textcolor{red}{\textbf{uniprogramming}}
    \item[\textbf{C.}] kernel
    \item[\textbf{D.}] privileged instruction
\end{itemize}

\vspace{0.5cm}

\vspace{0.5cm}

\textbf{Question 131:} A \underline{\hspace{1.5cm}} is a special type of programming language used to provide instructions to the monitor.

\begin{itemize}
    \item[\textbf{A.}] \textcolor{red}{\textbf{job control language}}
    \item[\textbf{B.}] multiprogram
    \item[\textbf{C.}] kernel
    \item[\textbf{D.}] utility
\end{itemize}

\vspace{0.5cm}

\vspace{0.5cm}

\textbf{Question 155:} The \underline{\hspace{1.5cm}} scheduler determines which programs are admitted to the system for processing

\begin{itemize}
    \item[\textbf{A.}] \textcolor{red}{\textbf{long-term}}
    \item[\textbf{B.}] medium-term
    \item[\textbf{C.}] short-term
    \item[\textbf{D.}] I/O
\end{itemize}

\vspace{0.5cm}

\vspace{0.5cm}

\textbf{Question 252:} Uniprogramming is the central theme of modern operating systems.

\begin{itemize}
    \item[\textbf{A.}] True
    \item[\textbf{B.}] \textcolor{red}{\textbf{False}}
\end{itemize}

\vspace{0.5cm}

\vspace{0.5cm}

\textbf{Question 289:} In MASM32, which OPCODE is used to compare two values?

\begin{itemize}
    \item[\textbf{A.}] COM
    \item[\textbf{B.}] \textcolor{red}{\textbf{CMP}}
    \item[\textbf{C.}] IF ... ELSE
    \item[\textbf{D.}] TEST
\end{itemize}

\vspace{0.5cm}

\vspace{0.5cm}

\textbf{Question 330:} What role does an Application Programming Interface (API) play in software development?

\begin{itemize}
    \item[\textbf{A.}] \textcolor{red}{\textbf{It allows program access to hardware resources using high-level language libraries}}
    \item[\textbf{B.}] It defines low-level machine instructions
    \item[\textbf{C.}] It provides a standard for binary portability
    \item[\textbf{D.}] It manages system resources for the operating system and machine language instructions
\end{itemize}

\vspace{0.5cm}

\vspace{0.5cm}

\textbf{Question 469:} What is the purpose of using an API?

\begin{itemize}
    \item[\textbf{A.}] \textcolor{red}{\textbf{To enable application software to be ported easily to other systems that support the same API}}
    \item[\textbf{B.}] To define the repertoire of machine language instructions
    \item[\textbf{C.}] To standardize binary portability across programs
    \item[\textbf{D.}] To manage system resources
\end{itemize}

\vspace{0.5cm}

\vspace{0.5cm}

\textbf{Question 492:} If a program accesses data in a specified file, how to manage the file's properties such as disk position, file's attributes?

\begin{itemize}
    \item[\textbf{A.}] \textcolor{red}{\textbf{They are managed by the current operating system.}}
    \item[\textbf{B.}] They are managed using variables of the running program.
    \item[\textbf{C.}] They are stored on the hard disk.
    \item[\textbf{D.}] They are managed by the programmer.
\end{itemize}

\vspace{0.5cm}

\vspace{0.5cm}

\textbf{Question 515:} In the MASM, how to receive numbers from user?

\begin{itemize}
    \item[\textbf{A.}] \textcolor{red}{\textbf{Raw data from keyboard are string. The function sval(num-string) will convert num-string to signed number.}}
    \item[\textbf{B.}] Raw data from keyboard are string. The function str\$(num-string) will convert num-string to signed number.
    \item[\textbf{C.}] Raw data from keyboard are string. The function chr\$(num-string) will convert num-string to signed number.
    \item[\textbf{D.}] Raw data from keyboard are string. The function chr\$(num-string) will convert num-string to unsigned number.
\end{itemize}

\vspace{0.5cm}

\vspace{0.5cm}

\section{Chapter 10: Number Systems}
\textit{Total questions in chapter: 12}

\textbf{Question 418:} Convert the 64 from decimal to their binary number equivalents.

\begin{itemize}
    \item[\textbf{A.}] \textcolor{red}{\textbf{100000}}
    \item[\textbf{B.}] 100100
    \item[\textbf{C.}] 111000
    \item[\textbf{D.}] 101010
\end{itemize}

\vspace{0.5cm}

\vspace{0.5cm}

\textbf{Question 395:} Consider an expression: 1000 0000 OR 1111 1010 What is the result of this expression?

\begin{itemize}
    \item[\textbf{A.}] 1000 0000
    \item[\textbf{B.}] \textcolor{red}{\textbf{1111 1010}}
    \item[\textbf{C.}] 1001 1010
    \item[\textbf{D.}] 1001 0101
\end{itemize}

\vspace{0.5cm}

\vspace{0.5cm}

\textbf{Question 119:} A computer has memory of 256k words, how many bits are required to specify the address p

\begin{itemize}
    \item[\textbf{A.}] 16 bits
    \item[\textbf{B.}] 12 bits
    \item[\textbf{C.}] 8 bits
    \item[\textbf{D.}] \textcolor{red}{\textbf{18 bits}}
    \item[\textbf{E.}] 10 bits
\end{itemize}

\vspace{0.5cm}

\vspace{0.5cm}

\textbf{Question 284:} Express a sign integer number (+18) in the sign magnitude representation.

\begin{itemize}
    \item[\textbf{A.}] \textcolor{red}{\textbf{00010010}}
    \item[\textbf{B.}] 10010010
    \item[\textbf{C.}] 11110010
    \item[\textbf{D.}] 01110010
\end{itemize}

\vspace{0.5cm}

\vspace{0.5cm}

\textbf{Question 312:} Which representation is most efficient to perform arithmetic operations on the signed integer numbers?

\begin{itemize}
    \item[\textbf{A.}] Sign-magnitude
    \item[\textbf{B.}] \textcolor{red}{\textbf{2's complement}}
    \item[\textbf{C.}] 1's \& 2's compliment
    \item[\textbf{D.}] 1's complement
\end{itemize}

\vspace{0.5cm}

\vspace{0.5cm}

\textbf{Question 336:} Express an integer number +18 (using 8-bits length) in twos complement representation.

\begin{itemize}
    \item[\textbf{A.}] \textcolor{red}{\textbf{00010010}}
    \item[\textbf{B.}] 10010010
    \item[\textbf{C.}] 00001101
    \item[\textbf{D.}] 10011101
\end{itemize}

\vspace{0.5cm}

\vspace{0.5cm}

\textbf{Question 339:} What is the output of Left Shift Operator << on (00011000<<2)?

\begin{itemize}
    \item[\textbf{A.}] \textcolor{red}{\textbf{01100000}}
    \item[\textbf{B.}] 11000000
    \item[\textbf{C.}] 00000110
    \item[\textbf{D.}] 00000011
\end{itemize}

\vspace{0.5cm}

\vspace{0.5cm}

\textbf{Question 358:} What is the maximum addressable memory of a 32-bit microprocessor with 24-bit address?

\begin{itemize}
    \item[\textbf{A.}] 16 GB
    \item[\textbf{B.}] \textcolor{red}{\textbf{16 MB}}
    \item[\textbf{C.}] 16 Gbits
    \item[\textbf{D.}] 16 Mbits
\end{itemize}

\vspace{0.5cm}

\vspace{0.5cm}

\textbf{Question 371:} What is result left rotate of 10110101 by 2 bit?

\begin{itemize}
    \item[\textbf{A.}] 01101011
    \item[\textbf{B.}] 10101101
    \item[\textbf{C.}] 01101101
    \item[\textbf{D.}] \textcolor{red}{\textbf{11010110}}
    \item[\textbf{E.}] 11010100
\end{itemize}

\vspace{0.5cm}

\vspace{0.5cm}

\textbf{Question 403:} What is result logical right shift of 10110000 by 2 bit?

\begin{itemize}
    \item[\textbf{A.}] \textcolor{red}{\textbf{00101100}}
    \item[\textbf{B.}] 01011000
    \item[\textbf{C.}] 01001100
    \item[\textbf{D.}] 00110000
\end{itemize}

\vspace{0.5cm}

\vspace{0.5cm}

\textbf{Question 417:} If you have an integer number +18 in sign magnitude representation, which is 00010010, what is the correct option for -18?

\begin{itemize}
    \item[\textbf{A.}] 00010010
    \item[\textbf{B.}] \textcolor{red}{\textbf{10010010}}
    \item[\textbf{C.}] 11110010
    \item[\textbf{D.}] 01110010
\end{itemize}

\vspace{0.5cm}

\vspace{0.5cm}

\textbf{Question 537:} In an arithmetic left shift operation, what is appended to the rightmost bits?

\begin{itemize}
    \item[\textbf{A.}] \textcolor{red}{\textbf{Zeros}}
    \item[\textbf{B.}] Ones
    \item[\textbf{C.}] The original leftmost bit
    \item[\textbf{D.}] Random bits
\end{itemize}

\vspace{0.5cm}

\vspace{0.5cm}

\section{Chapter 11: Computer Arithmetic}
\textit{Total questions in chapter: 5}

\textbf{Question 334:} (R1) = 01110110, (R2) = 11011111, the result of (R1) XOR (R2) is:

\begin{itemize}
    \item[\textbf{A.}] 11011011
    \item[\textbf{B.}] 00010110
    \item[\textbf{C.}] \textcolor{red}{\textbf{10101001}}
    \item[\textbf{D.}] 11001101
\end{itemize}

\vspace{0.5cm}

\vspace{0.5cm}

\textbf{Question 343:} What is the primary function of the Arithmetic and Logic Unit (ALU) in a processor?

\begin{itemize}
    \item[\textbf{A.}] \textcolor{red}{\textbf{Perform actual computations and data processing}}
    \item[\textbf{B.}] Control the movement of data and instructions
    \item[\textbf{C.}] Act as an interface to the system bus
    \item[\textbf{D.}] Manage the internal processor memory
\end{itemize}

\vspace{0.5cm}

\vspace{0.5cm}

\textbf{Question 370:} What is result of 10100101 xor 11001001?

\begin{itemize}
    \item[\textbf{A.}] 11101101
    \item[\textbf{B.}] 10000001
    \item[\textbf{C.}] \textcolor{red}{\textbf{01101100}}
    \item[\textbf{D.}] 10101100
\end{itemize}

\vspace{0.5cm}

\vspace{0.5cm}

\textbf{Question 379:} Regarding the ALU(Arithmetic Logic Unit), besides basic arithmetic operations, what operations can it perform? (choose two correct answers)

\begin{itemize}
    \item[\textbf{A.}] \textcolor{red}{\textbf{It handles logical operations such as AND, OR, XOR, NOT.}}
    \item[\textbf{B.}] It handles data transfer operations like MOVE, GO, JUMP.
    \item[\textbf{C.}] It handles decoding operations after an instruction is fetched.
    \item[\textbf{D.}] \textcolor{red}{\textbf{It handles bit shifting operations like multi and div operations by powers of two.}}
\end{itemize}

\vspace{0.5cm}

\vspace{0.5cm}

\textbf{Question 316:} Consider an expression: NOT (1111 1010) What is the result of this expression?

\begin{itemize}
    \item[\textbf{A.}] 0000 1010
    \item[\textbf{B.}] \textcolor{red}{\textbf{0000 0101}}
    \item[\textbf{C.}] 1111 0101
    \item[\textbf{D.}] 1111 1010
\end{itemize}

\vspace{0.5cm}

\vspace{0.5cm}

\section{Chapter 12: Digital Logic}
\textit{Total questions in chapter: 46}

\textbf{Question 10:} The operation \underline{\hspace{1.5cm}} yields true if either or both of its operands are true.

\begin{itemize}
    \item[\textbf{A.}] NOT
    \item[\textbf{B.}] AND
    \item[\textbf{C.}] NAND
    \item[\textbf{D.}] \textcolor{red}{\textbf{OR}}
\end{itemize}

\vspace{0.5cm}

\vspace{0.5cm}

\textbf{Question 105:} A \underline{\hspace{1.5cm}} is an electronic circuit that produces an output signal that is a simple Boolean op signals.

\begin{itemize}
    \item[\textbf{A.}] \textcolor{red}{\textbf{gate}}
    \item[\textbf{B.}] decoder
    \item[\textbf{C.}] counter
    \item[\textbf{D.}] flip-flop
\end{itemize}

\vspace{0.5cm}

\vspace{0.5cm}

\textbf{Question 217:} The operation \underline{\hspace{1.5cm}} yields true if either or both of its operands are true.

\begin{itemize}
    \item[\textbf{A.}] AND
    \item[\textbf{B.}] NOT
    \item[\textbf{C.}] \textcolor{red}{\textbf{OR}}
    \item[\textbf{D.}] NAND
\end{itemize}

\vspace{0.5cm}

\vspace{0.5cm}

\textbf{Question 230:} \underline{\hspace{1.5cm}} are used in digital circuits to control signal and data routing.

\begin{itemize}
    \item[\textbf{A.}] \textcolor{red}{\textbf{Multiplexers}}
    \item[\textbf{B.}] Program counters
    \item[\textbf{C.}] Flip-flops
    \item[\textbf{D.}] Gates
\end{itemize}

\vspace{0.5cm}

\vspace{0.5cm}

\textbf{Question 276:} For the following Boolean expression $AB + AB'$ and the truth table:

\begin{center}
\renewcommand{\arraystretch}{1.3}
\begin{tabular}{|c|c|c|c|}
\hline
\textbf{A} & \textbf{B} & \textbf{B'} & \textbf{Output} \\ \hline
1 & 0 & 1 & ? \\ \hline
1 & 1 & 0 & ? \\ \hline
0 & 0 & 1 & ? \\ \hline
0 & 1 & 0 & ? \\ \hline
\end{tabular}
\end{center}

Choose the correct option to replace at ``?'' (order top to bottom):

\begin{itemize}
    \item[\textbf{A.}] \textcolor{red}{\textbf{1,1,1,0}}
    \item[\textbf{B.}] 1,1,1,1
    \item[\textbf{C.}] 1,0,1,0
    \item[\textbf{D.}] 0,1,1,0
\end{itemize}

\vspace{0.5cm}

\textbf{Question 277:} If you have a boolean function with 3 variables, how may rows are there in the truth table?

\begin{itemize}
    \item[\textbf{A.}] \textcolor{red}{\textbf{8 rows}}
    \item[\textbf{B.}] 3 rows
    \item[\textbf{C.}] 6 rows
    \item[\textbf{D.}] 12 rows
\end{itemize}

\vspace{0.5cm}

\vspace{0.5cm}

\textbf{Question 282:} How does Boolean algebra contribute to the design of digital circuits?

\begin{itemize}
    \item[\textbf{A.}] \textcolor{red}{\textbf{It simplifies the implementation of desired functions}}
    \item[\textbf{B.}] It helps in the analysis of economic data
    \item[\textbf{C.}] It facilitates the design of analog circuits
    \item[\textbf{D.}] It is primarily used for chemical engineering and physical engineering
\end{itemize}

\vspace{0.5cm}

\vspace{0.5cm}

\textbf{Question 283:} When both inputs are 1, what is the result of a NAND gate?

\begin{itemize}
    \item[\textbf{A.}] \textcolor{red}{\textbf{0}}
    \item[\textbf{B.}] 1
    \item[\textbf{C.}] 2
    \item[\textbf{D.}] Undefined
    \item[\textbf{E.}] \#NA
\end{itemize}

\vspace{0.5cm}

\vspace{0.5cm}

\textbf{Question 335:} What is the output of a NOT gate when the input is 0?

\begin{itemize}
    \item[\textbf{A.}] 0
    \item[\textbf{B.}] \textcolor{red}{\textbf{1}}
    \item[\textbf{C.}] Undefined
    \item[\textbf{D.}] 2
\end{itemize}

\vspace{0.5cm}

\vspace{0.5cm}

\textbf{Question 366:} Consider the expression : NOT(A + B) = ? Apply DeMorgan's Theorem to replace at "?"

\begin{itemize}
    \item[\textbf{A.}] \textcolor{red}{\textbf{NOT A AND NOT B}}
    \item[\textbf{B.}] NOT A OR NOT B
    \item[\textbf{C.}] NOT A AND B
    \item[\textbf{D.}] A OR NOT B
\end{itemize}

\vspace{0.5cm}

\vspace{0.5cm}

\textbf{Question 385:} Assume the following truth table:

\begin{center}
\renewcommand{\arraystretch}{1.3}
\begin{tabular}{|c|c|c|}
\hline
\textbf{X} & \textbf{Y} & \textbf{OUTPUT} \\ \hline
1 & 1 & 1 \\ \hline
1 & 0 & 1 \\ \hline
0 & 1 & 1 \\ \hline
0 & 0 & 0 \\ \hline
\end{tabular}
\end{center}

Which basic operator matches the table above?

\begin{itemize}
    \item[\textbf{A.}] AND
    \item[\textbf{B.}] \textcolor{red}{\textbf{OR}}
    \item[\textbf{C.}] NAND
    \item[\textbf{D.}] NOR
\end{itemize}

\vspace{0.5cm}

\textbf{Question 387:} When both inputs are 0, what is the result of a NOR gate?

\begin{itemize}
    \item[\textbf{A.}] 0
    \item[\textbf{B.}] \textcolor{red}{\textbf{1}}
    \item[\textbf{C.}] 2
    \item[\textbf{D.}] Undefined
    \item[\textbf{E.}] \#NA
\end{itemize}

\vspace{0.5cm}

\vspace{0.5cm}

\textbf{Question 408:} Consider an expression: 1000 0000 AND 1111 1010 What is the result of this expression?

\begin{itemize}
    \item[\textbf{A.}] \textcolor{red}{\textbf{1000 0000}}
    \item[\textbf{B.}] 1111 0000
    \item[\textbf{C.}] 1001 1010
    \item[\textbf{D.}] 1001 0101
\end{itemize}

\vspace{0.5cm}

\vspace{0.5cm}

\textbf{Question 414:} Consider an expression: NOT (A OR B) Choose the correct expression that is equal to the given expression.

\begin{itemize}
    \item[\textbf{A.}] \textcolor{red}{\textbf{A NOR B}}
    \item[\textbf{B.}] A NAND B
    \item[\textbf{C.}] NOT A NOR B
    \item[\textbf{D.}] NOT A OR B
\end{itemize}

\vspace{0.5cm}

\vspace{0.5cm}

\textbf{Question 416:} Which gate can be used to create an inverted output of an input signal in digital logic?

\begin{itemize}
    \item[\textbf{A.}] \textcolor{red}{\textbf{NOT gate}}
    \item[\textbf{B.}] OR gate
    \item[\textbf{C.}] AND gate
    \item[\textbf{D.}] XOR gate
\end{itemize}

\vspace{0.5cm}

\vspace{0.5cm}

\textbf{Question 420:} Which of the following Boolean expressions is equivalent to $F(A, B, C) = \Sigma(0, 1, 3, 4, 5)$

\begin{itemize}
    \item[\textbf{A.}] $F = AC + B'$
    \item[\textbf{B.}] \textcolor{red}{\textbf{$F = A'C + B'$}}
    \item[\textbf{C.}] $F = A'C' + B$
    \item[\textbf{D.}] $F = AC' + B$
\end{itemize}

\vspace{0.5cm}

\vspace{0.5cm}

\textbf{Question 424:} \underline{\hspace{1.5cm}}Simplify the following Boolean expression: $(A + B + C)(A + B' + C)(A' + B + C')(A' + B' + C')$

\begin{itemize}
    \item[\textbf{A.}] $A'BC + AB'C'$
    \item[\textbf{B.}] \textcolor{red}{\textbf{$A'C + AC'$}}
    \item[\textbf{C.}] $A'D + AD'$
    \item[\textbf{D.}] $AD' + BC'$
\end{itemize}

\vspace{0.5cm}

\vspace{0.5cm}

\textbf{Question 426:} Given the Karnaugh map with inputs $CD$ (rows) and $AB$ (columns) as shown below:

\begin{center}
\renewcommand{\arraystretch}{1.4}
\begin{tabular}{|c|c|c|c|c|}
\hline
$CD \backslash AB$ & \textbf{00} & \textbf{01} & \textbf{11} & \textbf{10} \\ \hline
\textbf{00} & 1 & 1 & & \\ \hline
\textbf{01} & 1 & 1 & & \\ \hline
\textbf{11} & & & 1 & 1 \\ \hline
\textbf{10} & & & 1 & 1 \\ \hline
\end{tabular}
\end{center}

Choose the corresponding canonical sum-of-products expression:

\begin{itemize}
    \item[\textbf{A.}] F(A,B,C,D) = A'B'C'D' + A'B'C'D + A'BC'D' + A'BC'D + AB'CD + ABC'D' + ABCD' + ABCD
    \item[\textbf{B.}] \textcolor{red}{\textbf{F(A,B,C,D) = A'B'C'D' + A'B'C'D + A'BC'D' + A'BC'D + AB'CD + AB'CD' + ABCD' + ABCD}}
    \item[\textbf{C.}] F(A,B,C,D) = A'B'C'D' + A'B'C'D + A'BC'D' + A'BC'D + AB'CD + AB'C'D' + ABCD' + ABCD
    \item[\textbf{D.}] F(A,B,C,D) = A'B'C'D' + A'B'C'D + A'BC'D' + A'BC'D + AB'CD + AB'CD' + ABC'D + ABCD
\end{itemize}

\vspace{0.5cm}

\textbf{Question 427:} The Boolean expression $F = A'C' + AC' + BC = B + C'$. true or false

\begin{itemize}
    \item[\textbf{A.}] \textcolor{red}{\textbf{true}}
    \item[\textbf{B.}] false
\end{itemize}

\vspace{0.5cm}

\vspace{0.5cm}

\textbf{Question 428:} Use Boolean algebra to find the most simplified SOP expression for F = ABD + CD + ACD + ABC + ABCD.

\begin{itemize}
    \item[\textbf{A.}] $F = AC + AD$
    \item[\textbf{B.}] $F = CD + AD$
    \item[\textbf{C.}] \textcolor{red}{\textbf{$F = ABD + ABC + CD$}}
    \item[\textbf{D.}] $F = BC + AB$
\end{itemize}

\vspace{0.5cm}

\vspace{0.5cm}

\textbf{Question 429:} Applying DeMorgan's theorem to the expression $((A + B + C)D)'$, we get \underline{\hspace{1.5cm}}.

\begin{itemize}
    \item[\textbf{A.}] $ABCD'$
    \item[\textbf{B.}] $(ABC)'D$
    \item[\textbf{C.}] \textcolor{red}{\textbf{$A'B'C' + D'$}}
    \item[\textbf{D.}] $A'B'C'D'$
\end{itemize}

\vspace{0.5cm}

\vspace{0.5cm}

\textbf{Question 430:} Which logic gate has an output of 1 only if both inputs are 0 or both inputs are 1?

\begin{itemize}
    \item[\textbf{A.}] NOT gate
    \item[\textbf{B.}] AND gate
    \item[\textbf{C.}] \textcolor{red}{\textbf{XNOR gate}}
    \item[\textbf{D.}] OR gate
\end{itemize}

\vspace{0.5cm}

\vspace{0.5cm}

\textbf{Question 431:} \underline{\hspace{1.5cm}}For the SOP expression $AB + B'C$, how many 0s are in the truth table's output column?

\begin{itemize}
    \item[\textbf{A.}] \textcolor{red}{\textbf{4}}
    \item[\textbf{B.}] zero
    \item[\textbf{C.}] 1
    \item[\textbf{D.}] 5
\end{itemize}

\vspace{0.5cm}

\vspace{0.5cm}

\textbf{Question 433:} \underline{\hspace{1.5cm}}Karnaugh map is a systematic way of reducing which type of expression?

\begin{itemize}
    \item[\textbf{A.}] those with overbars
    \item[\textbf{B.}] \textcolor{red}{\textbf{sum-of-products}}
    \item[\textbf{C.}] product-of-sums
    \item[\textbf{D.}] exclusive NOR
\end{itemize}

\vspace{0.5cm}

\vspace{0.5cm}

\textbf{Question 435:} Which logic gate is also known as a "buffer" and simply passes its input to its output?

\begin{itemize}
    \item[\textbf{A.}] XOR gate
    \item[\textbf{B.}] \textcolor{red}{\textbf{NOT gate}}
    \item[\textbf{C.}] AND gate
    \item[\textbf{D.}] OR gate
\end{itemize}

\vspace{0.5cm}

\vspace{0.5cm}

\textbf{Question 437:} Which of the following Boolean expressions is equivalent to F=AB'+ABC'+AB'C

\begin{itemize}
    \item[\textbf{A.}] $F = AC + AB'$
    \item[\textbf{B.}] $F = A'C + AB'$
    \item[\textbf{C.}] $F = AC + AB$
    \item[\textbf{D.}] \textcolor{red}{\textbf{$F = AC' + AB'$}}
\end{itemize}

\vspace{0.5cm}

\vspace{0.5cm}

\textbf{Question 438:} For the SOP expression AB'C+A'BC+ABC', how many 1s are in the truth table's output column?

\begin{itemize}
    \item[\textbf{A.}] 5
    \item[\textbf{B.}] \textcolor{red}{\textbf{3}}
    \item[\textbf{C.}] 1
    \item[\textbf{D.}] 2
\end{itemize}

\vspace{0.5cm}

\vspace{0.5cm}

\textbf{Question 440:} \underline{\hspace{1.5cm}}Simplify the following Boolean expression: $AB'C' + A'BC + ABC + ABC'$

\begin{itemize}
    \item[\textbf{A.}] $AC' + BC$
    \item[\textbf{B.}] $AC + BC'$
    \item[\textbf{C.}] \textcolor{red}{\textbf{$BC + AB + AB'C'$}}
    \item[\textbf{D.}] $A'B + B'C$
\end{itemize}

\vspace{0.5cm}

\vspace{0.5cm}

\textbf{Question 441:} Applying DeMorgan's theorem to the expression (ABC)', we get \underline{\hspace{1.5cm}}.

\begin{itemize}
    \item[\textbf{A.}] (A+B+C)'
    \item[\textbf{B.}] A+B'+CC'
    \item[\textbf{C.}] \textcolor{red}{\textbf{A' + B' + C'}}
    \item[\textbf{D.}] A(B+C)
\end{itemize}

\vspace{0.5cm}

\vspace{0.5cm}

\textbf{Question 442:} \underline{\hspace{1.5cm}}What is the output of the following Boolean expression: $(A' + B + C)\text{ AND }(A + B' + C')$?

\begin{itemize}
    \item[\textbf{A.}] \textcolor{red}{\textbf{1}}
    \item[\textbf{B.}] 0
    \item[\textbf{C.}] A XOR B
    \item[\textbf{D.}] Depends on the values of A, B, and C
\end{itemize}

\vspace{0.5cm}

\vspace{0.5cm}

\textbf{Question 445:} Write the Boolean expression for each of the logic circuits.

\begin{itemize}
    \item[\textbf{A.}] \textcolor{red}{\textbf{$F = A'C' + AB + AB'C$}}
    \item[\textbf{B.}] $F = AC + A'B + AB'C$
    \item[\textbf{C.}] $F = A'C' + A'B + AB'C$
    \item[\textbf{D.}] $F = AC' + A'B + AB'C$
\end{itemize}

\vspace{0.5cm}

\vspace{0.5cm}

\textbf{Question 448:} A Karnaugh map is a systematic way of reducing which type of expression?

\begin{itemize}
    \item[\textbf{A.}] product-of-sums
    \item[\textbf{B.}] those with overbars
    \item[\textbf{C.}] exclusive NOR
    \item[\textbf{D.}] \textcolor{red}{\textbf{sum-of-products}}
\end{itemize}

\vspace{0.5cm}

\vspace{0.5cm}

\textbf{Question 449:} \underline{\hspace{1.5cm}}Simplify the following Boolean expression using De Morgan's Law: $(AB')'$

\begin{itemize}
    \item[\textbf{A.}] $A + B$
    \item[\textbf{B.}] $A'B$
    \item[\textbf{C.}] $AB'$
    \item[\textbf{D.}] \textcolor{red}{\textbf{$A' + B$}}
\end{itemize}

\vspace{0.5cm}

\vspace{0.5cm}

\textbf{Question 452:} \underline{\hspace{1.5cm}}Which of the following is NOT a Boolean identity?

\begin{itemize}
    \item[\textbf{A.}] A AND (A OR B) = A
    \item[\textbf{B.}] A XOR B = (A AND B') OR (A' AND B)
    \item[\textbf{C.}] (A OR B)' = A' AND B'
    \item[\textbf{D.}] \textcolor{red}{\textbf{A + A' = 1}}
\end{itemize}

\vspace{0.5cm}

\vspace{0.5cm}

\textbf{Question 493:} Which of the following is a characteristic of sequential circuits compared to combinational circuits?

\begin{itemize}
    \item[\textbf{A.}] Do not use gates
    \item[\textbf{B.}] Output depends only on current inputs
    \item[\textbf{C.}] \textcolor{red}{\textbf{Output depends on current inputs and past inputs}}
    \item[\textbf{D.}] Simpler to design and implement
\end{itemize}

\vspace{0.5cm}

\vspace{0.5cm}

\textbf{Question 494:} Given the following Karnaugh map, what is the optimal expression in the sum of product format ($X'$ means NOT $X$)?

\begin{center}
\renewcommand{\arraystretch}{1.4}
\begin{tabular}{|c|c|c|c|c|}
\hline
$AB \backslash CD$ & \textbf{00} & \textbf{01} & \textbf{11} & \textbf{10} \\ \hline
\textbf{00} & & & & \\ \hline
\textbf{01} & 1 & 1 & 1 & 1 \\ \hline
\textbf{11} & & 1 & & 1 \\ \hline
\textbf{10} & & 1 & & 1 \\ \hline
\end{tabular}
\end{center}

\begin{itemize}
    \item[\textbf{A.}] \textcolor{red}{\textbf{A'B + AC'D + ACD'}}
    \item[\textbf{B.}] AB + AC' + ACD'
    \item[\textbf{C.}] AB' + AD + ACD'
    \item[\textbf{D.}] A'B + C'D + BCD'
\end{itemize}

\vspace{0.5cm}

\textbf{Question 521:} Given the following Karnaugh map, what is the optimal expression in the sum of product format ($X'$ means NOT $X$)?

\begin{center}
\renewcommand{\arraystretch}{1.4}
\begin{tabular}{|c|c|c|c|c|}
\hline
$A \backslash BC$ & \textbf{00} & \textbf{01} & \textbf{11} & \textbf{10} \\ \hline
\textbf{0} & 1 & 1 & 1 & 1 \\ \hline
\textbf{1} & 1 & & & 1 \\ \hline
\end{tabular}
\end{center}

\begin{itemize}
    \item[\textbf{A.}] \textcolor{red}{\textbf{$A' + C'$}}
    \item[\textbf{B.}] $A + B$
    \item[\textbf{C.}] $B + C'$
    \item[\textbf{D.}] $A + B'C' + BC'$
\end{itemize}

\vspace{0.5cm}

\textbf{Question 531:} Consider the truth table:

\begin{center}
\renewcommand{\arraystretch}{1.3}
\begin{tabular}{|c|c|c|}
\hline
\textbf{A} & \textbf{B} & \textbf{F} \\ \hline
0 & 0 & 1 \\ \hline
0 & 1 & 0 \\ \hline
1 & 0 & 0 \\ \hline
1 & 1 & 0 \\ \hline
\end{tabular}
\end{center}

Choose the correct Algebraic Function matching the given truth table:

\begin{itemize}
    \item[\textbf{A.}] \textcolor{red}{\textbf{F = NOT (A + B)}}
    \item[\textbf{B.}] F = A XOR B
    \item[\textbf{C.}] F = A AND NOT B
    \item[\textbf{D.}] F = NOT A AND B
\end{itemize}

\vspace{0.5cm}

\textbf{Question 534:} Given the following Karnaugh map, what is the optimal expression in the sum of product format ($X'$ means NOT $X$)?

\begin{center}
\renewcommand{\arraystretch}{1.4}
\begin{tabular}{|c|c|c|c|c|}
\hline
$A \backslash BC$ & \textbf{00} & \textbf{01} & \textbf{11} & \textbf{10} \\ \hline
\textbf{0} & & & 1 & 1 \\ \hline
\textbf{1} & 1 & & 1 & 1 \\ \hline
\end{tabular}
\end{center}

\begin{itemize}
    \item[\textbf{A.}] \textcolor{red}{\textbf{B + AC'}}
    \item[\textbf{B.}] B + AB'C'
    \item[\textbf{C.}] B + AC
    \item[\textbf{D.}] AB + B'C'
\end{itemize}

\vspace{0.5cm}

\textbf{Question 388:} The operation \underline{\hspace{1.5cm}} yields true (binary value 1) if and only if both of its operands are true.

\begin{itemize}
    \item[\textbf{A.}] OR
    \item[\textbf{B.}] \textcolor{red}{\textbf{AND}}
    \item[\textbf{C.}] XOR
    \item[\textbf{D.}] NAND
\end{itemize}

\vspace{0.5cm}

\vspace{0.5cm}

\textbf{Question 31:} The unary operation \underline{\hspace{1.5cm}} inverts the value of its operand.

\begin{itemize}
    \item[\textbf{A.}] XOR
    \item[\textbf{B.}] NAND
    \item[\textbf{C.}] \textcolor{red}{\textbf{NOT}}
    \item[\textbf{D.}] OR
\end{itemize}

\vspace{0.5cm}

\vspace{0.5cm}

\textbf{Question 47:} Counters can be designated as \underline{\hspace{1.5cm}}.

\begin{itemize}
    \item[\textbf{A.}] asynchronous
    \item[\textbf{B.}] neither asynchronous or synchronous
    \item[\textbf{C.}] synchronous
    \item[\textbf{D.}] \textcolor{red}{\textbf{both asynchronous and synchronous}}
\end{itemize}

\vspace{0.5cm}

\vspace{0.5cm}

\textbf{Question 227:} Counters can be designated as \underline{\hspace{1.5cm}}.

\begin{itemize}
    \item[\textbf{A.}] asynchronous
    \item[\textbf{B.}] synchronous
    \item[\textbf{C.}] \textcolor{red}{\textbf{both asynchronous and synchronous}}
    \item[\textbf{D.}] neither asynchronous or synchronous
\end{itemize}

\vspace{0.5cm}

\vspace{0.5cm}

\textbf{Question 327:} Consider an expression: $A + (B.C)$ What expression is equal to the given expression?

\begin{itemize}
    \item[\textbf{A.}] \textcolor{red}{\textbf{$(A + B).(A + C)$}}
    \item[\textbf{B.}] $(A + B).C$
    \item[\textbf{C.}] $A.(B + C)$
    \item[\textbf{D.}] $NOT(A.(B + C))$
\end{itemize}

\vspace{0.5cm}

\vspace{0.5cm}

\textbf{Question 443:} \underline{\hspace{1.5cm}}Which of the following is an important feature of the sum-of-products (SOP) form of expression?

\begin{itemize}
    \item[\textbf{A.}] No signal must pass through more than two gates, not including inverters.
    \item[\textbf{B.}] \textcolor{red}{\textbf{All logic circuits are reduced to nothing more than simple AND and OR gates.}}
    \item[\textbf{C.}] The delay times are greatly reduced over other forms.
    \item[\textbf{D.}] The maximum number of gates that any signal must pass through is reduced by a factor of two.
\end{itemize}

\vspace{0.5cm}

\vspace{0.5cm}

\textbf{Question 454:} Which output expression might indicate a product-of-sums circuit construction?

\begin{itemize}
    \item[\textbf{A.}] $X = CD + EG$
    \item[\textbf{B.}] $X = (CD)'(D + E + F)$
    \item[\textbf{C.}] $X = (AC + BD + EF)'$
    \item[\textbf{D.}] \textcolor{red}{\textbf{$X = (C + D)(E + G)$}}
\end{itemize}

\vspace{0.5cm}

\vspace{0.5cm}

\section{Chapter 13: Instruction Sets: Characteristics and Functions}
\textit{Total questions in chapter: 11}

\textbf{Question 523:} Question 32: Which of the following is not an operand type?

\begin{itemize}
    \item[\textbf{A.}] Addresses
    \item[\textbf{B.}] Logical data
    \item[\textbf{C.}] \textcolor{red}{\textbf{Value}}
    \item[\textbf{D.}] Characters
\end{itemize}

\vspace{0.5cm}

\vspace{0.5cm}

\textbf{Question 35:} \underline{\hspace{1.5cm}} is a principle by which two variables are independent of each other.

\begin{itemize}
    \item[\textbf{A.}] Autoindexing
    \item[\textbf{B.}] Opcode
    \item[\textbf{C.}] Completeness
    \item[\textbf{D.}] \textcolor{red}{\textbf{Orthogonality}}
\end{itemize}

\vspace{0.5cm}

\vspace{0.5cm}

\textbf{Question 96:} In reference to access time to a two-level memory, a \underline{\hspace{1.5cm}} occurs if an accessed word is found in the faster memory.

\begin{itemize}
    \item[\textbf{A.}] miss
    \item[\textbf{B.}] \textcolor{red}{\textbf{hit}}
    \item[\textbf{C.}] line
    \item[\textbf{D.}] tag
\end{itemize}

\vspace{0.5cm}

\vspace{0.5cm}

\textbf{Question 126:} Backward compatible means that the programs written for the older machines can be executed on the new machine.

\begin{itemize}
    \item[\textbf{A.}] \textcolor{red}{\textbf{True}}
    \item[\textbf{B.}] False
\end{itemize}

\vspace{0.5cm}

\vspace{0.5cm}

\textbf{Question 239:} The major cost in the life cycle of a system is hardware.

\begin{itemize}
    \item[\textbf{A.}] True
    \item[\textbf{B.}] \textcolor{red}{\textbf{False}}
\end{itemize}

\vspace{0.5cm}

\vspace{0.5cm}

\textbf{Question 246:} The \underline{\hspace{1.5cm}} defines the repertoire of machine language instructions that a computer can follow.

\begin{itemize}
    \item[\textbf{A.}] ABI
    \item[\textbf{B.}] API
    \item[\textbf{C.}] HLL
    \item[\textbf{D.}] \textcolor{red}{\textbf{ISA}}
\end{itemize}

\vspace{0.5cm}

\vspace{0.5cm}

\textbf{Question 285:} Why is it essential to use symbolic representation of machine instructions?

\begin{itemize}
    \item[\textbf{A.}] \textcolor{red}{\textbf{It makes machine instructions more human-readable and understandable}}
    \item[\textbf{B.}] It reduces the overall complexity of computer systems and user programs
    \item[\textbf{C.}] It minimizes the need for memory storage for the user programs
    \item[\textbf{D.}] It enables fastest execution of high level language instructions
\end{itemize}

\vspace{0.5cm}

\vspace{0.5cm}

\textbf{Question 287:} What is a branch instrtuction?

\begin{itemize}
    \item[\textbf{A.}] The intructions that are used to divide a program into multiple subprograms
    \item[\textbf{B.}] \textcolor{red}{\textbf{The intructions that have as one of its operands the address of the next instruction to be executed}}
    \item[\textbf{C.}] The intructions that are used to pause the program
    \item[\textbf{D.}] The intructions that are used to return to the beginning of the program
\end{itemize}

\vspace{0.5cm}

\vspace{0.5cm}

\textbf{Question 338:} What are the most important general categories of data that machine instructions operate on?

\begin{itemize}
    \item[\textbf{A.}] \textcolor{red}{\textbf{Addresses, numbers, characters, and logical data}}
    \item[\textbf{B.}] Text, images, and audio
    \item[\textbf{C.}] Variables, functions, and arrays
    \item[\textbf{D.}] Instructions, control signals, and registers
\end{itemize}

\vspace{0.5cm}

\vspace{0.5cm}

\textbf{Question 471:} What type of instruction is used for converting data formats?

\begin{itemize}
    \item[\textbf{A.}] Transfer
    \item[\textbf{B.}] Arithmetic
    \item[\textbf{C.}] Control
    \item[\textbf{D.}] \textcolor{red}{\textbf{Conversion}}
\end{itemize}

\vspace{0.5cm}

\vspace{0.5cm}

\textbf{Question 496:} With respect to elements of a machine instruction, what does the instruction code mean?

\begin{itemize}
    \item[\textbf{A.}] \textcolor{red}{\textbf{It specifies the operation to be performed.}}
    \item[\textbf{B.}] It specifies data which will be processed.
    \item[\textbf{C.}] It specifies memory access.
    \item[\textbf{D.}] It specifies an IO device.
\end{itemize}

\vspace{0.5cm}

\vspace{0.5cm}

\section{Chapter 14: Instruction Sets: Addressing Modes and Formats}
\textit{Total questions in chapter: 34}

\textbf{Question 32:} Three of the most common uses of stack addressing are relative addressing, base-register addressing, and indexing.

\begin{itemize}
    \item[\textbf{A.}] \textcolor{red}{\textbf{False}}
    \item[\textbf{B.}] True
\end{itemize}

\vspace{0.5cm}

\vspace{0.5cm}

\textbf{Question 68:} The method of calculating the EA is the same for both base-register addressing and indexing

\begin{itemize}
    \item[\textbf{A.}] False
    \item[\textbf{B.}] \textcolor{red}{\textbf{True}}
\end{itemize}

\vspace{0.5cm}

\vspace{0.5cm}

\textbf{Question 183:} The advantages of \underline{\hspace{1.5cm}} addressing are that only a small address field is needed in the instruction and no time-consuming memory references are required.

\begin{itemize}
    \item[\textbf{A.}] indirect
    \item[\textbf{B.}] displacement
    \item[\textbf{C.}] direct
    \item[\textbf{D.}] \textcolor{red}{\textbf{register}}
\end{itemize}

\vspace{0.5cm}

\vspace{0.5cm}

\textbf{Question 511:} In the context of addressing modes with both indirect addressing and indexing, what is postindexing?

\begin{itemize}
    \item[\textbf{A.}] \textcolor{red}{\textbf{Indexing is performed after the indirection}}
    \item[\textbf{B.}] Indexing is performed before the indirection
    \item[\textbf{C.}] Both indexing and indirection are avoided
    \item[\textbf{D.}] Indexing and indirection occur simultaneously
\end{itemize}

\vspace{0.5cm}

\vspace{0.5cm}

\textbf{Question 2:} \underline{\hspace{1.5cm}} has the advantage of flexibility, but the disadvantage of complexity.

\begin{itemize}
    \item[\textbf{A.}] Stack addressing
    \item[\textbf{B.}] \textcolor{red}{\textbf{Displacement addressing}}
    \item[\textbf{C.}] Direct addressing
    \item[\textbf{D.}] Register addressing
\end{itemize}

\vspace{0.5cm}

\vspace{0.5cm}

\textbf{Question 82:} Indexing performed after the indirection is \underline{\hspace{1.5cm}} .

\begin{itemize}
    \item[\textbf{A.}] relative addressing
    \item[\textbf{B.}] autoindexing
    \item[\textbf{C.}] \textcolor{red}{\textbf{postindexing}}
    \item[\textbf{D.}] preindexing
\end{itemize}

\vspace{0.5cm}

\vspace{0.5cm}

\textbf{Question 128:} \underline{\hspace{1.5cm}} has the advantage of large address space, however it has the disadvantage of multiple memory references.

\begin{itemize}
    \item[\textbf{A.}] \textcolor{red}{\textbf{Indirect addressing}}
    \item[\textbf{B.}] Direct addressing
    \item[\textbf{C.}] Immediate addressing
    \item[\textbf{D.}] Stack addressing
\end{itemize}

\vspace{0.5cm}

\vspace{0.5cm}

\textbf{Question 141:} The principal advantage of \underline{\hspace{1.5cm}} addressing is that it is a very simple form of addressing.

\begin{itemize}
    \item[\textbf{A.}] displacement
    \item[\textbf{B.}] register
    \item[\textbf{C.}] stack
    \item[\textbf{D.}] \textcolor{red}{\textbf{direct}}
\end{itemize}

\vspace{0.5cm}

\vspace{0.5cm}

\textbf{Question 149:} The advantage of \underline{\hspace{1.5cm}} is that no memory reference other than the instruction fetch is required to obtain the operand.

\begin{itemize}
    \item[\textbf{A.}] direct addressing
    \item[\textbf{B.}] \textcolor{red}{\textbf{immediate addressing}}
    \item[\textbf{C.}] register addressing
    \item[\textbf{D.}] stack addressing
\end{itemize}

\vspace{0.5cm}

\vspace{0.5cm}

\textbf{Question 152:} The only form of addressing for branch instructions is \underline{\hspace{1.5cm}} addressing.

\begin{itemize}
    \item[\textbf{A.}] register
    \item[\textbf{B.}] relative
    \item[\textbf{C.}] base
    \item[\textbf{D.}] \textcolor{red}{\textbf{immediate}}
\end{itemize}

\vspace{0.5cm}

\vspace{0.5cm}

\textbf{Question 288:} The effective address of \underline{\hspace{1.5cm}} is
$$EA = A + (R)$$
$$(R) \leftarrow (R) + 1$$

\begin{itemize}
    \item[\textbf{A.}] relative addressing
    \item[\textbf{B.}] \textcolor{red}{\textbf{autoindexing}}
    \item[\textbf{C.}] postindexing
    \item[\textbf{D.}] preindexing
\end{itemize}

\vspace{0.5cm}

\vspace{0.5cm}

\textbf{Question 390:} Which of the following statement is correct about addressing modes?

\begin{itemize}
    \item[\textbf{A.}] \textcolor{red}{\textbf{They define how the operands of an instruction are specified, including immediate, register, direct, and indirect addressing modes}}
    \item[\textbf{B.}] Addressing modes are irrelevant in computer architecture, and all instructions only operate on values stored in registers
    \item[\textbf{C.}] Addressing modes are limited to only immediate and direct modes; register and indirect addressing modes are not used in modern computer systems
    \item[\textbf{D.}] All instructions in computer architecture use indirect addressing modes, making it the only relevant mode for operand specification
\end{itemize}

\vspace{0.5cm}

\vspace{0.5cm}

\textbf{Question 54:} With direct addressing, the length of the address field is usually less than the word length, thus limiting the address range.

\begin{itemize}
    \item[\textbf{A.}] \textcolor{red}{\textbf{True}}
    \item[\textbf{B.}] False
\end{itemize}

\vspace{0.5cm}

\vspace{0.5cm}

\textbf{Question 85:} Register indirect addressing uses the same number of memory references as indirect addressing.

\begin{itemize}
    \item[\textbf{A.}] True
    \item[\textbf{B.}] \textcolor{red}{\textbf{False}}
\end{itemize}

\vspace{0.5cm}

\vspace{0.5cm}

\textbf{Question 94:} The value of the mode field determines which addressing mode is to be used.

\begin{itemize}
    \item[\textbf{A.}] \textcolor{red}{\textbf{True}}
    \item[\textbf{B.}] False
\end{itemize}

\vspace{0.5cm}

\vspace{0.5cm}

\textbf{Question 154:} The disadvantage of immediate addressing is that the size of the number is restricted to the size of the address field.

\begin{itemize}
    \item[\textbf{A.}] \textcolor{red}{\textbf{True}}
    \item[\textbf{B.}] False
\end{itemize}

\vspace{0.5cm}

\vspace{0.5cm}

\textbf{Question 241:} Register addressing is similar to direct addressing with the only difference being that the address field refers to a register rather than a main memory address.

\begin{itemize}
    \item[\textbf{A.}] \textcolor{red}{\textbf{True}}
    \item[\textbf{B.}] False
\end{itemize}

\vspace{0.5cm}

\vspace{0.5cm}

\textbf{Question 247:} Almost all RISC instructions use simple register addressing.

\begin{itemize}
    \item[\textbf{A.}] \textcolor{red}{\textbf{True}}
    \item[\textbf{B.}] False
\end{itemize}

\vspace{0.5cm}

\vspace{0.5cm}

\textbf{Question 340:} Which addressing mode allows direct specification of the memory address within the instruction?

\begin{itemize}
    \item[\textbf{A.}] \textcolor{red}{\textbf{Direct}}
    \item[\textbf{B.}] Indirect
    \item[\textbf{C.}] Register Indirect
    \item[\textbf{D.}] Displacement
\end{itemize}

\vspace{0.5cm}

\vspace{0.5cm}

\textbf{Question 372:} The effective address of Register indirect addressing mode is \underline{\hspace{1.5cm}}.

\begin{itemize}
    \item[\textbf{A.}] $EA = R$
    \item[\textbf{B.}] \textcolor{red}{\textbf{$EA = (R)$}}
    \item[\textbf{C.}] $EA = (R) + A$
    \item[\textbf{D.}] $EA = (R) + (A)$
\end{itemize}

\vspace{0.5cm}

\vspace{0.5cm}

\textbf{Question 497:} What is the primary advantage of indirect addressing mode?

\begin{itemize}
    \item[\textbf{A.}] Faster access to operands
    \item[\textbf{B.}] Ability to access a larger data set
    \item[\textbf{C.}] \textcolor{red}{\textbf{Flexibility in data location}}
    \item[\textbf{D.}] Reduced instruction size
\end{itemize}

\vspace{0.5cm}

\vspace{0.5cm}

\textbf{Question 522:} Which type of instruction format is more flexible but can lead to inefficient use of memory?

\begin{itemize}
    \item[\textbf{A.}] Fixed-Length Instruction Format
    \item[\textbf{B.}] \textcolor{red}{\textbf{Variable-Length Instruction Format}}
    \item[\textbf{C.}] Compact Instruction Format
    \item[\textbf{D.}] Standardized Instruction Format
\end{itemize}

\vspace{0.5cm}

\vspace{0.5cm}

\textbf{Question 524:} Indirect addressing, the length of the address field is usually greater than the word length, because unlimiting the address range.

\begin{itemize}
    \item[\textbf{A.}] True
    \item[\textbf{B.}] \textcolor{red}{\textbf{False}}
\end{itemize}

\vspace{0.5cm}

\vspace{0.5cm}

\textbf{Question 345:} When considering the number of pipeline stages, what trade-offs must be made in computer architecture?

\begin{itemize}
    \item[\textbf{A.}] \textcolor{red}{\textbf{Trade-offs between potential speedup and increased cost and delays}}
    \item[\textbf{B.}] Trade-offs between software and hardware
    \item[\textbf{C.}] Trade-offs between speed and efficiency
    \item[\textbf{D.}] Trade-offs between branching and executing instructions with conditions
\end{itemize}

\vspace{0.5cm}

\vspace{0.5cm}

\textbf{Question 344:} The PC, IR, MAR, MBR registers belong to which of the following groups?

\begin{itemize}
    \item[\textbf{A.}] \textcolor{red}{\textbf{Control and Status Registers}}
    \item[\textbf{B.}] User-Visible Registers
    \item[\textbf{C.}] General Registers
    \item[\textbf{D.}] Handle Registers
\end{itemize}

\vspace{0.5cm}

\vspace{0.5cm}

\textbf{Question 74:} For \underline{\hspace{1.5cm}}, the address field references a main memory address and the referenced register contains a positive displacement from that address.

\begin{itemize}
    \item[\textbf{A.}] \textcolor{red}{\textbf{indexing}}
    \item[\textbf{B.}] base-register addressing
    \item[\textbf{C.}] relative addressing
    \item[\textbf{D.}] all of the above
\end{itemize}

\vspace{0.5cm}

\vspace{0.5cm}

\textbf{Question 90:} Typically an instruction set will include both preindexing and postindexing.

\begin{itemize}
    \item[\textbf{A.}] True
    \item[\textbf{B.}] \textcolor{red}{\textbf{False}}
\end{itemize}

\vspace{0.5cm}

\vspace{0.5cm}

\textbf{Question 223:} Assume a stack-oriented processor that includes the stack operations PUSH and POP. Arithmetic operations automatically involve the top one or two stack elements. Begin with an empty stack. What stack elements remain after the following instructions are executed? PUSH 3 ; PUSH 4 ; ADD ; PUSH 5 ; MUL.

\begin{itemize}
    \item[\textbf{A.}] \textcolor{red}{\textbf{35}}
    \item[\textbf{B.}] 60
    \item[\textbf{C.}] 34
    \item[\textbf{D.}] 39
    \item[\textbf{E.}] 345
\end{itemize}

\vspace{0.5cm}

\vspace{0.5cm}

\textbf{Question 28:} Assume an instruction set that uses a fixed 14-bit instruction length. Operand specifiers are 6 bits length. What is the maximum number of one-operand instructions that can be supported?

\begin{itemize}
    \item[\textbf{A.}] 256
    \item[\textbf{B.}] 128
    \item[\textbf{C.}] 32
    \item[\textbf{D.}] 512
    \item[\textbf{E.}] \textcolor{red}{\textbf{64}}
\end{itemize}

\vspace{0.5cm}

\vspace{0.5cm}

\textbf{Question 100:} Which of the following interrelated factors go into determining the use of the addressing bits?

\begin{itemize}
    \item[\textbf{A.}] number of operands
    \item[\textbf{B.}] number of register sets
    \item[\textbf{C.}] address range
    \item[\textbf{D.}] \textcolor{red}{\textbf{all of the above}}
\end{itemize}

\vspace{0.5cm}

\vspace{0.5cm}

\textbf{Question 229:} The advantage of \underline{\hspace{1.5cm}} is that no memory reference other than the instruction fetch is required to obtain the operand.

\begin{itemize}
    \item[\textbf{A.}] direct addressing
    \item[\textbf{B.}] \textcolor{red}{\textbf{immediate addressing}}
    \item[\textbf{C.}] register addressing
    \item[\textbf{D.}] stack addressing
\end{itemize}

\vspace{0.5cm}

\vspace{0.5cm}

\textbf{Question 472:} The effective address of \underline{\hspace{1.5cm}} is EA=top of stack

\begin{itemize}
    \item[\textbf{A.}] Immediate
    \item[\textbf{B.}] Direct
    \item[\textbf{C.}] Indirect
    \item[\textbf{D.}] Displacement
    \item[\textbf{E.}] \textcolor{red}{\textbf{Implicit}}
\end{itemize}

\vspace{0.5cm}

\vspace{0.5cm}

\textbf{Question 495:} In a conditional branch instruction, what happens if the specified condition is met?

\begin{itemize}
    \item[\textbf{A.}] \textcolor{red}{\textbf{The program counter is updated to equal the address specified in the operand.}}
    \item[\textbf{B.}] The program counter increments as usual.
    \item[\textbf{C.}] The instruction is skipped.
    \item[\textbf{D.}] The program terminates.
\end{itemize}

\vspace{0.5cm}

\vspace{0.5cm}

\textbf{Question 525:} Assume an instruction set that uses a fixed 16-bit instruction length. Operand specifiers are 6 bits in length. What is the maximum number of one-operand instructions that can be supported?

\begin{itemize}
    \item[\textbf{A.}] 256
    \item[\textbf{B.}] 512
    \item[\textbf{C.}] \textcolor{red}{\textbf{1024}}
    \item[\textbf{D.}] 2048
\end{itemize}

\vspace{0.5cm}

\vspace{0.5cm}

\section{Chapter 15: Assembly Language and Related Topics}
\textit{Total questions in chapter: 4}

\textbf{Question 342:} In MASM32, which command is incorrect?

\begin{itemize}
    \item[\textbf{A.}] ADD EAX, a
    \item[\textbf{B.}] ADD EAX, EBX
    \item[\textbf{C.}] ADD a, EAX
    \item[\textbf{D.}] \textcolor{red}{\textbf{ADD a, b}}
\end{itemize}

\vspace{0.5cm}

\vspace{0.5cm}

\textbf{Question 419:} What does the x86 assembly instruction "jmp label" typically do?

\begin{itemize}
    \item[\textbf{A.}] Jumps to the memory address stored in the label
    \item[\textbf{B.}] Jumps to the next instruction
    \item[\textbf{C.}] Jumps to the label if a specific condition is met
    \item[\textbf{D.}] \textcolor{red}{\textbf{Moves the label's address to the EIP register}}
\end{itemize}

\vspace{0.5cm}

\vspace{0.5cm}

\textbf{Question 526:} Which instruction is used by MASM32 to return from a procedure?

\begin{itemize}
    \item[\textbf{A.}] \textcolor{red}{\textbf{RET}}
    \item[\textbf{B.}] EXIT
    \item[\textbf{C.}] END
    \item[\textbf{D.}] RES
\end{itemize}

\vspace{0.5cm}

\vspace{0.5cm}

\textbf{Question 347:} What is one of the advantages of using a register file in computer architecture?

\begin{itemize}
    \item[\textbf{A.}] \textcolor{red}{\textbf{Reduction in memory accesses, saving time}}
    \item[\textbf{B.}] More efficient use of space due to dynamic adaptation
    \item[\textbf{C.}] Efficient handling of both local and global variables
    \item[\textbf{D.}] Easier management of cache residency
\end{itemize}

\vspace{0.5cm}

\vspace{0.5cm}

\section{Chapter 16: Processor Structure and Function}
\textit{Total questions in chapter: 53}

\textbf{Question 37:} The register file employs much shorter addresses than addresses for cache and memory.

\begin{itemize}
    \item[\textbf{A.}] \textcolor{red}{\textbf{True}}
    \item[\textbf{B.}] False
\end{itemize}

\vspace{0.5cm}

\vspace{0.5cm}

\textbf{Question 71:} The \underline{\hspace{1.5cm}} is a small cache memory associated with the instruction fetch stage of the pipeline.

\begin{itemize}
    \item[\textbf{A.}] dynamic branch
    \item[\textbf{B.}] loop table
    \item[\textbf{C.}] \textcolor{red}{\textbf{branch history table}}
    \item[\textbf{D.}] flag
\end{itemize}

\vspace{0.5cm}

\vspace{0.5cm}

\textbf{Question 112:} The \underline{\hspace{1.5cm}} stage includes ALU operations, cache access, and register update.

\begin{itemize}
    \item[\textbf{A.}] decode
    \item[\textbf{B.}] \textcolor{red}{\textbf{execute}}
    \item[\textbf{C.}] fetch
    \item[\textbf{D.}] write back
\end{itemize}

\vspace{0.5cm}

\vspace{0.5cm}

\textbf{Question 406:} In the context of instruction execution, how is a product on an assembly line conceptually similar to an instruction in a pipeline?

\begin{itemize}
    \item[\textbf{A.}] \textcolor{red}{\textbf{Both undergo multiple stages of production}}
    \item[\textbf{B.}] Both are executed in a single clock cycle
    \item[\textbf{C.}] Both follow a linear sequence of tasks
    \item[\textbf{D.}] Both are processed by the control unit
\end{itemize}

\vspace{0.5cm}

\vspace{0.5cm}

\textbf{Question 45:} \underline{\hspace{1.5cm}} registers may be used only to hold data and cannot be employed in the calculation of an address.

\begin{itemize}
    \item[\textbf{A.}] General purpose
    \item[\textbf{B.}] Condition code
    \item[\textbf{C.}] \textcolor{red}{\textbf{Data}}
    \item[\textbf{D.}] Address
\end{itemize}

\vspace{0.5cm}

\vspace{0.5cm}

\textbf{Question 150:} The processor needs to store instructions and data temporarily while an instruction is being executed.

\begin{itemize}
    \item[\textbf{A.}] \textcolor{red}{\textbf{True}}
    \item[\textbf{B.}] False
\end{itemize}

\vspace{0.5cm}

\vspace{0.5cm}

\textbf{Question 291:} How do data registers and address registers differ in some computer systems?

\begin{itemize}
    \item[\textbf{A.}] \textcolor{red}{\textbf{Address registers can be employed in calculating operand addresses, while data registers hold data.}}
    \item[\textbf{B.}] Data registers are only used for stack-related operations
    \item[\textbf{C.}] Data registers are used for indexed addressing, while address registers are used for data storage
    \item[\textbf{D.}] Address registers are reserved for segmented addressing, while data registers are general-purpose
\end{itemize}

\vspace{0.5cm}

\vspace{0.5cm}

\textbf{Question 375:} Control and status registers are used by which entities to control the operation of the processor?

\begin{itemize}
    \item[\textbf{A.}] \textcolor{red}{\textbf{Privileged, operating system programs}}
    \item[\textbf{B.}] Machine or assembly language programmers
    \item[\textbf{C.}] External I/O devices
    \item[\textbf{D.}] Main memory modules
\end{itemize}

\vspace{0.5cm}

\vspace{0.5cm}

\textbf{Question 455:} Which portion of a computer consists of an ALU, a control unit, and registers, uses for fetching and executing instructions?

\begin{itemize}
    \item[\textbf{A.}] Main memory
    \item[\textbf{B.}] \textcolor{red}{\textbf{CPU}}
    \item[\textbf{C.}] I/O
    \item[\textbf{D.}] System interconnection
\end{itemize}

\vspace{0.5cm}

\vspace{0.5cm}

\textbf{Question 22:} Pipelining is a means of introducing parallelism into the essentially sequential nature of a machine-instruction program.

\begin{itemize}
    \item[\textbf{A.}] \textcolor{red}{\textbf{True}}
    \item[\textbf{B.}] False
\end{itemize}

\vspace{0.5cm}

\vspace{0.5cm}

\textbf{Question 108:} One of the major problems in designing an instruction pipeline is assuring a steady flow of in initial stages of the pipeline.

\begin{itemize}
    \item[\textbf{A.}] \textcolor{red}{\textbf{True}}
    \item[\textbf{B.}] False
\end{itemize}

\vspace{0.5cm}

\vspace{0.5cm}

\textbf{Question 135:} The register file is on the same chip as the ALU and control unit.

\begin{itemize}
    \item[\textbf{A.}] \textcolor{red}{\textbf{True}}
    \item[\textbf{B.}] False
\end{itemize}

\vspace{0.5cm}

\vspace{0.5cm}

\textbf{Question 171:} The cycle time of an instruction pipeline is the time needed to advance a set of instructions one stage through the pipeline.

\begin{itemize}
    \item[\textbf{A.}] \textcolor{red}{\textbf{True}}
    \item[\textbf{B.}] False
\end{itemize}

\vspace{0.5cm}

\vspace{0.5cm}

\textbf{Question 250:} A control hazard occurs when two or more instructions that are already in the pipeline need the same resource.

\begin{itemize}
    \item[\textbf{A.}] True
    \item[\textbf{B.}] \textcolor{red}{\textbf{False}}
\end{itemize}

\vspace{0.5cm}

\vspace{0.5cm}

\textbf{Question 263:} In the CPU, what is the functionality of the control unit?

\begin{itemize}
    \item[\textbf{A.}] To decode program instructions
    \item[\textbf{B.}] \textcolor{red}{\textbf{To controls the sequence of operations}}
    \item[\textbf{C.}] To store program instructions
    \item[\textbf{D.}] To transfer data to primary storage
\end{itemize}

\vspace{0.5cm}

\vspace{0.5cm}

\textbf{Question 290:} What is the role of the control unit in a processor?

\begin{itemize}
    \item[\textbf{A.}] The control unit's primary role is to perform arithmetic and logical operations within the processor, orchestrating the manipulation of data
    \item[\textbf{B.}] The control unit only manages the flow of data between the CPU and external devices and does not play a significant role in executing instructions
    \item[\textbf{C.}] The control unit is solely responsible for managing the flow of instructions from secondary storage to RAM and does not have a role in the internal operation of the CPU
    \item[\textbf{D.}] \textcolor{red}{\textbf{The control unit in a processor directs and coordinates the execution of instructions, interpreting and managing the flow of operations within the CPU}}
\end{itemize}

\vspace{0.5cm}

\vspace{0.5cm}

\textbf{Question 355:} What is the most important function of the control unit (CU)?

\begin{itemize}
    \item[\textbf{A.}] It manages the order of running instructions.
    \item[\textbf{B.}] It will read and process data from main memory.
    \item[\textbf{C.}] \textcolor{red}{\textbf{It directs the operation of the other CPU components.}}
    \item[\textbf{D.}] It will read instructions from main memory then decoding them.
\end{itemize}

\vspace{0.5cm}

\vspace{0.5cm}

\textbf{Question 376:} Which of the following statement is correct in the context of Instruction Pipelining?

\begin{itemize}
    \item[\textbf{A.}] Instruction Pipelining reduces the efficiency of instruction execution by introducing delays and dependencies between instructions
    \item[\textbf{B.}] Instruction Pipelining is only effective for specific types of instructions and has no impact on the overall efficiency of instruction execution
    \item[\textbf{C.}] Instruction Pipelining enhances efficiency by enabling simultaneous execution of multiple instructions in different stages, boosting overall throughput
    \item[\textbf{D.}] \textcolor{red}{\textbf{Instruction Pipelining improves efficiency by processing multiple instructions simultaneously, reducing execution time. However, it can face challenges like hazards, introducing delays and impacting overall performance}}
\end{itemize}

\vspace{0.5cm}

\vspace{0.5cm}

\textbf{Question 410:} In the context of the basic instruction cycle, when the fifth instruction is executing, which of the following is the correct statement?

\begin{itemize}
    \item[\textbf{A.}] \textcolor{red}{\textbf{The sixth instruction is fetched.}}
    \item[\textbf{B.}] The fifth instruction is fetched.
    \item[\textbf{C.}] The fourth instruction is fetched.
    \item[\textbf{D.}] All program's instructions are fetched.
\end{itemize}

\vspace{0.5cm}

\vspace{0.5cm}

\textbf{Question 500:} What is a drawback of Instruction Pipelining?

\begin{itemize}
    \item[\textbf{A.}] It increases the instruction throughput
    \item[\textbf{B.}] It allows more CPU throughput than a multicycle computer at a given ck rate
    \item[\textbf{C.}] \textcolor{red}{\textbf{It may increase latency due to the added overhead of the pipelining process itself}}
    \item[\textbf{D.}] It decreases the complexity of the computer system
\end{itemize}

\vspace{0.5cm}

\vspace{0.5cm}

\textbf{Question 77:} \underline{\hspace{1.5cm}} determines the control and pipeline organization.

\begin{itemize}
    \item[\textbf{A.}] Calculation
    \item[\textbf{B.}] \textcolor{red}{\textbf{Execution sequencing}}
    \item[\textbf{C.}] Operations performed
    \item[\textbf{D.}] Operands used
\end{itemize}

\vspace{0.5cm}

\vspace{0.5cm}

\textbf{Question 359:} \underline{\hspace{1.5cm}} interprets the instructions in memory and causes them to be executed.

\begin{itemize}
    \item[\textbf{A.}] Registers
    \item[\textbf{B.}] CPU interconnection
    \item[\textbf{C.}] Arithmetic and Logic Unit (ALU)
    \item[\textbf{D.}] I/O Modules
    \item[\textbf{E.}] \textcolor{red}{\textbf{Control Unit (CU)}}
\end{itemize}

\vspace{0.5cm}

\vspace{0.5cm}

\textbf{Question 473:} Steps for executing a machine instruction are concerned, choose the correct statement which specifies the role of the fetching instruction step.

\begin{itemize}
    \item[\textbf{A.}] \textcolor{red}{\textbf{The processor reads an instruction from memory (cache, main memory).}}
    \item[\textbf{B.}] The instruction is decoded to determine what action is required.
    \item[\textbf{C.}] The execution of an instruction may require performing some arithmetic or logical operation on data.
    \item[\textbf{D.}] The results of an execution may require writing data to memory or an I/O module.
\end{itemize}

\vspace{0.5cm}

\vspace{0.5cm}

\textbf{Question 499:} Steps for executing a machine instruction are concerned, choose the most suitable statement for the role of the interpret instruction step.

\begin{itemize}
    \item[\textbf{A.}] \textcolor{red}{\textbf{The instruction is decoded to determine what action is required.}}
    \item[\textbf{B.}] The processor reads an instruction from memory (cache, main memory).
    \item[\textbf{C.}] The execution of an instruction may require performing some arithmetic or logical operation on data.
    \item[\textbf{D.}] The results of an execution may require writing data to memory or an I/O module.
\end{itemize}

\vspace{0.5cm}

\vspace{0.5cm}

\textbf{Question 156:} The allocation of control information between registers and memory are not considered to be issue.

\begin{itemize}
    \item[\textbf{A.}] True
    \item[\textbf{B.}] \textcolor{red}{\textbf{False}}
\end{itemize}

\vspace{0.5cm}

\vspace{0.5cm}

\textbf{Question 25:} Which one of the following CPU registers holds the address of the storage location being accessed?

\begin{itemize}
    \item[\textbf{A.}] \textcolor{red}{\textbf{MAR (Memory address register)}}
    \item[\textbf{B.}] MBR (Memory Buffer Register)
    \item[\textbf{C.}] AC (Accumulator)
    \item[\textbf{D.}] IR (Instruction Register)
\end{itemize}

\vspace{0.5cm}

\vspace{0.5cm}

\textbf{Question 44:} The predict-never-taken approach is the most popular of all the branch prediction methods.

\begin{itemize}
    \item[\textbf{A.}] \textcolor{red}{\textbf{True}}
    \item[\textbf{B.}] False
\end{itemize}

\vspace{0.5cm}

\vspace{0.5cm}

\textbf{Question 53:} What is stored in the Stack Pointer?

\begin{itemize}
    \item[\textbf{A.}] Address of next instruction
    \item[\textbf{B.}] Stack data values
    \item[\textbf{C.}] Addressing method
    \item[\textbf{D.}] Operations
    \item[\textbf{E.}] \textcolor{red}{\textbf{Address of top item}}
\end{itemize}

\vspace{0.5cm}

\vspace{0.5cm}

\textbf{Question 62:} The \underline{\hspace{1.5cm}} determines the opcode and the operand specifiers.

\begin{itemize}
    \item[\textbf{A.}] \textcolor{red}{\textbf{decode instruction}}
    \item[\textbf{B.}] fetch operands
    \item[\textbf{C.}] calculate operands
    \item[\textbf{D.}] execute instruction
\end{itemize}

\vspace{0.5cm}

\vspace{0.5cm}

\textbf{Question 66:} \underline{\hspace{1.5cm}} are a set of storage locations.

\begin{itemize}
    \item[\textbf{A.}] \textcolor{red}{\textbf{Registers}}
    \item[\textbf{B.}] Control units
    \item[\textbf{C.}] PSWs
    \item[\textbf{D.}] Processors
\end{itemize}

\vspace{0.5cm}

\vspace{0.5cm}

\textbf{Question 67:} The \underline{\hspace{1.5cm}} holds the address of the next instruction to be fetched.

\begin{itemize}
    \item[\textbf{A.}] MBR
    \item[\textbf{B.}] IR
    \item[\textbf{C.}] \textcolor{red}{\textbf{PC}}
    \item[\textbf{D.}] MAR
\end{itemize}

\vspace{0.5cm}

\vspace{0.5cm}

\textbf{Question 75:} \underline{\hspace{1.5cm}} are a set of storage locations.

\begin{itemize}
    \item[\textbf{A.}] Processors
    \item[\textbf{B.}] PSWs
    \item[\textbf{C.}] \textcolor{red}{\textbf{Registers}}
    \item[\textbf{D.}] Control units
\end{itemize}

\vspace{0.5cm}

\vspace{0.5cm}

\textbf{Question 78:} The \underline{\hspace{1.5cm}} contains the address of an instruction to be fetched.

\begin{itemize}
    \item[\textbf{A.}] instruction register
    \item[\textbf{B.}] memory address register
    \item[\textbf{C.}] memory buffer register
    \item[\textbf{D.}] \textcolor{red}{\textbf{program counter}}
\end{itemize}

\vspace{0.5cm}

\vspace{0.5cm}

\textbf{Question 121:} The disadvantage of the software poll is that it is time consuming.

\begin{itemize}
    \item[\textbf{A.}] \textcolor{red}{\textbf{True}}
    \item[\textbf{B.}] False
\end{itemize}

\vspace{0.5cm}

\vspace{0.5cm}

\textbf{Question 139:} What is stored in the Stack Pointer?

\begin{itemize}
    \item[\textbf{A.}] Operations
    \item[\textbf{B.}] Addressing method
    \item[\textbf{C.}] Stack data values
    \item[\textbf{D.}] \textcolor{red}{\textbf{Address of top item}}
    \item[\textbf{E.}] Address of next instruction
\end{itemize}

\vspace{0.5cm}

\vspace{0.5cm}

\textbf{Question 142:} \underline{\hspace{1.5cm}} provide storage internal to the CPU.

\begin{itemize}
    \item[\textbf{A.}] Control units
    \item[\textbf{B.}] ALUs
    \item[\textbf{C.}] Main memory
    \item[\textbf{D.}] \textcolor{red}{\textbf{Registers}}
\end{itemize}

\vspace{0.5cm}

\vspace{0.5cm}

\textbf{Question 163:} The register that keeps the address of next instruction to be executed is \underline{\hspace{1.5cm}}.

\begin{itemize}
    \item[\textbf{A.}] AC
    \item[\textbf{B.}] \textcolor{red}{\textbf{PC}}
    \item[\textbf{C.}] IR
    \item[\textbf{D.}] MBR
    \item[\textbf{E.}] MQ
\end{itemize}

\vspace{0.5cm}

\vspace{0.5cm}

\textbf{Question 186:} Instruction register stores\underline{\hspace{1.5cm}}

\begin{itemize}
    \item[\textbf{A.}] Data of the current instruction
    \item[\textbf{B.}] Address of the next instruction
    \item[\textbf{C.}] \textcolor{red}{\textbf{Instruction which is currently executed}}
    \item[\textbf{D.}] Address of the current instruction
\end{itemize}

\vspace{0.5cm}

\vspace{0.5cm}

\textbf{Question 197:} The \underline{\hspace{1.5cm}} determines the opcode and the operand specifiers.

\begin{itemize}
    \item[\textbf{A.}] fetch operands
    \item[\textbf{B.}] calculate operands
    \item[\textbf{C.}] \textcolor{red}{\textbf{decode instruction}}
    \item[\textbf{D.}] execute instruction
\end{itemize}

\vspace{0.5cm}

\vspace{0.5cm}

\textbf{Question 199:} \underline{\hspace{1.5cm}} provide storage internal to the CPU.

\begin{itemize}
    \item[\textbf{A.}] Control units
    \item[\textbf{B.}] ALUs
    \item[\textbf{C.}] Main memory
    \item[\textbf{D.}] \textcolor{red}{\textbf{Registers}}
\end{itemize}

\vspace{0.5cm}

\vspace{0.5cm}

\textbf{Question 202:} Which one of the following CPU registers holds the address of the storage location being accessed?

\begin{itemize}
    \item[\textbf{A.}] MBR (Memory Buffer Register)
    \item[\textbf{B.}] IR (Instruction Register)
    \item[\textbf{C.}] \textcolor{red}{\textbf{MAR (Memory address register)}}
    \item[\textbf{D.}] AC (Accumulator)
\end{itemize}

\vspace{0.5cm}

\vspace{0.5cm}

\textbf{Question 209:} Instruction register stores\underline{\hspace{1.5cm}}

\begin{itemize}
    \item[\textbf{A.}] Address of the next instruction
    \item[\textbf{B.}] \textcolor{red}{\textbf{Instruction which is currently executed}}
    \item[\textbf{C.}] Address of the current instruction
    \item[\textbf{D.}] Data of the current instruction
\end{itemize}

\vspace{0.5cm}

\vspace{0.5cm}

\textbf{Question 243:} Condition codes facilitate multiway branches.

\begin{itemize}
    \item[\textbf{A.}] \textcolor{red}{\textbf{True}}
    \item[\textbf{B.}] False
\end{itemize}

\vspace{0.5cm}

\vspace{0.5cm}

\textbf{Question 267:} Which of the following components of CPU is responsible to direct the system to execute instructions?

\begin{itemize}
    \item[\textbf{A.}] Arithmetic and Logic Unit (ALU)
    \item[\textbf{B.}] \textcolor{red}{\textbf{Control Unit (CU)}}
    \item[\textbf{C.}] Registers
    \item[\textbf{D.}] Random Access Memory (RAM)
\end{itemize}

\vspace{0.5cm}

\vspace{0.5cm}

\textbf{Question 328:} \underline{\hspace{1.5cm}} : Determine the address of the next instruction to be executed. Usually, this involves adding a fixed number to the address of the previous instruction.

\begin{itemize}
    \item[\textbf{A.}] Instruction fetch
    \item[\textbf{B.}] Instruction operation decoding
    \item[\textbf{C.}] \textcolor{red}{\textbf{Instruction address calculation}}
    \item[\textbf{D.}] Operand fetch
    \item[\textbf{E.}] Operand address calculation
\end{itemize}

\vspace{0.5cm}

\vspace{0.5cm}

\textbf{Question 337:} Which registers can be assigned to a variety of functions by the programmer?

\begin{itemize}
    \item[\textbf{A.}] Data registers
    \item[\textbf{B.}] \textcolor{red}{\textbf{General purpose registers}}
    \item[\textbf{C.}] Address registers
    \item[\textbf{D.}] Condition codes (flags)
\end{itemize}

\vspace{0.5cm}

\vspace{0.5cm}

\textbf{Question 369:} \underline{\hspace{1.5cm}}: If the operation involves reference to an operand in memory or available via I/O, then determine the address of the operand.

\begin{itemize}
    \item[\textbf{A.}] Operand fetch
    \item[\textbf{B.}] Data operation
    \item[\textbf{C.}] Operand store
    \item[\textbf{D.}] \textcolor{red}{\textbf{Operand address calculation}}
\end{itemize}

\vspace{0.5cm}

\vspace{0.5cm}

\textbf{Question 392:} What is the role of the registers in a processor?

\begin{itemize}
    \item[\textbf{A.}] \textcolor{red}{\textbf{Registers in a processor provide fast, temporary storage for data and instructions, facilitating efficient access during instruction execution}}
    \item[\textbf{B.}] Registers are only used to store data temporarily during the execution of a program and do not contribute to the processing of instructions
    \item[\textbf{C.}] Registers are solely responsible for storing data from the main memory and have no involvement in holding instructions or facilitating data manipulation
    \item[\textbf{D.}] Registers are only necessary when the CPU is idle and have no impact on the speed or efficiency of instruction execution
\end{itemize}

\vspace{0.5cm}

\vspace{0.5cm}

\textbf{Question 405:} How many general-purpose registers are there in the Microprocessor Register Organizations of Intel 80386–Pentium 4 ?

\begin{itemize}
    \item[\textbf{A.}] \textcolor{red}{\textbf{8}}
    \item[\textbf{B.}] 16
    \item[\textbf{C.}] 4
    \item[\textbf{D.}] 12
\end{itemize}

\vspace{0.5cm}

\vspace{0.5cm}

\textbf{Question 536:} Choose the register that controls sequential execution.

\begin{itemize}
    \item[\textbf{A.}] \textcolor{red}{\textbf{Program counter (PC)}}
    \item[\textbf{B.}] Instruction register(IR)
    \item[\textbf{C.}] Memory address register (MAR)
    \item[\textbf{D.}] Memory buffer register (MBR)
\end{itemize}

\vspace{0.5cm}

\vspace{0.5cm}

\textbf{Question 498:} What is the primary purpose of index registers within the processor?

\begin{itemize}
    \item[\textbf{A.}] To manage control signals
    \item[\textbf{B.}] To hold condition codes
    \item[\textbf{C.}] \textcolor{red}{\textbf{To facilitate indexed addressing}}
    \item[\textbf{D.}] To execute arithmetic operations
\end{itemize}

\vspace{0.5cm}

\vspace{0.5cm}

\textbf{Question 513:} Steps for executing a machine instruction are concerned, given an assembly code: ADD EAX, V. What step will access the variable V?

\begin{itemize}
    \item[\textbf{A.}] \textcolor{red}{\textbf{The fetch data step.}}
    \item[\textbf{B.}] The fetch instruction step.
    \item[\textbf{C.}] The process data step.
    \item[\textbf{D.}] The interpret instruction step.
\end{itemize}

\vspace{0.5cm}

\vspace{0.5cm}

\textbf{Question 309:} What electronic component is used to govern operations such as fetching, decoding, and performing arithmetic operations executed by a processor?

\begin{itemize}
    \item[\textbf{A.}] \textcolor{red}{\textbf{Using a system clock}}
    \item[\textbf{B.}] Using a quartz crystal
    \item[\textbf{C.}] Using a analog to digital converter
    \item[\textbf{D.}] Using a counter
\end{itemize}

\vspace{0.5cm}

\vspace{0.5cm}

\section{Chapter 17: Reduced Instruction Set Computers (RISC)}
\textit{Total questions in chapter: 19}

\textbf{Question 294:} In the concept of Register Windows, how many register groups are there?

\begin{itemize}
    \item[\textbf{A.}] 4
    \item[\textbf{B.}] \textcolor{red}{\textbf{3}}
    \item[\textbf{C.}] 2
    \item[\textbf{D.}] No distinction
\end{itemize}

\vspace{0.5cm}

\vspace{0.5cm}

\textbf{Question 501:} What is a primary disadvantage of CISC (Complex Instruction Set Computer) architectures?

\begin{itemize}
    \item[\textbf{A.}] Increased implementation complexity
    \item[\textbf{B.}] Limited instruction set
    \item[\textbf{C.}] \textcolor{red}{\textbf{Slower execution due to complex instructions}}
    \item[\textbf{D.}] Higher energy consumption
\end{itemize}

\vspace{0.5cm}

\vspace{0.5cm}

\textbf{Question 42:} The instruction location immediately following the delayed branch is referred to as the \underline{\hspace{1.5cm}}

\begin{itemize}
    \item[\textbf{A.}] delay load
    \item[\textbf{B.}] delay file
    \item[\textbf{C.}] delay register
    \item[\textbf{D.}] \textcolor{red}{\textbf{delay slot}}
\end{itemize}

\vspace{0.5cm}

\vspace{0.5cm}

\textbf{Question 95:} A tactic similar to the delayed branch is the \underline{\hspace{1.5cm}}, which can be used on LOAD instructions

\begin{itemize}
    \item[\textbf{A.}] \textcolor{red}{\textbf{delayed load}}
    \item[\textbf{B.}] delayed program
    \item[\textbf{C.}] delayed slot
    \item[\textbf{D.}] delayed register
\end{itemize}

\vspace{0.5cm}

\vspace{0.5cm}

\textbf{Question 293:} Which statement is incorrect about RISC and CISC architecture?

\begin{itemize}
    \item[\textbf{A.}] CISC architecture is more convenient for programmers than RISC architecture.
    \item[\textbf{B.}] CISC architecture has more operands in a intruction compared to RISC architecture.
    \item[\textbf{C.}] CISC architecture has a more flexible instruction set than RISC architecture.
    \item[\textbf{D.}] \textcolor{red}{\textbf{CISC architecture requires more general-purpose registers than RISC architecture.}}
\end{itemize}

\vspace{0.5cm}

\vspace{0.5cm}

\textbf{Question 295:} What is the main benefit of using RISC over CISC?

\begin{itemize}
    \item[\textbf{A.}] RISC has more instructions and addressing modes than CISC
    \item[\textbf{B.}] \textcolor{red}{\textbf{RISC has faster instruction execution and simpler instruction decoding than CISC}}
    \item[\textbf{C.}] RISC has variable-length instruction formats and direct memory access than CISC
    \item[\textbf{D.}] RISC has more registers and pipelines than CISC
\end{itemize}

\vspace{0.5cm}

\vspace{0.5cm}

\textbf{Question 346:} What does CISC stand for?

\begin{itemize}
    \item[\textbf{A.}] \textcolor{red}{\textbf{Complex Instruction Set Computer}}
    \item[\textbf{B.}] Computer Instruction Set Complex
    \item[\textbf{C.}] Complex Instruction Summarize Computer
    \item[\textbf{D.}] Computer Instruction Summarize Complex
\end{itemize}

\vspace{0.5cm}

\vspace{0.5cm}

\textbf{Question 470:} What type of addressing is primarily used by almost all RISC instructions?

\begin{itemize}
    \item[\textbf{A.}] Immediate addressing
    \item[\textbf{B.}] \textcolor{red}{\textbf{Register addressing}}
    \item[\textbf{C.}] Indirect addressing
    \item[\textbf{D.}] Direct addressing
\end{itemize}

\vspace{0.5cm}

\vspace{0.5cm}

\textbf{Question 474:} Which of the following describes a RISC (Reduced Instruction Set Computer) architecture?

\begin{itemize}
    \item[\textbf{A.}] Many complex instructions
    \item[\textbf{B.}] \textcolor{red}{\textbf{Fixed instruction length}}
    \item[\textbf{C.}] Extensive use of memory addressing modes
    \item[\textbf{D.}] Emphasis on multi-tasking capabilities
\end{itemize}

\vspace{0.5cm}

\vspace{0.5cm}

\textbf{Question 348:} How does pipelining in a RISC architecture handle branch instruction?

\begin{itemize}
    \item[\textbf{A.}] \textcolor{red}{\textbf{By using NOOP instructions inserted by the compiler or assembler}}
    \item[\textbf{B.}] By eliminating branch instructions from the instruction stream
    \item[\textbf{C.}] By executing branch instructions in a separate pipeline
    \item[\textbf{D.}] By delaying all instructions until branch instructions are executed
\end{itemize}

\vspace{0.5cm}

\vspace{0.5cm}

\textbf{Question 298:} To enhance performance in a superscalar processor, which method(s) should we apply?

\begin{itemize}
    \item[\textbf{A.}] Duplication of resources.
    \item[\textbf{B.}] Out-of-order issue.
    \item[\textbf{C.}] Renaming registers
    \item[\textbf{D.}] \textcolor{red}{\textbf{All of the mentioned.}}
\end{itemize}

\vspace{0.5cm}

\vspace{0.5cm}

\textbf{Question 514:} In various statement types in a high-level programming language, what is not the most common statement?

\begin{itemize}
    \item[\textbf{A.}] \textcolor{red}{\textbf{GOTO statement.}}
    \item[\textbf{B.}] IF statement.
    \item[\textbf{C.}] CLOOP statement.
    \item[\textbf{D.}] CALL procedure statement.
\end{itemize}

\vspace{0.5cm}

\vspace{0.5cm}

\textbf{Question 33:} When using graph coloring, nodes that share the same color cannot be assigned to the same [register]

\begin{itemize}
    \item[\textbf{A.}] True
    \item[\textbf{B.}] \textcolor{red}{\textbf{False}}
\end{itemize}

\vspace{0.5cm}

\vspace{0.5cm}

\textbf{Question 83:} \underline{\hspace{1.5cm}} instructions are used to position quantities in registers temporarily for computational operations.

\begin{itemize}
    \item[\textbf{A.}] \textcolor{red}{\textbf{Load-and-store}}
    \item[\textbf{B.}] Window
    \item[\textbf{C.}] Complex
    \item[\textbf{D.}] Branch
\end{itemize}

\vspace{0.5cm}

\vspace{0.5cm}

\textbf{Question 219:} \underline{\hspace{1.5cm}} instructions are used to position quantities in registers temporarily for computational operations.

\begin{itemize}
    \item[\textbf{A.}] \textcolor{red}{\textbf{Load-and-store}}
    \item[\textbf{B.}] Window
    \item[\textbf{C.}] Complex
    \item[\textbf{D.}] Branch
\end{itemize}

\vspace{0.5cm}

\textbf{Question 234:} \underline{\hspace{1.5cm}} are bits set by the processor hardware as the result of operations.

\begin{itemize}
    \item[\textbf{A.}] MIPS
    \item[\textbf{B.}] \textcolor{red}{\textbf{Condition codes}}
    \item[\textbf{C.}] Stacks
    \item[\textbf{D.}] PSWs
\end{itemize}

\vspace{0.5cm}

\vspace{0.5cm}

\textbf{Question 292:} What is the main benefit of using ARM processors over other processors?

\begin{itemize}
    \item[\textbf{A.}] \textcolor{red}{\textbf{Low cost and low power consumption}}
    \item[\textbf{B.}] Higher degree of multi-tasking
    \item[\textbf{C.}] Lower error or glitches
    \item[\textbf{D.}] Efficient memory management
\end{itemize}

\vspace{0.5cm}

\vspace{0.5cm}

\textbf{Question 373:} What is the main difference between x86 and ARM instruction formats?

\begin{itemize}
    \item[\textbf{A.}] \textcolor{red}{\textbf{x86 instructions are variable in length, while ARM instructions are fixed}}
    \item[\textbf{B.}] x86 instructions are fixed in length, while ARM instructions are variable
    \item[\textbf{C.}] Both x86 and ARM instructions are fixed in length
    \item[\textbf{D.}] Both x86 and ARM instructions are variable in length
\end{itemize}

\vspace{0.5cm}

\vspace{0.5cm}

\textbf{Question 475:} Choose an INCORRECT statement about high-level programming languages (HLLs).

\begin{itemize}
    \item[\textbf{A.}] \textcolor{red}{\textbf{HLLs support directives for accessing the CPU's registers.}}
    \item[\textbf{B.}] HLLs allow the programmer to express algorithms more concisely.
    \item[\textbf{C.}] HLLs allow the compiler to take care of details that are not important in the programmer's expression of algorithms.
    \item[\textbf{D.}] HLLs often support the use of structured programming and/or object-oriented design.
\end{itemize}

\vspace{0.5cm}

\vspace{0.5cm}

\section{Chapter 18: Instruction-Level Parallelism and Superscalar Processors}
\textit{Total questions in chapter: 11}

\textbf{Question 407:} \underline{\hspace{1.5cm}} there are multiple functional units, each of which is implemented as a pipeline.

\begin{itemize}
    \item[\textbf{A.}] Scalar
    \item[\textbf{B.}] Superpipelined
    \item[\textbf{C.}] \textcolor{red}{\textbf{Superscalar}}
    \item[\textbf{D.}] Vector
\end{itemize}

\vspace{0.5cm}

\vspace{0.5cm}

\textbf{Question 476:} What role does branch prediction play in superscalar processors?

\begin{itemize}
    \item[\textbf{A.}] It reduces the need for registers
    \item[\textbf{B.}] It prevents the execution of any branches
    \item[\textbf{C.}] \textcolor{red}{\textbf{It minimizes pipeline stalling by guessing the paths of branches}}
    \item[\textbf{D.}] It reduces the number of functional units required
\end{itemize}

\vspace{0.5cm}

\vspace{0.5cm}

\textbf{Question 161:} A \underline{\hspace{1.5cm}} architecture is one that makes use of more, and more fine-grained pipeline stages.

\begin{itemize}
    \item[\textbf{A.}] parallel
    \item[\textbf{B.}] \textcolor{red}{\textbf{superpipelined}}
    \item[\textbf{C.}] superscalar
    \item[\textbf{D.}] hybrid
\end{itemize}

\vspace{0.5cm}

\vspace{0.5cm}

\textbf{Question 296:} What is the benefit of using a superscalar organization over a scalar organization?

\begin{itemize}
    \item[\textbf{A.}] \textcolor{red}{\textbf{It increases the instruction throughput and improves the performance}}
    \item[\textbf{B.}] It reduces the power consumption and the heat dissipation
    \item[\textbf{C.}] It simplifies the instruction set and the compiler design
    \item[\textbf{D.}] All of the mentioned
    \item[\textbf{E.}] None of the mentioned
\end{itemize}

\vspace{0.5cm}

\vspace{0.5cm}

\textbf{Question 297:} What does the term "instruction-level parallelism" refer to in computer architecture?

\begin{itemize}
    \item[\textbf{A.}] \textcolor{red}{\textbf{The degree to which instructions in a program can be executed in parallel}}
    \item[\textbf{B.}] The number of processor cores in a multi-core CPU with multiple resources
    \item[\textbf{C.}] The complexity of the instruction set architecture
    \item[\textbf{D.}] The length of an instruction cycle with high level programing language
\end{itemize}

\vspace{0.5cm}

\vspace{0.5cm}

\textbf{Question 393:} In SuperScalar computer architecture, the primary goal of utilizing multiple processors concurrently is to

\begin{itemize}
    \item[\textbf{A.}] Increase processing speed by increasing CPU frequency.
    \item[\textbf{B.}] \textcolor{red}{\textbf{Improve performance by executing more than one instruction per machine cycle.}}
    \item[\textbf{C.}] Reduce the size of the CPU to conserve energy.
    \item[\textbf{D.}] Enhance computational power by increasing the number of CPU cores.
\end{itemize}

\vspace{0.5cm}

\vspace{0.5cm}

\textbf{Question 180:} A \underline{\hspace{1.5cm}} is a small, very-high-speed memory maintained by the instruction fetch stage of the pipeline and containing the n most recently fetched instructions in sequence.

\begin{itemize}
    \item[\textbf{A.}] delayed branch
    \item[\textbf{B.}] \textcolor{red}{\textbf{loop buffer}}
    \item[\textbf{C.}] branch prediction
    \item[\textbf{D.}] multiple stream
\end{itemize}

\vspace{0.5cm}

\vspace{0.5cm}

\textbf{Question 349:} \underline{\hspace{1.5cm}} is (are) determined by the number of instructions that can be fetched and executed at the same time (the number of parallel pipelines) and by the speed and sophistication of the mechanisms that the processor uses to find independent instructions.

\begin{itemize}
    \item[\textbf{A.}] Instruction-level parallelism
    \item[\textbf{B.}] \textcolor{red}{\textbf{Machine parallelism}}
    \item[\textbf{C.}] Both instruction-level parallelism and machine parallelism
    \item[\textbf{D.}] None of the mentioned
\end{itemize}

\vspace{0.5cm}

\vspace{0.5cm}

\textbf{Question 240:} A \underline{\hspace{1.5cm}} hazard occurs when there is a conflict in the access of an operand location.

\begin{itemize}
    \item[\textbf{A.}] resource
    \item[\textbf{B.}] \textcolor{red}{\textbf{data}}
    \item[\textbf{C.}] structural
    \item[\textbf{D.}] control
\end{itemize}

\vspace{0.5cm}

\vspace{0.5cm}

\textbf{Question 502:} What is the primary method used to cope with storage conflicts in registers?

\begin{itemize}
    \item[\textbf{A.}] Increasing the size of the instruction set
    \item[\textbf{B.}] \textcolor{red}{\textbf{Register renaming}}
    \item[\textbf{C.}] Using shared memory
    \item[\textbf{D.}] Using technique of out-of-order issue with out-of-order completion
\end{itemize}

\vspace{0.5cm}

\vspace{0.5cm}

\textbf{Question 231:} A \underline{\hspace{1.5cm}} instruction can be used to account for data and branch delays.

\begin{itemize}
    \item[\textbf{A.}] SUB
    \item[\textbf{B.}] \textcolor{red}{\textbf{NOOP}}
    \item[\textbf{C.}] JUMP
    \item[\textbf{D.}] all of the above
\end{itemize}

\vspace{0.5cm}

\vspace{0.5cm}

\section{Chapter 19: Control Unit Operation and Microprogrammed Control}
\textit{Total questions in chapter: 2}

\textbf{Question 300:} How does multithreading improve the performance of a processor?

\begin{itemize}
    \item[\textbf{A.}] It increases the instruction-level parallelism by issuing multiple instructions from different threads in the same cycle
    \item[\textbf{B.}] It increases the thread-level parallelism by executing multiple threads on different cores or processors
    \item[\textbf{C.}] It increases the utilization of the processor resources by hiding the latency of long-latency events such as cache misses or branch mispredictions
    \item[\textbf{D.}] \textcolor{red}{\textbf{All of the mentioned}}
\end{itemize}

\vspace{0.5cm}

\vspace{0.5cm}

\textbf{Question 535:} With respect to multi-threading technique, choose an INCORRECT statement.

\begin{itemize}
    \item[\textbf{A.}] A thread is a dispatchable unit of work within a process.
    \item[\textbf{B.}] Each thread needs an individual program counter and stack pointer.
    \item[\textbf{C.}] Each thread has its own data area for a stack (to enable subroutine branching).
    \item[\textbf{D.}] \textcolor{red}{\textbf{A thread executes sequentially and interrupts cannot be allowed.}}
\end{itemize}

\vspace{0.5cm}

\vspace{0.5cm}

\section{Chapter 20: Parallel Processing}
\textit{Total questions in chapter: 5}

\textbf{Question 404:} Which architectural concept involves replicating register banks to facilitate the sharing of pipeline resources among multiple threads?

\begin{itemize}
    \item[\textbf{A.}] Pipelining
    \item[\textbf{B.}] Superscalar
    \item[\textbf{C.}] \textcolor{red}{\textbf{Simultaneous multithreading}}
    \item[\textbf{D.}] Superpipelining
\end{itemize}

\vspace{0.5cm}

\vspace{0.5cm}

\textbf{Question 351:} What is one advantage of Nonuniform Memory Access (NUMA) over Uniform Memory Access (UMA)?

\begin{itemize}
    \item[\textbf{A.}] \textcolor{red}{\textbf{NUMA provides each processor with its own local memory, reducing memory access times}}
    \item[\textbf{B.}] NUMA allows all processors to access the same memory location simultaneously
    \item[\textbf{C.}] NUMA is easier to implement than UMA
    \item[\textbf{D.}] NUMA provides limited memory capacity
\end{itemize}

\vspace{0.5cm}

\vspace{0.5cm}

\textbf{Question 377:} "Multiple processors share a single memory or pool of memory by means of a shared bus or other interconnection mechanism; a distinguishing feature is that the memory access time to any region of memory is approximately the same for each processor". Which concept does the statement belong to?

\begin{itemize}
    \item[\textbf{A.}] \textcolor{red}{\textbf{Symmetric multiprocessor (SMP)}}
    \item[\textbf{B.}] Nonuniform memory access (NUMA)
    \item[\textbf{C.}] Cluster
    \item[\textbf{D.}] Single instruction, multiple data (SIMD)
\end{itemize}

\vspace{0.5cm}

\vspace{0.5cm}

\textbf{Question 503:} What is a scalar instruction?

\begin{itemize}
    \item[\textbf{A.}] \textcolor{red}{\textbf{An instruction in which all operands must be single values.}}
    \item[\textbf{B.}] An instruction in which all operands must be groups such as arrays.
    \item[\textbf{C.}] An instruction in which all operands can be single values or groups.
    \item[\textbf{D.}] An instruction in which has no operand.
\end{itemize}

\vspace{0.5cm}

\vspace{0.5cm}

\textbf{Question 168:} Microprogramming eases the task of designing and implementing the control unit and provid family concept.

\begin{itemize}
    \item[\textbf{A.}] \textcolor{red}{\textbf{True}}
    \item[\textbf{B.}] False
\end{itemize}

\vspace{0.5cm}

\vspace{0.5cm}

\section{Chapter 21: Multicore Computers}
\textit{Total questions in chapter: 6}

\textbf{Question 302:} Which cache is not a shared cache?

\begin{itemize}
    \item[\textbf{A.}] L4 cache
    \item[\textbf{B.}] L3 cache
    \item[\textbf{C.}] L2 cache
    \item[\textbf{D.}] \textcolor{red}{\textbf{L1 cache}}
\end{itemize}

\vspace{0.5cm}

\vspace{0.5cm}

\textbf{Question 480:} What is a characteristic of Many Integrated Core (MIC) chips?

\begin{itemize}
    \item[\textbf{A.}] They include only specialized cores for specific tasks.
    \item[\textbf{B.}] \textcolor{red}{\textbf{They consist of homogeneous general-purpose processors on a single chip.}}
    \item[\textbf{C.}] They are designed to eliminate the need for GPUs in high-performance applications.
    \item[\textbf{D.}] They have limited scalability beyond four cores per chip.
\end{itemize}

\vspace{0.5cm}


\end{document}

\vspace{0.5cm}

\textbf{Question 350:} Which is the correct choice for the decription: "A single machine instruction controls the simultaneous execution of a number of processing elements such as vector and array processors"?

\begin{itemize}
    \item[\textbf{A.}] Single instruction, single data (SISD)
    \item[\textbf{B.}] \textcolor{red}{\textbf{Single instruction, multiple data (SIMD)}}
    \item[\textbf{C.}] Multiple instruction, single data (MISD)
    \item[\textbf{D.}] Multiple instruction, multiple data (MIMD)
\end{itemize}

\vspace{0.5cm}

\vspace{0.5cm}

\textbf{Question 299:} How many common classifications of parallel systems are there as proposed by Flynn?

\begin{itemize}
    \item[\textbf{A.}] 2
    \item[\textbf{B.}] 3
    \item[\textbf{C.}] \textcolor{red}{\textbf{4}}
    \item[\textbf{D.}] 5
\end{itemize}

\vspace{0.5cm}

\vspace{0.5cm}

\textbf{Question 477:} Which is the correct choice for the description: "A single processor executes a single instruction stream to operate on data stored in a single memory. "?

\begin{itemize}
    \item[\textbf{A.}] \textcolor{red}{\textbf{Single instruction, single data (SISD)}}
    \item[\textbf{B.}] Single instruction, multiple data (SIMD)
    \item[\textbf{C.}] Multiple instruction, single data (MISD)
    \item[\textbf{D.}] Multiple instruction, multiple data (MIMD)
\end{itemize}

\vspace{0.5cm}

\vspace{0.5cm}

\textbf{Question 504:} Which computing model is characterized by a large number of processors working on different parts of a problem simultaneously?

\begin{itemize}
    \item[\textbf{A.}] Single instruction, single data (SISD)
    \item[\textbf{B.}] Single instruction, multiple data (SIMD)
    \item[\textbf{C.}] \textcolor{red}{\textbf{Multiple instruction, multiple data (MIMD)}}
    \item[\textbf{D.}] Multiple instruction, single data (MISD)
\end{itemize}

\vspace{0.5cm}

\vspace{0.5cm}

\end{document}
